\pdfoutput=1
\documentclass[11pt]{article}
\usepackage[T1]{fontenc}
\usepackage[utf8]{inputenc}
\usepackage{lmodern,microtype}
\usepackage[a4paper,margin=27mm,headheight=14pt]{geometry}
\usepackage{amsmath,amssymb,amsthm,mathtools,bm}
\usepackage{graphicx,xcolor,booktabs,array,longtable,enumitem}
\usepackage[numbers,sort&compress]{natbib}
\renewcommand{\bibfont}{\small}
\usepackage{fancyhdr}
\usepackage{hyperref}
\usepackage[nameinlink,noabbrev]{cleveref}
\definecolor{paperblue}{HTML}{0077B5}
\definecolor{paperorange}{HTML}{DD7D20}
\hypersetup{colorlinks=true,linkcolor=paperblue,citecolor=paperblue,urlcolor=paperblue,
 pdftitle={Moment geometry and positive operator calculus},pdfauthor={}}
\graphicspath{{figures/}}
\setlength{\parskip}{3pt}
\setlength{\parindent}{1.2em}
\setlength{\emergencystretch}{2.5em}
\setcounter{tocdepth}{2}
\setcounter{secnumdepth}{3}
\AddToHook{cmd/section/before}{\par\begingroup
  \dimen0=\pagegoal\advance\dimen0 by-\pagetotal
  \ifdim\dimen0<7\baselineskip\newpage\fi\endgroup}
\setlist{itemsep=2pt,topsep=4pt}
\numberwithin{equation}{section}
\newtheorem{theorem}{Theorem}[section]
\newtheorem{lemma}[theorem]{Lemma}
\newtheorem{proposition}[theorem]{Proposition}
\newtheorem{corollary}[theorem]{Corollary}
\newtheorem{conjecture}[theorem]{Conjecture}
\theoremstyle{definition}
\newtheorem{definition}[theorem]{Definition}
\newtheorem{example}[theorem]{Example}
\theoremstyle{remark}
\newtheorem{remark}[theorem]{Remark}
\DeclareMathOperator{\tr}{tr}
\DeclareMathOperator{\ad}{ad}
\DeclareMathOperator{\Ent}{Ent}
\DeclareMathOperator{\Sym}{Sym}
\DeclareMathOperator{\supp}{supp}
\DeclareMathOperator{\rank}{rank}
\DeclareMathOperator{\diag}{diag}
\DeclareMathOperator{\Span}{span}
\DeclareMathOperator{\CP}{CP}
\DeclareMathOperator{\HS}{HS}
\newcommand{\C}{\mathbb C}
\newcommand{\R}{\mathbb R}
\newcommand{\N}{\mathbb N}
\newcommand{\Pp}{\mathbb P}
\newcommand{\Id}{\mathrm{Id}}
\newcommand{\dd}{\mathrm d}
\newcommand{\ddc}{dd^c}
\newcommand{\ket}[1]{|#1\rangle}
\newcommand{\bra}[1]{\langle#1|}
\newcommand{\norm}[1]{\left\|#1\right\|}
\newcommand{\abs}[1]{\left|#1\right|}
\newcommand{\one}{\mathbf 1}
\newcommand{\fig}[3]{\begin{figure}[tbp]\centering\includegraphics[width=\linewidth]{#1.pdf}\caption{#2}\label{#3}\end{figure}}
\pagestyle{fancy}
\fancyhf{}
\fancyhead[L]{\small Moment geometry and positive operator calculus}
\fancyhead[R]{\small\thepage}
\renewcommand{\headrulewidth}{.3pt}
\fancypagestyle{plain}{\fancyhf{}\fancyfoot[C]{\thepage}\renewcommand{\headrulewidth}{0pt}}

\title{Moment geometry and positive operator calculus\\[5pt]
\large Coherent states, algebraic support, and quantitative quantum certificates}
\author{}
\date{}
\begin{document}
\maketitle
\begin{abstract}
Positive moment constraints turn ensembles of pure quantum states into a
tractable meeting point of Hermitian algebra, complex geometry, and probability.
We develop this perspective together with a completely positive calculus for
modular response and tensor-network boundary maps. For balanced measures on
projective space, the optimal entropy--pluricomplex-energy coefficient is
expressed by harmonic numbers; on the projective plane it is $12/7$.
Positive moment repair yields smooth constrained equilibria and constrained
repulsive large deviations. A compactness theorem for conditional densities,
combined with a two-sided entropy estimate, also proves full attractive
determinant large deviations below the normalized plurisubharmonic
integrability threshold. On $\mathbb P^q$ this includes every fixed-order
shrinking design constraint for $0\le\gamma<1/(q+1)$.
A positive coherent-moment formula for the finite-energy Hessian also
yields a nonlinear Green representation of Calabi--Yau potentials by
random interpolation sections.
A classification of cubic supports of projective second designs follows by
turning a determinantal representation into a balanced three-qutrit tensor
and applying the Cartan theory associated with $E_6$. Its Hesse specialization
gives a robust ensemble certificate using three matched copies. A second,
ten-point design produces an explicit K3 surface with transcendental lattice
$A_2(7)$. For random coherent Gram matrices, directed spanning-tree counts
give all limiting spectral moments, imply unbounded support at every fixed
dimension, and explain the emergence of the
Marchenko--Pastur law as the complex dimension increases.
A separate Newton-polytope argument turns the phase-averaged Slater energy
into an explicitly certified relative-entropy assignment problem.
Anticanonical curvature gives finite-particle concentration and an exact
first-moment susceptibility.
On the operator side, signed heap logarithms lead to positive modular
resummation, with Gamma-clock approximations certified by local nested
commutators. Completely positive order estimates for primitive wires yield
dimension-independent core bounds and logarithmic sufficient decoration for
bounded-degree AKLT networks. An anisotropic contact theorem gives uniform
neighborhoods of qutrit measurements simulated at visibility $5/12$.
The paper explains how these constructions
connect geometric invariant theory, quantum information, random matrices,
and noncommutative analysis. Reproducible exact-arithmetic certificates,
illustrated computational tools, and theorem-level dependency records
accompany the proofs.
\end{abstract}
\vspace{5pt}
\noindent\textbf{Keywords.} Coherent states; projective designs; K\"ahler energy;
geometric invariant theory; quantum channels; tensor networks; modular response;
spectral gaps; random Gram matrices; large deviations; optimal transport.
\clearpage
\begingroup
\small
\setlength{\parskip}{0pt}
\tableofcontents
\endgroup
\clearpage
\section{Introduction: what positivity remembers}
\label{sec:introduction}

A probability measure on pure quantum states has several mathematical lives.
It is an ensemble for an experimentalist, a positive integration functional
for a geometer, and a sequence of symmetric tensor moments for an operator
theorist. Passing between these descriptions is useful only if the relevant
structure survives. The support of the measure, the Hermitian metric on the
ambient vector space, and the positivity of the integration weights all matter.
For example, two ensembles can have identical two-copy statistics and very
different algebraic supports. A third-copy observable can then detect a
geometric distinction that is invisible to every two-copy measurement.

The same discipline is useful for many-body operators. A small transfer-map
error in an ordinary norm need not remain small after contraction with an
arbitrary tensor environment. Completely positive order does survive that
contraction. Likewise, a scalar cluster expansion need not control an ordered
operator logarithm, whereas a positive average of modular-time conjugations
preserves complete positivity term by term. The central theme of this paper is
to identify representations in which the required positivity is preserved and
then extract quantitative consequences from them.

\subsection{The problem and established approaches}

The motivating physical question asks when a local frustration-free
Hamiltonian has a spectral gap bounded away from zero as the system grows.
The AKLT chain has a classical rigorous solution through its valence-bond
structure \citep{q:AKLT1988}. In higher dimensions, decorated AKLT models
admit rigorous tensor-network gap proofs
\citep{q:ARLLNY2020,LuciaYoung2023}; boundary-state factorization provides
another route, developed for the hexagonal model in the recent preprint
\citet{Kastoryano2026}. These arguments make transfer operators, overlaps
of local ground spaces, and the normalization of the physical energy
central objects. We seek estimates that remain valid for arbitrary
positive tensor environments, so that a finite channel calculation can
certify a whole family of Hamiltonians.

The geometric question has an analogous emphasis on preserved structure:
what can a positive moment constraint force about support, concentration,
and equilibrium? Classical projective designs fix low tensor moments,
geometric invariant theory controls balanced tensors, and pluripotential
theory converts concentration into a nonlinear energy. Each has its own
established proof methods. The aim here is to build explicit interfaces
between them, retaining the Hermitian metric and positivity throughout.
This produces exact support classifications and quantitative certificates,
as well as an analytic entropy threshold.

For the associated particle systems, positive-temperature determinant
large deviations are established by \citet{prob_berman2017}. On complex
curves, a Vandermonde factorization also relates attractive determinant
ensembles to logarithmic pair interactions. In higher dimension the
determinant has collective singularities, and its microscopic pressure is
a distinct issue from macroscopic coercivity
\citep{AndreassonBermanSvensson2026}. We state the sign, constraint, and
scope of each probability result explicitly. The same care distinguishes
known scalar subordination laws from the positive operator approximations
and observable-dependent error bounds constructed from them.

\subsection{Three interfaces and eight principal contributions}

The first interface is between \emph{moments and algebraic support}.
The $t$th moment of an ensemble is a positive matrix on a symmetric tensor
power. A polynomial vanishing on the support gives a kernel vector of that
matrix. A Hermitian polynomial nonnegative on the support gives a separating
functional. These are concrete connections to geometric quantization,
Hermitian sums of squares, and quantum de Finetti theory
\citep{Rao2024,GouveiaParriloThomas2010}.
The second interface is between \emph{determinants and concentration}:
vanishing orders of evaluation determinants measure energy gained by
concentrating a particle ensemble. The third is between \emph{completely positive
order and local quantum constraints}: finite matrix inequalities become
certificates for large tensor networks without introducing the condition
number of the tensor core.

The contributions can be read in the following order.
\begin{enumerate}[label=\arabic*.]
\item \textbf{A sharp balanced entropy--energy coefficient.}
For first-moment-balanced measures on $\Pp^d$, the strict-subcritical
coercivity threshold is
\[
\Gamma_d^{\mathrm{bal}}=\frac{d}{d+1-H_{d+1}},
\qquad H_m=\sum_{j=1}^m\frac1j.
\]
The proof uses minimal jet determinants, a moment-map argument, a conditioned
determinantal coupling, and a fixed-degree quantization estimate. Harmonic
numbers enter through the largest coordinate of a uniform simplex point.
Positive moment repair and exact energy comparison then give smooth
constrained equilibria on the projective plane and a full large-deviation
principle for repulsive determinant ensembles in shrinking design tubes.

\item \textbf{Coherent moments control determinant ensembles and geometric potentials.}
For a positive polarized compact complex manifold, we prove the full
attractive determinant Laplace principle below the exponential integrability
threshold of normalized plurisubharmonic potentials. On $\mathbb P^q$ with
anticanonical polarization the range is $0\le\gamma<1/(q+1)$, and every
fixed-order shrinking design tube is allowed. The proof combines compact
one-particle conditional densities, coarse-grained entropy, and a two-sided
small-coupling bound that imports mean-energy information from the known
repulsive sign. A separate curvature argument gives finite-particle
concentration. For $D=q+1$ and $M_1=N^{-1}\sum_iP_{z_i}$, a projective
Ward identity yields the exact susceptibility
\[
 \mathbb E\left\|M_1-I/D\right\|_{\rm HS}^2
 =\frac{D-1}{DN(1-\gamma)}.
\]
The identity holds whenever the unconditioned partition is finite and
$\gamma<1$; the LDP range above requires no such additional assumption.
The product mean energy has a positive Hessian expansion in squared
coherent moments of rank-two interpolation pencils. Its convexity
allows differentiation of the mean-energy limit and gives a logarithmic
interpolation formula for the inverse complex Monge--Amp\`ere operator.
The resulting potential estimator has an explicit sampling-variance
bound, uniform in the polynomial level.

\item \textbf{Cubic support is controlled by three-qutrit balance.}
We classify all plane cubics supporting a positive projective second design
in their given Hermitian embedding. A quartic rank witness completes a
two-factor-balanced determinantal representation to a fully balanced tensor.
The $E_6$ Cartan theory then forces a unitary Hesse form. Smooth members have
a unique supporting design; the singular case gives orthogonal triangles
with an explicit Fourier description.

\item \textbf{Algebraic rigidity becomes an operational certificate.}
A rank-two observable on three matched qutrit copies, together with the
one-copy moment defect, bounds transport distance to the Hesse ensemble.
The copy number is minimal in this access model. The geometry of a
complete intersection supplies a pointwise error bound, and the simplex
structure of the finite ensemble controls its weights.

\item \textbf{A finite design determines a K3 surface.}
A ten-point Singer design yields three disjoint $(-3)$-curves on a rational
surface. A cyclic triple cover and three contractions produce a K3 surface
with transcendental lattice $A_2(7)$. The calculation makes the relation
between moment constraints, special point configurations, and lattice
geometry explicit.

\item \textbf{Coherent Gram spectra interpolate between ordinary and
noncrossing partitions.}
At unit sampling density in complex dimension $d$, the limiting $p$th moment is
\[
m_p(d)=\sum_{\pi\in\mathcal P(p)}\tau(G_\pi)^{-d}.
\]
Here $G_\pi$ is the directed cycle quotient associated with a set partition
and $\tau$ counts rooted directed spanning trees. Crossing partitions
survive at every fixed dimension, but disappear in the limit $d\to\infty$.
For logarithmic energy, an independent Newton-polytope argument replaces
torus averaging by a finite assignment with Kullback--Leibler cost, with
an explicit error of order $\log N/k$. This connects coherent determinants
to tropical geometry and optimal transport through a computation that
admits exact rational dual certificates.

\item \textbf{Positive operator calculus gives quantitative compilers.}
We relate ordered logarithms to order polynomials, represent compatible
modular responses by positive random times, and give local-commutator error
bounds for finite clock approximations. A separate completely positive
wire comparison controls angles between tensor-network constraints and
yields logarithmic sufficient AKLT decoration in the maximum graph degree.

\item \textbf{Moment contact controls measurement simulation.}
The same simplex expected maximum fixes the Haar-parent visibility
$\eta_D=(H_D-1)/(D-1)$. A quantitative anisotropic contact analysis gives
an open simulation neighborhood of projective-simulable measurements in
every fixed dimension and outcome number, and of the larger convex hull
of reducible qutrit measurements at visibility $5/12$.
\end{enumerate}

These results share tools, but their proof dependencies are deliberately
explicit. For instance, the spectral-gap certificate uses completely positive
comparison and subspace geometry. It does not require a thermodynamic limit
for a determinant ensemble. The Hesse certificate uses finite-dimensional
algebra and transport, so it can be implemented without solving a nonlinear
equilibrium problem.

\subsection{Position within the literature}

The use of symmetry breaking by moment constraints in Fano geometry is
developed by Andreasson--Berman--Svensson \citep{AndreassonBermanSvensson2026}.
Our balanced energy calculation concerns their analytic, orbit-reduced
coercivity framework. The determinant ensemble has additional microscopic
structure. Theorems~\ref{att:full-laplace}--\ref{att:polarized-extension}
give an attractive large-deviation principle below the section integrability
threshold; extension to the canonical coupling is posed in the final
section. This separates a geometric stability coefficient from a statement
about every singular stratum of a many-particle determinant.

The mean-energy approach to singular Gibbs measures, including its
positive- and negative-temperature distinction, is developed by
\citet{Berman2018meanenergy}. Our argument supplies the requisite entropy
and energy convergence for the actual attractive tilted determinant laws
in a definite range. In particular, it uses the established repulsive
theorem without assuming product-law energy convergence as an input.
For these Slater Hamiltonians, this supplies the smooth-reference
mean-energy limit considered in \citet[Problem 4.10]{prob_berman2017}.
The nonlinear Green formula in Theorem~\ref{cy:green-formula} is the
representation formulated in \citet[Conjecture 6.1]{prob_berman2017}.
A later conditional proof uses a complex zero-free neighborhood for a
two-parameter partition function
\citep[Theorems 3.2--3.3]{Berman2023zeta}. Our proof obtains the formula
from real mean-energy convergence and the positive projective-kernel
Hessian, without asserting that zero-free property. This is a concrete
connection between coherent tensor moments, interpolation theory, and
the nonlinear equation solved by Yau's theorem \citep{Yau1978}.
Toric pluripotential duality and complex optimal transport are established
in \citet{BermanWittNystrom2014}; the finite assignment error bound makes
that interface computational. The curvature estimate uses classical
H\"ormander theory, while the projective change of variables evaluates a
specific covariance exactly.

Moment-improved Sobolev and Moser--Trudinger inequalities have a substantial
history. Recent sharp results on spheres under finite-order moment constraints
are given by Zhou--Zou \citep{ZhouZou2026}. The energy here is the nonlinear
dual Monge--Amp\`ere energy, and the proof uses projective jets and balanced
determinantal couplings. The role of positive cubature remains a useful
conceptual comparison, with the geometric energy and its normalization kept
distinct.

The invariant theory used in the cubic classification is established:
Kempf--Ness converts zero moment into closed orbits and unitary uniqueness,
and the three-qutrit Cartan representation is treated explicitly by
Di Trani--de Graaf--Marrani \citep{geo_ditrani2023}. The additional bridge is
the rank witness that forces a design-supported pencil into that representation's
balanced locus. The Hesse pencil and its classical incidence geometry are
described by Artebani--Dolgachev \citep{geo_hesse2009}.

For quantum spin systems, tensor-network methods already prove gaps for
decorated AKLT models. Lucia--Young obtain a sufficient decoration linear
in the maximum degree \citep{LuciaYoung2023}. The quantitative comparison
developed here uses the positive order of the half-wire Choi matrices;
the cancellation between two ports gives logarithmic sufficient decoration.
The elementary finite-dimensional certificate and the transfer back to the
specified microscopic Hamiltonian are stated separately, including the
finite-patch gap factor.

The scalar laws underlying modular clocks are also established. Hyperbolic
infinite divisibility and Brownian subordination are treated by Pitman--Yor
\citep{PitmanYor2003}. We use them as positive operator representations and
combine them with observable-specific nested-commutator estimates. This
distinguishes the scalar sampling identity from the quantitative operator
approximation it supports.

\subsection{Reading routes and conventions}

Readers coming from quantum information can start with the moment operators
below, then read the cubic classification and ensemble certificate.
Readers from complex geometry can begin with the balanced energy theorem and
the Singer surface construction, then the nonlinear Green formula in
Section~\ref{cy:section}. The spectral-law and finite assignment
sections are largely self-contained after the definition of coherent
vectors. For the probability results, Section~\ref{att:entropy-section}
gives the entropy mechanism and the full attractive LDP, while
Section~\ref{ac:section} gives finite-particle bounds and the susceptibility
identity. The positive
calculus and network sections can be read independently using ordinary
finite-dimensional operator theory.

All logarithms are natural. Pure states are rays represented by unit vectors,
and $P_z=zz^*$ is independent of their phase. Our pure-state distance is
\[
d_{\mathrm{tr}}(P_z,P_w)=\tfrac12\|P_z-P_w\|_1
=\sqrt{1-|\langle z,w\rangle|^2}.
\]
On $\Pp^d$, Fubini--Study measure has total mass one, and the geometric
energy uses the anticanonical polarization $(d+1)H$. Spectral gaps of blocked
parents refer to sums of norm-one projections. The conversion to a
nearest-neighbor Hamiltonian keeps its local energy scale explicitly.

The supplemental source provides full proof modules for long geometric
arguments, exact arithmetic certificates, figure-generation code, and an
index of theorem dependencies. The verification section describes the
mathematical meaning of those computations and their trusted arithmetic
assumptions.

\section{Coherent states, moments, and the geometry of positivity}
\label{sec:foundations}

\subsection{From an ensemble to a bosonic moment state}

Let $V=\C^D$ and $X=\Pp(V)$. For a probability measure $\mu$ on $X$ define
\[
M_t(\mu)=\int_X P_z^{\otimes t}\,d\mu(z)
\quad\text{on }\Sym^t V,
\qquad n_t=\binom{D+t-1}{t}.
\]
This is a trace-one positive operator. Schur's lemma and the trace give
$M_t(\sigma)=I_{n_t}/n_t$ for the unitary-invariant measure $\sigma$.
A weighted complex projective $t$-design is exactly a measure with this
moment. Positivity of the weights is part of the definition.

The vectors $z^{\otimes t}$ form the Veronese image. The set of possible
moment states is the convex hull of their projectors, a structured subset
of all density matrices. If $f$ is a homogeneous degree-$t$ polynomial
represented in the coherent/Bombieri inner product by a vector $v_f$, then
\[
\tr\bigl(|v_f\rangle\langle v_f|M_t(\mu)\bigr)
=\int|f(z)|^2\,d\mu(z).
\]
Thus a polynomial vanishes on the support exactly when its vector lies in
the kernel of the moment state. This simple identity is the mechanism by
which algebraic geometry becomes experimentally accessible through matched
copies of a hidden pure preparation.

More general Hermitian forms test inequalities on the support. One must
distinguish a matrix that is positive on all vectors of $\Sym^t V$ from a
matrix whose symbol is positive only on coherent vectors. The latter is a
positivity problem on an algebraic realization variety. Sum-of-squares
relaxations and ideal compression exploit that distinction
\citep{Rao2024,GouveiaParriloThomas2010}.

\fig{fig01_moment_fibres}{Moment equality preserves less information than
the ensemble itself. (a) The map $[z]\mapsto(|z_0|^2,|z_1|^2,|z_2|^2)$ sends
the nine Hesse rays to three edge midpoints, with three relative phases
above each. It sends the seven noncoordinate Singer rays to the barycenter;
the three other Singer rays are the vertices. This is a probability-simplex
projection, not a drawing of complex projective incidence. Hesse weights are
$1/9$; Singer vertex weights are $1/12$ and its other weights are $3/28$.
(b) Both ensembles and Haar measure have $M_2=I_6/6$, but the rank-two
three-copy projector $Q$ defined in the ensemble-certification theorem has
expectations $0$, $5/18$, and $1/5$, respectively. All displayed expectation
values are exact identities; decimals merely label the bars.}{fig:moments}

\subsection{Positive moment repair as a finite-dimensional inverse problem}

A useful feature of coherent moments is that small moment errors can be
repaired by adding a positive density. On $\Pp^2$ let $Q_z=P_z^{\otimes2}$,
$\tau=I_6/6$, and
\[
L(A)=(A\otimes I+I\otimes A)|_{\Sym^2\C^3},\qquad
B_z=90Q_z-30L(P_z)+5I_6.
\]
For traceless Hermitian $D$ on $\Sym^2\C^3$, the coherent fourth-moment
identity gives
\begin{equation}\label{eq:repairframe}
\int\tr(DB_z)\,d\sigma=0,\qquad
\int\tr(DB_z)Q_z\,d\sigma=D,\qquad
|\tr(DB_z)|\le51\|D\|_{\mathrm{HS}}.
\end{equation}
The number $51$ is a convenient rational upper bound: the eigenvalues of
$B_z$ are $35,-25,-25,5,5,5$, whose squared sum is $2550<51^2$.
Equation~\eqref{eq:repairframe} is an explicit right inverse of the moment
map, with a positivity radius.

\begin{proposition}[Positive repair]\label{prop:repair}
Let $\delta=\|M_2(\mu)-\tau\|_{\mathrm{HS}}>0$ and
$p=(1+102\delta)^{-1}$. There is a probability $\nu$ with
$1/2\le d\nu/d\sigma\le3/2$ such that
$\mu^*=p\mu+(1-p)\nu$ is a second design. Moreover,
\[
\Ent(\mu^*\mid\sigma)
\le p\Ent(\mu\mid\sigma)+(1-p)\log(3/2).
\]
\end{proposition}
\begin{proof}
For $D=M_2(\mu)-\tau$, take
$d\nu/d\sigma=1-\tr(DB_z)/(102\delta)$.
The preceding identities prove positivity, normalization, and exact moment
cancellation. Convexity of relative entropy proves the last inequality.
\end{proof}

This is an example of a general principle: invert the finite moment frame
first, then make the inverse perturbation small enough to remain a density.
It is stronger than an arbitrary signed correction, because positivity
allows entropy estimates and actual random sampling.

For later use the full frame operator is
\begin{equation}
 \mathsf S(B)=\int\tr(BQ_z)Q_z\,d\sigma
 =\frac{B+2L(\tr_2B)+\tr(B)I_6}{90}.
 \label{an:frame-spectrum}
\end{equation}
Here an operator on $\Sym^2\C^3$ is extended by zero on the antisymmetric
subspace before taking its partial trace. Haar's fourth tensor moment
gives the formula by permutation contraction. The scalar, adjoint, and
remaining irreducible sectors have real dimensions $1,8,27$ and frame
eigenvalues $1/6,1/15,1/90$, respectively. Thus the correction uses the
representation-theoretic decomposition of the moment space, while its
norm bound gives an explicit positive density.

\subsection{Multipoles explain the time-reversal grading}

For a spin-$j$ system, operator space decomposes under rotations as
\[
\mathcal B(V_j)\simeq\bigoplus_{\ell=0}^{2j}V_\ell.
\]
The spin-coherent symbol $a(n)=\langle n|A|n\rangle$ carries the same
multipole components, with nonzero normalization factors on each degree.
Time reversal sends $n$ to $-n$. For Hermitian $A$ invariant under time
reversal, the symbol is real and satisfies $a(n)=a(-n)$; therefore all odd
spherical-harmonic degrees vanish. On a register, simultaneous time reversal
removes odd \emph{total} multipole degree.

For qubits, the only one-site degrees are $0$ and $1$. Total degree then
coincides with body level, so odd body levels vanish. For spin one, each
nonidentity site can contribute degree $1$ or $2$, and a one-body
quadrupole is time-reversal even. The grading is therefore portable across
spins when expressed in multipoles. It should not be transported as an
odd-body rule beyond qubits. Standard irreducible-tensor and spin-coherent
conventions may be found in \citet{Radcliffe1971}.

\fig{fig11_multipoles}{Time reversal acts on multipole degree. The left two
panels show the exact restrictions of $n_z$ and $(3n_z^2-1)/2$ to a disk
parameterized by $n_z$, with a common signed color scale. Antipodal reversal
changes the dipole's sign and preserves the quadrupole. The right panel lists
possible spin-one pairs of body level and total degree for at most three
active sites, truncated at total degree four. Each active site contributes
degree one or two; blue marks even total
degree and orange marks odd total degree. In particular, the one-body,
degree-two entry survives time reversal. The identity is the entry $(0,0)$.}
{fig:multipoles}

\subsection{Two positivity notions that will recur}

For matrices $A,B$, $A\preceq B$ means $B-A$ is positive semidefinite.
For completely positive maps $\Phi,\Psi$, the notation
$\Phi\le_{\mathrm{CP}}\Psi$ means that $\Psi-\Phi$ is completely positive,
equivalently that their Choi matrices satisfy the corresponding Loewner
inequality. A completely positive comparison remains valid after tensoring
with an arbitrary identity map and applying the result to any positive
environment. This is why it is the appropriate input for network bounds.

On the geometric side, a positive polynomial witness is evaluated only on
the allowed coherent vectors. On the network side, a CP-order witness must
remain positive on all ancillary extensions. The two requirements are
different, but both replace a scalar coefficient test by positivity on the
correct realization space. This distinction is central to the later
certificates.

% Publication fragment: standard amsmath/amsthm environments.
% Bibliography: geometry_references.bib. Audit manifest: geometry_status.txt.
\section{Projective moment geometry}
\label{geo:section}

A probability distribution of pure quantum states can be studied through its
moments, through the algebraic equations satisfied by its support, or through
the singular potentials against which it is tested. These descriptions retain
different information. The first moment specifies an average density matrix;
the second moment specifies all quadratic experiments on matched preparations;
higher moments can distinguish ensembles whose lower moments coincide. On the
geometric side, a homogeneous equation of degree \(k\) becomes a positive test
on \(k\) matched copies. We use this relation between moments and equations
in both directions.

Throughout this section, \(P_z=zz^*/\|z\|^2\), and \(\sigma\) denotes the
unitary-invariant probability measure on the relevant complex projective
space. A projective \(k\)-design is a Borel probability measure satisfying
\begin{equation}
 M_k(\mu):=\int P_z^{\otimes k}\,d\mu(z)
 =\frac{I_{\operatorname{Sym}^k\mathbb C^D}}
 {\binom{D+k-1}{k}}.
 \label{geo:design}
\end{equation}
Partial trace gives all lower moment identities. The definition allows
weighted discrete designs and continuous measures. The connection between
projective second designs and informationally complete measurements is
standard \cite{geo_scott2006}; the results below exploit the additional
geometry of a design's support.

\subsection{A sharp first-moment entropy threshold in every dimension}
\label{geo:threshold}

The geometric energy is the Legendre transform of the Monge--Amp\`ere
energy. Its normalization matters. Put \(d=D-1\), \(X=\mathbb P^d\),
\(H=\mathcal O(1)\), \(L=-K_X=DH\), and \(\omega_0=D\omega_{\rm FS}\), with
\(\int_X\omega_{\rm FS}^d=1\). For a smooth strictly
\(\omega_0\)-plurisubharmonic potential \(u\), set
\begin{align}
 \mathcal E_L(u)
 &=\frac{1}{D^{d+1}}\sum_{j=0}^{d}
       \int_Xu\,\omega_u^j\wedge\omega_0^{d-j},
 &\omega_u&=\omega_0+dd^cu,\\
 E_L(\mu)
 &=\sup_u\left\{\mathcal E_L(u)-\int_Xu\,d\mu\right\}.
 \label{geo:energy}
\end{align}
Constants cancel in the second line. A measure is \emph{balanced} if
\(M_1(\mu)=I_D/D\). Balance is simultaneously a vanishing moment-map
condition and the assertion that the mean quantum state is maximally mixed.

\begin{theorem}[Sharp balanced entropy threshold]
\label{geo:sharp}
Let \(H_D=\sum_{j=1}^D j^{-1}\). The supremum of the coefficients
\(\gamma\) for which
\[
 \operatorname{Ent}(\mu\mid\sigma)\ge\gamma E_L(\mu)-C_\gamma
\]
holds uniformly over balanced finite-entropy measures is
\begin{equation}
 \Gamma_d^{\rm bal}=\frac{d}{D-H_D}.
 \label{geo:gamma}
\end{equation}
Every strict subcritical coefficient admits such a constant; every strict
supercritical coefficient fails. In particular, the first-moment constant
on \(\mathbb P^2\) is \(12/7\). The assertion concerns the supremal
coefficient and does not require boundedness at the endpoint.
\end{theorem}

The first values are
\[
 2,\qquad \frac{12}{7},\qquad \frac{36}{23},\qquad
 \frac{240}{163},\qquad \frac{100}{71}.
\]
The harmonic number has a concrete origin: for a uniform random point
\((x_0,\ldots,x_{D-1})\) of the probability simplex,
\(\mathbb E\max_i x_i=H_D/D\). The proof converts a balanced allocation of
monomials to projective-frame points into a statement about singularities,
then into an entropy inequality. This explains why the same constant
appears in a combinatorial allocation problem and in pluripotential theory.

\fig{fig04_jet_geometry}{Minimal jet allocation and the sharp balanced
coefficient. (a) Degree-24 monomials on $\mathbb P^2$ are plotted in their
normalized exponent simplex. The color is the minimal normalized order
$1-\max_i\alpha_i/24$ among the three coordinate centers. White boundaries
separate the maximizing-coordinate regions; ties can be allocated to meet
equal quotas. The continuum average is $7/18$. (b) The resulting
anticanonically normalized coefficient
$\Gamma_d^{\rm bal}=d/(d+1-H_{d+1})$ decreases toward one as complex
dimension grows. Its projective-plane value is $12/7$; the dashed line marks
one, not a separately proved microscopic integrability threshold.}{fig:jets}

\paragraph{Balanced jets and their invariant.}
Choose \(M\) with
\[
 N=\binom{M+d}{d}=Dn,\qquad
 S_M=\sum_{\substack{\alpha\in\mathbb N^D\\|\alpha|=M}}
                  (M-\max_i\alpha_i).
\]
There are arbitrarily large such levels: \(M=1+kD\) suffices. At coordinate
point \(e_i\), the order of \(z^\alpha\) is \(M-\alpha_i\). A joint jet
determinant using \(n\) rows at each of the \(D\) centers therefore has
total order at least \(S_M\). Equality is possible: assign each monomial
equally among its maximum-exponent centers, obtaining load \(N/D\) by
symmetry, and use integral max-flow to obtain the same integer quotas.

Collect the minimal-order minors into a nonzero vector \(P_0\) in
\[
 \left(\bigotimes_{i=0}^{D-1}
 \bigwedge^n\bigoplus_{\ell=0}^{M}
              \operatorname{Sym}^{\ell}(\mathbb C^d)^*\right)_{\deg S_M}.
\]
Use the invariant monomial norms \(\|x^\beta\|^2=\beta!/|\beta|!\).
After twisting by \(\prod_i\det_i^{-S_M/(Dd)}\), this vector has zero
moment map for \(\operatorname{GL}(d)^D\). A local root operator changes
one global monomial label, so its image has no overlap with the supported
assignments. Permutations of tangent coordinates kill the traceless
diagonal expectations, and permutations of centers distribute the total
degree equally. Polynomial pullback here is a right action; transposing
each matrix converts it to a left representation without changing the
determinant character.

Kempf--Ness, after clearing the rational character by a tensor power,
supplies a homogeneous semi-invariant \(F\) of some degree \(q>0\):
\begin{equation}
 F(P_0)\ne0,\qquad
 F(\rho(g)P)=\prod_i(\det g_i)^{qS_M/(Dd)}F(P).
 \label{geo:jet-semi-invariant}
\end{equation}
Let \(\mathcal U\subset X^D\) be the locus of \(D\) linearly independent
rays. Projective transformations act transitively on \(\mathcal U\).
All lower-total-order quota minors vanish identically there, and the
semi-invariant in \eqref{geo:jet-semi-invariant} is nonzero there.

\paragraph{From jets to log-canonical thresholds.}
For a basis \(s_1,\ldots,s_N\) of \(H^0(X,MH)\), form the balanced basis
ideal on \(X^D\),
\[
 \mathfrak b_s=
 \left\langle\prod_{i=0}^{D-1}\prod_{j\in B_i}s_j(z_i):
 \bigsqcup_i B_i=\{1,\ldots,N\},\quad |B_i|=n\right\rangle.
\]
The log-canonical threshold is the largest exponent for which the inverse
squared norm of ideal generators is locally integrable. On a resolution
it is computed by log discrepancy divided by ideal order; discrepancy
records the Jacobian cost of changing coordinates.

Let \(v_i\) be a monomial valuation on a resolving chart
\(F_i:y_i\mapsto x_i\), with weights \(a_{i1},\ldots,a_{id}\ge0\), and
\(J_i=\det dF_i\). Then
\[
 A_X(v_i)=\sum_j a_{ij}+v_i(J_i).
\]
The chain rule for the minimal jet vector has no surviving nonlinear
correction: every such term contains a lower-total-order minor, which
vanishes on \(\mathcal U\). Therefore
\[
 \widetilde P=\rho(dF_i)P(F_i),\qquad
 F(\widetilde P)=
 \prod_iJ_i^{qS_M/(Dd)}F(P(F_i)).
\]
The last factor is a valuation unit. Scaling jet rows by resolving
coordinates charges their differentiation costs to the same determinant
character. Some minimal minor consequently satisfies
\[
 \text{valuation of minor}+\text{differentiation cost}
 \le\frac{S_M}{Dd}\sum_i A_X(v_i).
\]
Apply its derivatives to the evaluation determinant in independent
particle variables, and only then identify the \(n\) variables within
each cluster. Identification can increase valuation. Thus the displayed
bound gives an upper bound for the relevant product valuation of the
independent-particle determinant.

Resolve the divisors of all \(s_j\) separately in each factor. The
pullback of \(\mathfrak b_s\) is a monomial ideal on the product resolution,
so its threshold is computed by valuations of this product form. Its
order is at most the determinant order. Hence
\begin{equation}
 \operatorname{lct}_{\mathcal U}(\mathfrak b_s)
 \ge\frac{Dd}{S_M}.
 \label{geo:basis-lct}
\end{equation}
When \(S_M=0\), the ideal is the unit ideal on \(\mathcal U\).
This statement concerns balanced basis ideals; arbitrary valuations on a
full particle configuration space need not be product valuations.

\paragraph{The analytic transfer.}
Fix a compact subset of \(\mathcal U\) and one degree \(M\). Strictly
below \eqref{geo:basis-lct}, inverse ideal norms are integrable uniformly
over the compact group of orthonormal section bases, by semicontinuity
of complex singularity exponents \cite{geo_demailly_kollar2001}.
For \(Q(z)=\sum_j e^{\lambda_j}|s_j(z)|^2\), expansion gives
\[
 \prod_iQ(z_i)^n\ge e^{\sum_j\lambda_j}
 \sum_B\left|\prod_i\prod_{j\in B_i}s_j(z_i)\right|^2.
\]
This yields the quantized local inequality for every
\(\beta<dMN/S_M\). No estimate uniform in growing \(M\) is used.
The normalized monomial lattice converges to the uniform simplex law:
\begin{equation}
 \frac{S_M}{MN}\longrightarrow\frac{D-H_D}{D}.
 \label{geo:monomial-limit}
\end{equation}
Indeed, integrating
\[
 \mathbb P(\max_i x_i>a)
 =\sum_{j=1}^D(-1)^{j+1}\binom Dj(1-ja)_+^{D-1}
\]
gives \(\mathbb E\max_i x_i=H_D/D\).

For the Hilbert map
\[
 H_\psi(s,t)=\int_Xh_{\rm FS}^M(s,t)e^{-M\psi}\,d\sigma,\qquad
 v=\operatorname{FS}(H_\psi),
\]
and \(E_M(v)=(MN)^{-1}\log(\det H_0/\det H_\psi)\), one has the exact
comparison
\begin{equation}
 \mathcal E_H(\psi)\le E_M(v),\qquad
 \int_Xe^{M(v-\psi)}\,d\sigma=N.
 \label{geo:exact-quantization}
\end{equation}
Here \(\mathcal E_H\) uses \(\omega_{\rm FS}\). To justify the exact
coefficient, apply the direct-image maximum principle of
\cite[Propositions 4.1--4.2]{geo_zhang2024} along the weak geodesic
from \(0\) to \(\psi\). The auxiliary bundle \(MH-K_X\) has curvature
\[
 M(\pi^*\omega_{\rm FS}+dd^c\psi)+D\pi^*\omega_{\rm FS}\ge0,
\]
and the reference Bergman density is exactly \(N\). Differentiating the
Hilbert/Bergman geodesic comparison at the reference endpoint gives the
first inequality; the second is the trace of Hilbert orthonormality.
This exact endpoint-derivative specialization is spelled out again in
Theorem~\ref{an:lower-comparison}; it is not obtained by setting the error
parameter to zero in an asymptotic quantization bound.

Given \(\lambda<dD/(D-H_D)\), choose a larger strict-subcritical \(\beta\)
and a single sufficiently large admissible \(M\) with
\(\lambda/M+\lambda/\beta\le1\). H\"older's inequality with these weights
transfers the quantized estimate using \eqref{geo:exact-quantization}.
Since \(\mathcal E_L(D\psi)=D\mathcal E_H(\psi)\), for
\(\gamma<d/(D-H_D)\) this gives
\begin{equation}
 \gamma\mathcal E_L(u)+\frac1D\sum_{i=0}^{D-1}
          \log\int_{U_i}e^{-\gamma u}\,d\sigma
 \le C_{\gamma,U},
 \label{geo:local-entropy}
\end{equation}
whenever the \(U_i\) have disjoint closures with closed product in
\(\mathcal U\).

\paragraph{Globalization by a determinantal coupling.}
For unit representatives set
\[
 \mathfrak D(z_1,\ldots,z_D)=|\det(z_1,\ldots,z_D)|,\qquad
 c_D=\frac{D^{D-1}}{(D-1)!}.
\]
If \(\nu\) is balanced, \(c_D\mathfrak D^2\nu^{\otimes D}\) has all
marginals \(\nu\). Assume first that \(\nu\ge\epsilon\sigma\).
The conditional operator and the derivative of common marginal scaling are
\[
 Tf(z)=\frac D{D-1}\left(\int f\,d\nu-
     \operatorname{tr}\left(P_z\int fP\,d\nu\right)\right),\qquad
 \mathcal L=I+(D-1)T.
\]
On traceless Hermitian matrices put
\(\mathcal B(M)=D\int\operatorname{tr}(MP)P\,d\nu\). Haar's second
moment and the variance inequality imply
\[
 0\le\mathcal B\le\left(1-\frac{D\epsilon}{D+1}\right)I.
\]
For \(m=\int g\,d\nu\), \(G=\int g(P-I/D)\,d\nu\), the inverse is
\begin{equation}
 \mathcal L^{-1}g=g-\frac{D-1}{D}m+
 D\operatorname{tr}\bigl((I-\mathcal B)^{-1}G\,P\bigr),\qquad
 \|\mathcal L^{-1}\|_{\infty\to\infty}\le\kappa:=\frac{D+3}{\epsilon}.
 \label{geo:marginal-inverse}
\end{equation}
Truncate the kernel to \(\mathfrak D\ge\delta\), where
\(\delta=(8\kappa D\sqrt{c_D})^{-1}\). Each marginal loses at most
\(\eta=c_D\delta^2\). Write the common scaling as \(a=1+h\). Its
equation is \(-r+\mathcal Lh+R(h)=0\), with \(\|r\|_\infty\le\eta\).
On \(\|h\|_\infty\le b=2\kappa\eta\),
\[
 \operatorname{Lip}(R)\le D\eta+2D(D-1)b,\qquad
 \kappa\operatorname{Lip}(R)<\frac18.
\]
Contraction gives the exact coupling
\[
 d\pi=\prod_i a(z_i)c_D\mathfrak D^2
          \mathbf1_{\{\mathfrak D\ge\delta\}}\,d\nu^{\otimes D}.
\]
A fixed finite partition of diameter at most \(\delta/(2D)\) in
projective chordal distance localizes it into finitely many frame
products, since the determinant is separately \(1\)-Lipschitz. Each good
closed product has determinant at least \(\delta/2\), hence disjoint
coordinate-cell closures. Slight enlargements give the \(U_i\) in
\eqref{geo:local-entropy}. Conditioning on labels expresses
\(\nu=\sum_jp_j\nu_j\), where \(\nu_j\) is the equal average of \(D\)
measures in one frame product, and the number \(J\) depends only on
\(D,\epsilon\). The finite-mixture entropy identity gives
\[
 \sum_jp_j\operatorname{Ent}(\nu_j\mid\sigma)
 \le\operatorname{Ent}(\nu\mid\sigma)+\log J.
\]
Entropy duality in \eqref{geo:local-entropy} and energy convexity now
prove the required inequality for \(\nu\ge\epsilon\sigma\).

To remove the Haar floor, normalize \(\int u\,d\sigma=0\) and set
\(J(u)=-\mathcal E_L(u)\). Integration by parts gives
\(I(u)\ge(1+1/d)J(u)\) and \(c\frac{d}{dc}J(cu)=I(cu)\), so
\(J(cu)\le c^{1+1/d}J(u)\). Taking \(c=\theta^d\) in the Legendre
transform proves the extended-valued inequality
\begin{equation}
 E_L(\theta\mu+(1-\theta)\sigma)\ge\theta^{d+1}E_L(\mu).
 \label{geo:dilution}
\end{equation}
Entropy convexity transfers the estimate to \(\mu\) with coefficient
\(\gamma\theta^d\). Letting \(\theta\uparrow1\) proves every
strict-subcritical case of Theorem~\ref{geo:sharp}.

\paragraph{Explicit sharpness.}
For \(t>\log2\), set
\[
 C_i(t)=\{[z]:1-|z_i|^2\le e^{-t}\},\qquad
 \mu_t=\frac{\mathbf1_{\cup_i C_i(t)}}{De^{-dt}}\,\sigma.
\]
The disjoint caps each have mass \(e^{-dt}\). Their torus and permutation
symmetries make \(\mu_t\) balanced, and its entropy is \(dt-\log D\).
For an orthonormal monomial basis and \(r_\alpha=M-\max_i\alpha_i\),
consider
\[
 \psi_t=\frac1M\log\sum_{|\alpha|=M}e^{tr_\alpha}|s_\alpha|^2.
\]
Subgeodesic convexity gives
\(\mathcal E_H(\psi_t)\ge(\log N)/M+tS_M/(MN)\). On cap \(i\),
\(r_\alpha\le M-\alpha_i\), so the multinomial theorem gives
\(\psi_t\le(\log N)/M+\log2\). Passing \(M\to\infty\) in energy duality,
\[
 E_L(\mu_t)\ge(D-H_D)t-D\log2.
\]
Every strict-supercritical coefficient fails along this family. Unitary
convolution and a small Haar admixture produce smooth positive balanced
examples with the same limiting obstruction.

\paragraph{Relation to stability theory.}
Moment-constrained entropy coercivity is the analytic threshold in the
probabilistic approach to K\"ahler--Einstein metrics with continuous
automorphisms. In the stated anticanonical normalization,
\cite[Propositions 3.6 and 3.10]{AndreassonBermanSvensson2026} identify this balanced
threshold on projective space with the analytic polystability threshold.
The theorem calculates that analytic quantity; microscopic integrability
of an interacting particle determinant requires additional information.

\subsection{Which plane cubics can carry a second design?}
\label{geo:cubic-classification}

A second design is highly isotropic, but its support need not fill
projective space. The following classification identifies exactly how
a cubic equation can coexist with that isotropy. Three structures enter
the proof essentially: a rank inequality for matrix pencils, the
moment-map description of multipartite entanglement, and the order-three
grading of the exceptional Lie algebra \(E_6\).

Let \(\omega=e^{2\pi i/3}\), and define the nine Hesse rays
\begin{equation}
 \mathcal H=\left\{
 \frac{(0,1,-\omega^j)}{\sqrt2},
 \frac{(1,0,-\omega^j)}{\sqrt2},
 \frac{(1,-\omega^j,0)}{\sqrt2}:j=0,1,2\right\},
 \qquad
 \mu_{\mathcal H}=\frac19\sum_{h\in\mathcal H}\delta_h.
 \label{geo:hesse-points}
\end{equation}
They form a symmetric informationally complete set:
\(\operatorname{tr}(P_hP_{h'})=1/4\) for \(h\ne h'\).
Their incidence geometry and linear dependencies are discussed in
\cite{geo_hesse2009,geo_dang2013}.

\begin{theorem}[Complete cubic-support classification]
\label{geo:cubic}
For a nonzero homogeneous ternary cubic \(F\), a projective second design
supported on \(\{F=0\}\) exists if and only if, after a unitary coordinate
change and multiplication by a nonzero scalar, \(F\) is either
\[
 z_0^3+z_1^3+z_2^3-\lambda z_0z_1z_2,\qquad \lambda^3\ne27,
\]
or \(z_0z_1z_2\).
In the smooth case the measure is uniquely \(\mu_{\mathcal H}\), in the
corresponding unitary coordinates. On the coordinate triangle every such
measure is
\[
 \mu=\frac13(\nu_0+\nu_1+\nu_2),
\]
where \(\nu_i\) is supported on \(z_i=0,\ |z_j|=|z_k|\), and its
first two phase Fourier coefficients vanish.
\end{theorem}

The conclusion is unitary, rather than merely projective: a general
projective transformation changes the prescribed Hermitian moments.
Nodal or cuspidal cubics, a line plus an irreducible conic, nonorthogonal
triangles, concurrent triples, and repeated components are excluded.

\begin{proof}
\emph{Step 1: a determinant with two balanced tensor factors.}
Every ternary cubic has a \(3\times3\) linear determinantal representation
\[
 F(z)=\det M(z),\qquad M(z)=\sum_{i=0}^2z_iM_i.
 \tag{a}
\]
For a smooth cubic, use the Hesse normal form
\cite[Lemma 1]{geo_hesse2009} and the Moore matrix
\[
 \mathsf M(a,z)=
 \begin{pmatrix}
 a_0z_0&a_1z_2&a_2z_1\\
 a_2z_2&a_0z_1&a_1z_0\\
 a_1z_1&a_2z_0&a_0z_2
 \end{pmatrix},
 \quad
 \det\mathsf M=S(a)T(z)-T(a)S(z),
\]
where \(S(z)=\sum_i z_i^3\), \(T(z)=z_0z_1z_2\). Choose \(a\) on the
cubic with \(T(a)\ne0\).
For a singular irreducible cubic, put a singular point at
\([0:0:1]\), so \(F=zQ_2(x,y)+Q_3(x,y)\). If \(Q_2\) has distinct
linear factors, normalize it to \(xy\) and absorb the terms divisible
by \(xy\) into \(z\), obtaining \(xyz+ax^3+dy^3\).
If \(Q_2\) is a square, normalize it to \(y^2\), absorb its multiples
into \(z\), then translate \(x\) to obtain \(x^3-y^2z\).
The two resulting determinants are
\[
 \begin{pmatrix}x&0&y\\0&y&x\\-dy&-ax&z\end{pmatrix},
 \qquad
 \begin{pmatrix}x&y&0\\0&x&z\\-y&0&x\end{pmatrix}.
\]
A reducible cubic is a line times a conic; represent the conic by a
\(2\times2\) linear determinant and add the line as a diagonal block.
This includes repeated factors.

Regard the coefficients of \(M\) as
\(\Psi\in(\mathbb C^3)^{\otimes3}\), and put
\[
 s=\|\Psi\|^2,\quad
 R=\sum_iM_iM_i^*,\quad
 C=\sum_iM_i^*M_i,\quad
 T_{ij}=\operatorname{tr}(M_iM_j^*).
\]
Up to a transpose for the second factor, these are its one-factor
reduced density matrices, each of trace \(s\).
The coefficients of \(\det M(z)\) are invariant under
\(M\mapsto AMB^{\mathsf T}\), \(A,B\in SL_3\).
Their nonvanishing prevents zero from belonging to the orbit closure.
The closure has a nonzero point of minimum Euclidean norm. Differentiating
its norm in all traceless Hermitian directions on the first two factors
gives
\[
 R=C=\frac{s}{3}I.
 \tag{b}
\]
The determinant remains exactly \(F\) at the limit. This argument does
not assume a finite invertible scaling of the original pencil.

\emph{Step 2: the design completes the balance.}
The Hermitian quartic
\[
 q_M(z)=2\operatorname{tr}\bigl((M(z)M(z)^*)^2\bigr)
       -\bigl(\operatorname{tr}(M(z)M(z)^*)\bigr)^2
\]
is nonnegative on \(\det M=0\). Indeed, a positive matrix of rank at most
two, with possibly zero eigenvalues \(b_1,b_2\), gives
\(q_M=(b_1-b_2)^2\). Haar's second-moment contractions give
\[
 12\int q_M\,d\sigma
 =2\operatorname{tr}R^2+2\operatorname{tr}C^2
      -\operatorname{tr}T^2-s^2.
 \tag{c}
\]
If a second design is supported on the cubic, then (b), (c), and its
moment identity imply
\[
 0\le12\int q_M\,d\mu
 =-\left\|T-\frac{s}{3}I\right\|_{\rm HS}^2\le0.
\]
Thus the third marginal is scalar. The measure constraint has forced a
three-qutrit tensor to be fully balanced.

\fig{fig03_tensor_balance}{The tensor used in the cubic classification.
The three matrix slices display the nonzero coefficients of a vector in
the Cartan subspace, with illustrative values $a=1$, $b=0.65$, and $c=0.35$.
After normalization, all three one-qutrit marginals are $I_3/3$, regardless
of these coefficients. The slice determinant is
$(a^3+b^3+c^3)z_0z_1z_2-abc(z_0^3+z_1^3+z_2^3)$, a member of the Hesse
pencil. The proof first obtains full marginal balance from the design
identity and only then invokes the $E_6$ Cartan representation.}{fig:tensor}

\emph{Step 3: the exceptional-Lie-algebra input.}
For the \(SL_3^3\) action, scalar one-factor marginals characterize
critical vectors. Their orbits are closed, and two critical vectors in
one orbit are locally unitarily equivalent, by Kempf--Ness
\cite[Theorem 14]{geo_sauerwein2018}. Vinberg's Cartan theorem for the
order-three \(E_6\) representation states that every closed orbit meets
\[
 \mathfrak c=\operatorname{span}\left\{
 |000\rangle+|111\rangle+|222\rangle,\,
 |012\rangle+|120\rangle+|201\rangle,\,
 |021\rangle+|210\rangle+|102\rangle\right\}.
\]
The statement includes nongeneric semisimple orbits
\cite[Section 7.1]{geo_ditrani2023}. Every vector in this space already
has scalar marginals: the three supports are disjoint Latin squares.
The tensor from (b) is consequently related to \(a u_1+b u_2+c u_3\)
by local unitaries. Its determinant is
\[
 (a^3+b^3+c^3)z_0z_1z_2
             -abc(z_0^3+z_1^3+z_2^3).
 \tag{d}
\]
The first two local unitaries change the determinant by a scalar; the
third gives a unitary change of its variables. Its nonzero singular
members are precisely orthogonal triangles. For example,
\[
 x^3+y^3+z^3-3xyz
 =(x+y+z)(x+\omega y+\omega^2z)(x+\omega^2y+\omega z),
\]
and the three coefficient vectors are orthogonal.

\emph{Step 4: rigidity on a smooth Hesse cubic.}
Define
\[
 r(z)=\sum_i|z_i|^4-\frac12\|z\|^4,\qquad
 D(z)=\sum_i z_i^2\overline{z_{i+1}z_{i+2}}.
\]
The Moore matrix satisfies
\[
 q_{\mathsf M(a,\cdot)}(z)
 =4r(a)r(z)+8\operatorname{Re}(D(a)\overline{D(z)}).
\]
On a common Hesse cubic, therefore, the vectors
\(v(z)=(r(z),\sqrt2\operatorname{Re}D(z),\sqrt2\operatorname{Im}D(z))\)
have pairwise nonnegative inner products. A second design has mean
\(v=0\). The nonnegative continuous kernel \(v(a)\cdot v(z)\) then has
zero double integral, forcing \(v=0\) on the entire support.

Suppose all coordinates of such a point are nonzero. For a unit
representative, put \(p_i=|z_i|^2\), \(u_i=z_i^3/p_i\).
The equations \(r=D=0\) imply
\[
 \sum_i u_i=0,\qquad |u_i|^2=p_i,\qquad \sum_i p_i^2=\frac12.
\]
Heron's identity gives zero area for the triangle of vectors \(u_i\).
Hence \(u_i=e^{i\theta}s_i\), where \(s_i\) are real,
\(\sum_i s_i=0\), and \(s_i^2=p_i\). Thus
\(z_i=e^{i\theta/3}\zeta_i s_i\), \(\zeta_i^3=1\), and
\[
 \frac{S(z)}{T(z)}=\frac3{\zeta_0\zeta_1\zeta_2}.
\]
This is a singular Hesse parameter. On a smooth member one coordinate
must therefore vanish; \(r=0\) equalizes the other two moduli, and the
cubic equation makes their cubes opposite. The support is contained
in \(\mathcal H\). Pairing \(M_2=I_6/6\) with \(P_i^{\otimes2}\)
gives
\[
 \frac16=w_i+\frac1{16}(1-w_i),
\]
so \(w_i=1/9\). Conversely those weights form a second design, by direct
summation or equality in \(\operatorname{tr}M_2^2\ge1/6\).

\emph{Step 5: the triangle.}
On \(xyz=0\), \(\sum_i|z_i|^4\ge1/2\), whose second-design expectation
is \(1/2\). Thus the support lies on the three midpoint circles.
The mixed moment \(\int|z_j|^2|z_k|^2\,d\mu=1/12\) gives mass \(1/3\)
to each circle. The off-diagonal identities are precisely vanishing
of its first and second phase Fourier coefficients. Conversely those
conditions give every bidegree-\((2,2)\) moment.
\end{proof}

\paragraph{Why this tensor format is special.}
For an \(n\times n\) pencil in \(d\) variables, the rank witness
\((n-1)\operatorname{tr}A^2-(\operatorname{tr}A)^2\) is nonnegative on
the determinant hypersurface. After balancing its first two factors,
a supported second design would imply
\[
 0\le\frac{n-2}{n}s^2-\operatorname{tr}T^2
 \le\left(\frac{n-2}{n}-\frac1d\right)s^2.
\]
The strictly negative range excludes designs. Equality forces the
third factor to balance, and for integers \(n,d\ge2\) occurs exactly at
\((n,d)=(3,3)\) and \((4,2)\). The binary-quartic case recovers a regular
tetrahedron on the Bloch sphere with uniform weights. Indeed, for a
second design on \(\mathbb P^1\), the nonnegative polynomial
\((3|\langle p,z\rangle|^2-1)^2\) has mean one and value four at \(p\),
so every atom has mass at most \(1/4\). A quartic has at most four
distinct zeros; equality therefore forces four equal weights, with
the polynomial zero at every other support point. Their squared
overlaps are \(1/3\), precisely the regular tetrahedron condition.
The three-qutrit case is a critical format of the rank inequality,
which explains the effectiveness of the exceptional-group description.

\subsection{An algebraic certificate that becomes a quantum experiment}
\label{geo:hesse-cert}

The cubic classification identifies an ensemble from exact equations.
The next result is quantitative and operational. A request of size \(k\)
means that a hidden ray \(z\sim\mu\) is drawn once and \(k\) copies of that
same state are prepared; independent requests use independent draws.
This is a matched-replica model with trusted dimension, basis, and
measurements. Its average request state is \(M_k(\mu)\). Redrawing the
hidden ray separately for every copy would instead reveal only \(M_1\).

Use pure-state trace distance
\[
 d(z,w)=\frac12\|P_z-P_w\|_1
       =\sqrt{1-|\langle z,w\rangle|^2}
\]
and its \(1\)-Wasserstein distance \(W_1\) on ensembles.
In \(\operatorname{Sym}^3\mathbb C^3\), define orthonormal vectors
\[
 \phi_F=\frac{|000\rangle+|111\rangle+|222\rangle}{\sqrt3},\qquad
 \phi_T=\frac1{\sqrt6}\sum_{\pi\in S_3}|\pi(012)\rangle,\qquad
 Q=|\phi_F\rangle\langle\phi_F|+|\phi_T\rangle\langle\phi_T|.
\]
Its coherent-state symbol is
\begin{equation}
 q(z)=\langle z^{\otimes3},Qz^{\otimes3}\rangle
 =\frac13|z_0^3+z_1^3+z_2^3|^2+6|z_0z_1z_2|^2.
 \label{geo:q}
\end{equation}
The two cubic equations cut out precisely the nine rays \(\mathcal H\).

\fig{fig02_cubic_witness}{A positive cubic witness locates the Hesse
ensemble. (a) The exact symbol $q$ on the slice
$z=(0,\cos s,e^{i\theta}\sin s)$ has three isolated zeros at equal
coordinate moduli and the Hesse phases. (b) For 2,500 independent complex
Gaussian rays normalized to unit length (fixed reproducible seed), the
distance to the nearest Hesse ray is plotted against $q$. The solid curve
is the proved bound $\min\{1,\sqrt{27q/4}\}$; the samples illustrate the
bound and do not establish it. The inset shows the same symbol near one
zero as a surface, with the plane $q=0.03$ for scale.}{fig:witness}

\begin{theorem}[Robust Hesse ensemble certificate]
\label{geo:hesse-bound}
For any Borel probability measure \(\mu\) on \(\mathbb P^2\), put
\[
 \epsilon_Q=\operatorname{tr}(QM_3(\mu)),\qquad
 \delta_1=\|M_1(\mu)-I_3/3\|_{\rm HS},\qquad K=1+\frac3{\sqrt2}.
\]
Then
\begin{equation}
 W_1(\mu,\mu_{\mathcal H})
 \le K\sqrt{\frac{27}{4}\epsilon_Q}+\frac32\delta_1.
 \label{geo:robust}
\end{equation}
In particular, \(M_1(\mu)=I_3/3\) and \(\epsilon_Q=0\) characterize
\(\mu_{\mathcal H}\) among all probability ensembles.
\end{theorem}

\begin{proof}
First we prove the global error bound
\begin{equation}
 d(z,\mathcal H)^2\le\frac{27}{4}q(z).
 \label{geo:pointwise}
\end{equation}
For a normalized binary vector \((u,v)\), let \(\Delta\le\pi/3\) be its
phase error from the nearest vector \((1,-\omega^j)/\sqrt2\).
Put \(t=2|uv|\), \(y=t\cos\Delta\). Its squared distance to this
three-point set is \((1-y)/2\), whereas
\[
 |u^3+v^3|^2
 =(1-y)\left(1+y+y^2-\frac34t^2\right).
\]
The second factor is at least one, since
\[
 y+y^2-\frac34t^2
 =(y-t/2)(y+t/2+1)+\frac12t(1-t)\ge0.
\]
Thus binary squared distance is at most \(|u^3+v^3|^2/2\).

For a qutrit let \(a\ge b\ge c\) be its squared coordinate moduli and
\(f=|\sum_i z_i^3|\). Delete the minimum coordinate and use the binary
estimate. Since \(c\le1/3\),
\[
 d(z,\mathcal H)^2
 \le c+\frac{(f+c^{3/2})^2}{2(1-c)^2}
 \le\frac54c+\frac94f^2.
\]
If \(a\le3/4\), then \(ab\ge3/32\), and this is at most
\(40abc/3+9f^2/4\le(27/4)q(z)\).
If \(a\ge3/4\), then
\[
 f\ge a^{3/2}-(1-a)^{3/2}\ge2a-1,\qquad
 d(z,\mathcal H)^2\le\frac{3-a}{4}.
\]
The identity
\(9(2a-1)^2-(3-a)=(4a-3)(9a-2)\ge0\) again proves
\eqref{geo:pointwise}.

Map each ray to a nearest point of \(\mathcal H\), obtaining weights
\(w_i\), and put \(L^2=\int d(z,\mathcal H)^2\,d\mu\).
This transport costs at most \(L\), and the first-moment error of the
transported measure is at most \(\delta_1+\sqrt2L\).
For \(u_i=w_i-1/9\), the SIC Gram matrix gives
\[
 \left\|\sum_i u_iP_i\right\|_{\rm HS}^2=\frac34\sum_i u_i^2.
\]
The weights' total-variation distance from uniform is at most
\(\sqrt3\,\delta_1+\sqrt6\,L\).
All pairwise Hesse distances are \(\sqrt3/2\), so correcting the weights
costs at most \((3/2)\delta_1+(3/\sqrt2)L\).
Combine the two transports and \(L^2\le27\epsilon_Q/4\).
\end{proof}

All measurements in the theorem are positive operators.
The Hesse third moment has the particularly simple form
\begin{equation}
 M_3(\mu_{\mathcal H})=\frac18(I_{10}-Q).
 \label{geo:hesse-third}
\end{equation}
The vectors \(h^{\otimes3}\) form a nine-vector equiangular tight frame
in the eight-dimensional kernel of \(Q\). Algebraically, the ideal
\((z_0^3+z_1^3+z_2^3,z_0z_1z_2)\) is a complete intersection, with
Hilbert series
\[
 \frac{(1-t^3)^2}{(1-t)^3}.
\]
Its Hilbert function starts \(1,3,6,8,9,9,\ldots\), giving the ranks
available to coherent tensors at successive replica numbers.
At three replicas the density matrix has rank eight, yet its
decomposition into pure coherent tensors is uniquely the nine-point
Hesse mixture. Indeed, the coherent mixtures annihilated by \(Q\) are
supported on \(\mathcal H\), and the Gram matrix of their nine
projectors is
\[
 \frac{63I_9+\mathbf1\mathbf1^{\mathsf T}}{64},
\]
so those projectors are linearly independent. They span an exposed
simplex face of the convex set of bosonic separable states.

Three matched replicas are also necessary for distribution-free
identification. Haar measure and \(\mu_{\mathcal H}\) have equal first
and second moments, so any protocol with independent request blocks of
size at most two receives identical quantum states in the two cases,
including joint processing of all blocks. At order three,
Haar has \(M_3=I_{10}/10\) and \(Q\)-probability \(1/5\), whereas Hesse
has \(Q\)-probability zero. The degree of the defining equations has
therefore acquired a direct experimental meaning.

\paragraph{Finite-sample certification.}
On \(N_3\) triple requests measure \(\{Q,I-Q\}\), observing \(s\) clicks.
Let \(q_U\) be the one-sided Clopper--Pearson upper bound at failure
probability \(\alpha/2\). For \(s<N_3\), it is the
\(1-\alpha/2\) quantile of \(\operatorname{Beta}(s+1,N_3-s)\); if
\(s=N_3\), put \(q_U=1\). In particular, zero clicks give
\[
 q_U=1-(\alpha/2)^{1/N_3}.
\]
On \(N_1\) single requests measure the SIC POVM \(\{P_h/3\}\),
with empirical frequencies \(\widehat p_h\), and set
\[
 \widehat\delta_1=\sqrt{12}
 \left\|\widehat p-\frac19\mathbf1\right\|_2,\qquad
 u_{N_1}=
 \frac{\sqrt{12}\bigl(1+\sqrt{\log(2/\alpha)}\bigr)}{\sqrt{N_1}}.
\]
Multinomial mean-square error and bounded differences imply, with
confidence at least \(1-\alpha\),
\begin{equation}
 W_1(\mu,\mu_{\mathcal H})
 \le\min\left\{1,\,
 K\sqrt{\frac{27}{4}q_U}
 +\frac32(\widehat\delta_1+u_{N_1})\right\}.
 \label{geo:confidence}
\end{equation}
The inverse SIC relation is \(M_1=4\sum_h p_hP_h-I\).
For example, the sufficient budgets
\[
 N_3\ge27K^2\varepsilon^{-2}\log(2/\alpha),\qquad
 N_1\ge432\varepsilon^{-2}
                 \bigl(1+\sqrt{\log(2/\alpha)}\bigr)^2
\]
make the acceptance rule \(s=0\),
\(\widehat\delta_1\le\varepsilon/6\) sound and complete for the ideal
Hesse ensemble at confidence \(1-\alpha\). The number of physical
copies is \(3N_3+N_1\).

The accuracy exponent \(\varepsilon^{-2}\) is optimal for fixed request
patterns. For \(0\le\theta\le\pi/6\), rotate the ensemble by
\(U_\theta=\operatorname{diag}(e^{i\theta},e^{-i\theta},1)\). Then
\[
 W_1(\mu_{\mathcal H},U_{\theta*}\mu_{\mathcal H})
 =\frac{\sin\theta+2\sin(\theta/2)}3,\qquad
 \epsilon_Q=
 \frac{\sin^2(3\theta)+2\sin^2(3\theta/2)}{18}.
\]
Coupling hidden labels and using tensor-product fidelity bounds the
trace distance of the two \(n\)-copy experiments by \(\sqrt n\sin\theta\).
Constant-success testing at ensemble distance \(\varepsilon\) thus
requires \(n=\Omega(\varepsilon^{-2})\). This assertion concerns
fixed request patterns; it makes no claim about every adaptive
stopping convention.

\paragraph{A portable principle.}
Let a finite ray set \(S=\{h_1,\ldots,h_N\}\) be cut out by degree-\(k\)
equations whose coefficient space has orthogonal projector \(Q\).
Suppose a certified global bound is available:
\[
 d(z,S)^2\le
 C\langle z^{\otimes k},Qz^{\otimes k}\rangle.
\]
Suppose also that the weight map
\(u\mapsto\sum_j u_jP_{h_j}^{\otimes r}\), restricted to
\(\sum_j u_j=0\), has least singular value \(s_r>0\). If the observed
\(r\)-th moment differs from that of target weights by \(\delta\)
in Hilbert--Schmidt norm, and the \(Q\)-probability is \(\epsilon\),
the same transport argument gives
\begin{equation}
 W_1(\mu,\mu_*)
 \le\sqrt{C\epsilon}
 +\frac{\operatorname{diam}(S)\sqrt N}{2s_r}
       \left(\delta+\sqrt{2rC\epsilon}\right).
 \label{geo:general-cert}
\end{equation}
Equations localize the support; moments condition its weights.
Transverse isolated zeros give local quadratic error bounds. An
effective global \(C\) additionally needs a Jacobian estimate near the
zeros and a positive minimum away from them. These are concrete
targets for exact algebra and certified numerics.

\subsection{A second design that obstructs geometric localization}
\label{geo:singer}

The first-moment proof used a coupling on well-conditioned projective
frames. A natural stronger localization at second moment would couple
nine points whose blowup has nef anticanonical class. Here \emph{nef}
means nonnegative intersection with every curve. The next example
shows that second-moment isotropy alone cannot enforce that geometry.

Let \(\zeta=e^{2\pi i/7}\), and set
\begin{equation}
 B=\{e_0,e_1,e_2\},\qquad
 T_j=\frac{(1,\zeta^j,\zeta^{3j})}{\sqrt3},\quad 0\le j<7,\qquad
 \mu_{\rm S}=\frac1{12}\sum_{b\in B}\delta_b+
                  \frac3{28}\sum_{j=0}^6\delta_{T_j}.
 \label{geo:singer-measure}
\end{equation}
This ten-point design is a known projective-toric construction
\cite[Example 4.2]{geo_iosue2024}. Its character exponents
\(\{0,2,6,1,3,4\}\) in the symmetric-square basis are distinct modulo
seven, giving \(M_2(\mu_{\rm S})=I_6/6\) by exact character cancellation.

\begin{theorem}[An exact coupling obstruction and repair threshold]
\label{geo:singer-nef}
There is no coupling with all marginals \(\mu_{\rm S}\) supported on
nine distinct points whose blowup of \(\mathbb P^2\) has nef
anticanonical class. For
\[
 \mu_t=t\mu_{\rm S}+(1-t)\sigma,
\]
such a coupling exists if and only if \(t\le20/21\).
The coupling may always be taken exchangeable.
\end{theorem}

\begin{proof}
The cubics
\[
 F_0=x^3-yz^2,\qquad F_1=y^3-x^2z,\qquad F_2=z^3-xy^2
\]
vanish at all \(T_j\). Each omits one basis point and has a cusp at
another. A nine-point subset obtained by deleting a basis point thus
has a cubic through eight points simply and one doubly; its strict
transform has anticanonical intersection \(9-10=-1\).
Deleting a torus point instead leaves nine points on a smooth cubic.
For deletion of \(T_0\), an example is
\[
 x^2y+x^2z+xy^2+y^2z+xz^2+yz^2+xyz=0.
\]
Its gradient differences are
\[
 G_x-G_y=(y-x)(x+y-z),\quad
 G_y-G_z=(z-y)(y+z-x),\quad
 G_z-G_x=(x-z)(z+x-y).
\]
At a common gradient zero, if \(x=y=z\), a gradient entry is \(7x^2\).
If exactly two coordinates agree, say \(x=y\ne z\), the remaining
differences force \(z=0\), but then \(G_z=3x^2\ne0\).
If all three are distinct, the differences force
\(x+y=z,\ y+z=x,\ z+x=y\), hence the zero vector.
The cubic is therefore smooth, and
the diagonal order-seven symmetry handles all other deletions.
A smooth cubic through nine points has a strict transform of square
zero; its intersections with every distinct irreducible curve are
nonnegative, so it is nef.

Thus a nef nine-subset of the ten-point support contains all three
basis points. Its expected number of basis points would be three,
whereas its prescribed marginals give \(9\mu_{\rm S}(B)=9/4\).

For a tuple with marginal \(\mu_t\), let \(n_T\) count entries in the
seven-point torus orbit, and \(n_O\) entries outside the ten-point
support. Every admissible tuple satisfies \(n_T\le6+n_O\): when
\(n_O=0\) it must omit a torus point, and otherwise distinctness gives
the assertion. Taking expectations yields
\[
 \frac{27t}{4}\le6+9(1-t),\qquad t\le\frac{20}{21}.
\]
At equality, mix two types and randomly permute the tuple:
with probability \(4/7\), all three basis points and six uniformly
chosen torus points; with probability \(3/7\), all seven torus points,
one uniform basis point, and one Haar point.
For the second type, cyclic symmetry reduces to the retained basis
point \(e_0\). Cubics through \(T\cup\{e_0\}\) form the pencil
\(\operatorname{span}(F_1,F_2)\), since the seven cubic evaluation
characters exhaust all residues modulo seven.
This pencil contains the smooth member
\[
 F_1+F_2=y^3-x^2z+z^3-xy^2.
\]
Indeed, its \(y\)-derivative is \(y(3y-2x)\). If \(y=0\), the other
gradient equations force the zero vector. Otherwise \(y=2x/3\),
the \(x\)-derivative forces \(z=-2x/9\), and the \(z\)-derivative is
\((4/27-1)x^2\ne0\). Only finitely many pencil members are singular,
so a Haar point chooses a smooth member almost surely.
The mixture has weights \(5/63\) at each basis point, \(5/49\) at
each torus point, and Haar weight \(1/21\), exactly \(\mu_{20/21}\).
For smaller \(t\), mix with independent Haar points, whose generic
nine-point configuration lies on a smooth cubic.
\end{proof}

Two measures with identical second moments can therefore have
different coupling geometry. The obstruction is not explained by a
large atom: every atom of \(\mu_{\rm S}\) has mass below \(1/9\).
It is encoded in the incidence and multiplicities of the three cubics.

\subsection{The same design produces a singular K3 surface}
\label{geo:k3}

Blow up the ten points and call the rational surface \(S\). The strict
transforms \(C_0,C_1,C_2\) of the preceding cubics are disjoint smooth
rational curves with
\[
 C_i^2=-3,\qquad C_iC_j=0\ (i\ne j),\qquad
 C_0+C_1+C_2=-3K_S,\qquad K_S^2=-1.
\]
The cyclic triple cover \(\pi:Z\to S\), using the line bundle
\(-K_S\) and branch divisor \(\sum C_i\), is smooth: locally it has
equation \(w^3=u\). Its ramification curves \(R_i\) have square
\(-1\), and the canonical bundle formula gives
\(K_Z=R_0+R_1+R_2\). Contract them to obtain \(X\).
Then \(K_X\) is trivial, and
\[
 e(Z)=3e(S)-2\sum_ie(C_i)=3\cdot13-2\cdot6=27,\qquad e(X)=24.
\]
Noether's formula gives \(\chi(\mathcal O_X)=2\); connectedness and
triviality of \(K_X\) imply \(h^1(\mathcal O_X)=0\).
Thus \(X\) is a smooth projective K3 surface.
The term \emph{singular K3} below means maximal Picard rank twenty.

\fig{fig05_singer_cover}{From a finite projective design to a K3 surface.
(a) The three coordinate rays and seven Singer rays are drawn as an
incidence schematic; their positions in the drawing are not a projective
embedding. (b) Blowing up the ten rays separates the three cubic transforms
into disjoint $(-3)$-curves. The triple cover makes their ramification
curves $(-1)$-curves; contracting these gives a smooth K3 surface with
Euler number 24. (c) The transcendental lattice has the Eisenstein shape
$A_2(7)$, Gram matrix $\left(\begin{smallmatrix}14&-7\\-7&14\end{smallmatrix}\right)$,
and discriminant 147. The arrows in (b) describe actual geometric
operations, while (c) is a lattice diagram.}{fig:k3}

\begin{theorem}[The lattice of the design cover]
\label{geo:k3-lattice}
The K3 surface obtained from this ten-point configuration satisfies
\[
 \rho(X)=20,\qquad
 T_X\cong A_2(7)=
 \begin{pmatrix}14&-7\\-7&14\end{pmatrix},\qquad
 \operatorname{disc}T_X=147.
\]
It admits a symplectic automorphism group \(C_7\rtimes C_3\) and a
commuting non-symplectic deck transformation of order three with
exactly three isolated fixed points.
\end{theorem}

\begin{proof}
A weighted-projective model is
\[
 w^3=(x^3-yz^2)(y^3-x^2z)(z^3-xy^2),\qquad
 \deg(x,y,z,w)=(1,1,1,3).
\]
On its smooth locus a two-form is
\[
 \Omega=\frac{x\,dy\wedge dz-y\,dx\wedge dz+z\,dx\wedge dy}{w}.
\]
At each blown-up point the total branch multiplicity is three, so
the blowup removes its apparent poles. The zeros are exactly the
ramification curves, removed by their contraction.
The deck map \(\delta:w\mapsto e^{2\pi i/3}w\) acts nontrivially on
\(\Omega\); its three fixed curves become three isolated fixed points.
The maps
\[
 g:(x,y,z,w)\mapsto(x,\zeta y,\zeta^3z,\zeta^4w),\qquad
 c:(x,y,z,w)\mapsto(y,z,x,w)
\]
preserve \(\Omega\), and satisfy \(g^7=c^3=1\),
\(cgc^{-1}=g^2\) in weighted projective coordinates.

A symplectic order-seven automorphism has rank-four invariant
cohomology and a rank-eighteen algebraic coinvariant lattice
\cite[Propositions 4.5--4.6]{geo_garbagnati_sarti2007}.
An invariant ample class therefore gives \(\rho(X)\ge19\).
The non-symplectic order-three action makes the transcendental
lattice a module over \(\mathbb Z[e^{2\pi i/3}]\), so its rank is
even. It is at most three and at least two, hence is two.
Thus \(\rho(X)=20\). Since the Eisenstein ring is principal, its
invariant positive even rank-two form is \(A_2(m)\) for an integer \(m>0\).

To determine \(m\), let
\[
 L_0=\{D\in\operatorname{Pic}(S):D C_i=0\text{ for every }i\}.
\]
The primitive closure of \(\langle C_0,C_1,C_2\rangle\) has index
three: it contains \(K_S=-(C_0+C_1+C_2)/3\), and its index squared
divides \(27\). Unimodularity of \(\operatorname{Pic}(S)\) gives
\(|\operatorname{disc}L_0|=3\).
Pullback scales its form by three. The fixed lattice of \(\delta\)
has rank eight and discriminant \(3^7\), since its fixed locus is
three isolated points \cite[Theorem 4.1]{geo_artebani_sarti_taki2011}.
Hence \(\pi^*L_0\) has index three in that fixed lattice.

Let \(H\) and \(E_b,E_{T_j}\) be the line and exceptional classes on \(S\).
The \(g\)-fixed part of \(L_0\) has basis
\[
 v_1=H-\sum_{b\in B}E_b,\qquad
 v_2=7H-3\sum_{j=0}^6E_{T_j},\qquad
 \bigl(v_iv_j\bigr)=\begin{pmatrix}-2&7\\7&-14\end{pmatrix}.
\]
The order-seven action on the index-three quotient is trivial, and
the obstruction in \(H^1(C_7,\pi^*L_0)\) is seven-torsion. Thus
taking \(g\)-invariants preserves index three.
The unique even index-three overlattice of the threefold-scaled
matrix adds \((v_2-v_1)/3\): an order-three discriminant element
\((\alpha v_1+\beta v_2)/3\) is isotropic exactly when
\(\alpha+\beta=0\) modulo three.
In the basis \(v_1+(v_2-v_1)/3,-(v_2-v_1)/3\) its matrix is
\[
 N=\begin{pmatrix}2&1\\1&-10\end{pmatrix}.
\]
A symplectic automorphism fixes \(T_X\) pointwise, since its rational
coinvariant classes are of type \((1,1)\).
The rank-four invariant lattice therefore has rank-two algebraic part.
This contains the rank-two lattice just computed in the \(\delta\)-fixed
part; consequently \(\delta\) acts trivially on its rational span, and
\(\operatorname{NS}(X)^g=N\).

The full invariant lattice is
\[
 H^2(X,\mathbb Z)^g
 \cong U(7)\oplus\begin{pmatrix}4&1\\1&2\end{pmatrix},
 \qquad |\operatorname{disc}|=7^3,
\]
with elementary discriminant group \((\mathbb Z/7)^3\)
\cite[Proposition 4.6]{geo_garbagnati_sarti2007}.
Both \(N\) and \(T_X\) are primitive in this lattice.
If \(h\) is their gluing index, the gluing group injects into both
discriminant groups, so \(h\mid21\). Determinants give
\[
 7^3=\frac{21\cdot3m^2}{h^2}.
\]
The only possibilities are \((m,h)=(7,3),(49,21)\).
In the latter case, the seven-primary discriminants of \(N,T_X\)
are \(\mathbb Z/7\) and \((\mathbb Z/49)^2\). The gluing subgroup
has generator \((n,7v)\) with \(n\ne0\).
Its orthogonal complement projects onto \((\mathbb Z/49)^2\),
since every pairing with \(7v\) can be canceled by a multiple of \(n\).
Lift an element of order \(49\). Seven times that lift has zero first
component and nonzero second component; it cannot belong to the
gluing subgroup. The glued discriminant would retain an element of
order \(49\), contradicting its elementary form. Therefore \(m=7\).
\end{proof}

The invariant lattice also makes the relevant nef geometry explicit.
Its norm is \(2(x^2+xy-5y^2)\), congruent to \(2(x-y)^2\) modulo three,
so it represents neither \(-2\) nor \(-14\).
A rational curve on \(X\) has square \(-2\); its \(g\)-orbit cannot
have length one. A length-seven orbit sum has square \(-14+14k\)
for an integer \(k\ge0\), and cannot have square \(-14\), hence is
nonnegative. Classes in the positive component of \(N_{\mathbb R}\)
therefore pair positively with every such orbit sum, and thus with
every curve, by the Hodge index theorem. The invariant nef cone is
the closure of that positive component. In invariant blowup
coordinates
\[
 dH-b\sum_{b\in B}E_b-c\sum_jE_{T_j},\qquad
 3d=3b+7c,
\]
its null boundary rays are \(d/b=(9\pm\sqrt{21})/2\).

The construction does not assert a new K3 isomorphism class. Its
content is the explicit passage from an isotropic quantum ensemble,
through a geometric localization obstruction, to a canonical cover
whose transcendental lattice is determined. The same cubic
multiplicities that obstruct nef localization make the cyclic cover
possible.

\subsection{Verification and reproducible mathematical content}
\label{geo:verification}

The conceptual inputs are Kempf--Ness, the three-qutrit Cartan
classification, semicontinuity of singularity exponents,
Hilbert/Bergman geodesic comparison, and the stated K3 automorphism
lattices. They enter with precise hypotheses and primary references.
Finite calculations are supplied as exact-arithmetic certificates:
Moore determinants and quartic contractions; Hesse moments, Gram
matrices and spectra; Singer character cancellation, cubic
smoothness and intersections; and finite minimal-jet moment-map
examples. Numerical sampling illustrates bounds and searches for
counterexamples. The theorem constants follow from the stated
inequalities and identities.

The operational supplement evaluates \eqref{geo:confidence} from
measurement counts and supplies seeded synthetic-data examples.
It exposes the equation-to-projector and moment-to-weight-conditioning
steps of \eqref{geo:general-cert} for reuse with other finite ray
sets. The computations verify finite identities and statistical
calculations; they are not a completed machine formalization of
the analytic and geometric proofs.



\section{Moment repair and constrained equilibria}
\label{sec:equilibria}

The geometric threshold becomes more useful when the constraint can be
enforced stably and its equilibria can be described by an equation.
This section provides those analytic interfaces on $X=\mathbb P^2$.
It also explains which part of the determinant problem is controlled by
ordinary compactness and which part requires information about very small
singular values.

We use $\omega_0=3\omega_{\rm FS}$, $V=\int\omega_0^2=9$,
$\sigma=\omega_0^2/V$, and $dd^c=i\partial\bar\partial/(2\pi)$.
Write $\mathcal A=\mathcal E_L$ and $E=E_L$ for the energies in
\eqref{geo:energy}, and $H(\mu)=\Ent(\mu\mid\sigma)$. A positive
K\"ahler potential changes the reference metric by
$\omega_u=\omega_0+dd^cu>0$. Its normalized Monge--Amp\`ere measure is
$\operatorname{MA}(u)=\omega_u^2/V$: in coordinates it is the determinant
of the Hermitian metric matrix, normalized to mass one. The energy $E$ is
the convex dual of the concave potential energy $\mathcal A$.

The finite-energy completion $\mathcal E^1$ allows singular potentials
with full Monge--Amp\`ere mass and finite Aubin--Mabuchi energy. We use the
variational theorem of \citet{BBGZ2013}: a finite-energy measure has a
normalized potential attaining the dual energy. We also use the strong
compactness of entropy sublevels and continuity of the inverse
Monge--Amp\`ere map from \citet[Theorem 2.17 and Corollary 2.19]{BBEGZ2019}.
These results supply the analytic completion in which the elementary
finite-dimensional moment correction can be used.

\subsection{Quantitative stability of the constraint}

We first record a uniform integrability estimate. If $u$ is
$\omega_0$-plurisubharmonic and $\sup u=0$, then
\begin{equation}
 \int e^{-\alpha u}\,d\sigma\le
 \frac{2}{(1-3\alpha)(2-3\alpha)},\qquad 0\le\alpha<1/3.
 \label{an:alpha-baseline}
\end{equation}
Here is an elementary derivation of the constant. Fix a maximum point
$a$ and disintegrate over complex lines through $a$, writing
$z=\sqrt{x}\,a+e^{i\theta}\sqrt{1-x}\,b$. The conditional density is
$2(1-x)dx\,d\theta/(2\pi)$. The line potential $u/3$ is a positive
average of logarithms of normalized linear forms plus a nonnegative
constant, by the Riesz representation and its value at $a$. Jensen
reduces the estimate to one form $\ell$.
If $b_0=|\ell(a)|^2$ and $s<1$, direct integration gives
\[
 \int|\ell|^{-2s}d\sigma_{\text{line},a}
 =\frac{2(1-s+sb_0)}{(1-s)(2-s)}
 \le\frac2{(1-s)(2-s)}.
\]
Setting $s=3\alpha$ proves \eqref{an:alpha-baseline}. Its derivative
at zero gives $\int(-u)d\sigma\le9/2$. Thus a mean-zero potential
satisfies $\int|u|d\sigma\le9$. Entropy duality, with $\alpha=1/6$,
also gives the global bound
\begin{equation}
 E(\mu)\le6H(\mu)+6\log(8/3).
 \label{an:entropy-to-energy}
\end{equation}
This ensures finite energy for every finite-entropy measure, before any
moment constraint is imposed.

\begin{proposition}[Energy retained by positive dilution]
\label{an:dilution}
Let $0<p\le1$, let $1/2\le d\nu/d\sigma\le3/2$, and put
\[
 c=c(p)=\frac{3-\sqrt{9-8p}}2,\qquad p=\frac{c(3-c)}2.
\]
For every probability $\mu$, allowing infinite energy,
\begin{equation}
 E(p\mu+(1-p)\nu)\ge pcE(\mu)-\frac92c(1-p).
 \label{an:dilution-bound}
\end{equation}
\end{proposition}

\begin{proof}
For a smooth test normalized by $\int u\,d\sigma=0$, write
$J(u)=-\mathcal A(u)$ and
\[
 a=V^{-1}\int du\wedge d^cu\wedge\omega_0,\qquad
 b=V^{-1}\int du\wedge d^cu\wedge\omega_u.
\]
Both are nonnegative. Integration by parts gives
$J(u)=(2a+b)/3$ and $J(cu)=c^2[(3-c)a+cb]/3$. Therefore
$pcJ(u)-J(cu)=c^2(1-c)b/2\ge0$. Also
$|\int u\,d\nu|\le\tfrac12\int|u|d\sigma\le9/2$.
Testing the diluted measure at $cu$ and taking the supremum proves
\eqref{an:dilution-bound}, including the infinite case.
\end{proof}

Suppose $H\ge qE-K$ on the exact second-design fibre, with $q>0$.
For every $\gamma<q$ this inequality persists, with a smaller coefficient
still above $\gamma$, on an explicit moment neighborhood. Indeed choose
$\delta_0>0$ so that, with $p_0=(1+102\delta_0)^{-1}$,
$q_0=q\,c(p_0)>\gamma$. Positive repair
(Proposition~\ref{prop:repair}) and \eqref{an:dilution-bound} give
\begin{equation}
 H(\mu)\ge q_0E(\mu)-K_1,
 \qquad
 K_1=\frac{K+\log(3/2)+9q/2}{p_0},
 \label{an:robust-coercivity}
\end{equation}
whenever $\|M_2(\mu)-I_6/6\|_{\rm HS}\le\delta_0$ and $H(\mu)<\infty$.
In particular $H-\gamma E\ge(1-\gamma/q_0)H-\gamma K_1/q_0$ there.
Theorem~\ref{geo:sharp} supplies every $q<12/7$, since every second
design is balanced. Thus this neighborhood result and the equilibrium
theorems below hold unconditionally for $\gamma<12/7$.

\subsection{Entropy projection and the finite multiplier}

A constrained equilibrium has an infinite-dimensional potential but only
a finite-dimensional correction enforcing its moments. This correction
is an exponential-family parameter, determined by a strictly convex
optimization problem.

Let $\mathfrak a=\operatorname{Herm}_0(\Sym^2\mathbb C^3)$ and
$h_Y(z)=\tr(YQ_z)$, with $Q_z=P_z^{\otimes2}$. For
$u\in\mathcal E^1$, $\sup u=0$, a bounded real $W$, and $\gamma\ge0$,
define
\[
 J_u(Y)=\log Z_u(Y),\qquad
 Z_u(Y)=\int e^{-\gamma u-W+h_Y}\,d\sigma.
\]

\begin{theorem}[Exact constrained entropy projection]
\label{an:entropy-projection}
The function $J_u$ has a unique minimizer $Y_u\in\mathfrak a$.
The probability
\begin{equation}
 d\eta_u=Z_u(Y_u)^{-1}e^{-\gamma u-W+h_{Y_u}}d\sigma
 \label{an:exponential-family}
\end{equation}
is an exact second design of finite entropy and energy. It uniquely
minimizes $H(\nu)+\int(\gamma u+W)d\nu$ among finite-entropy measures
on that fibre, and
\begin{equation}
 \Ent(\nu\mid\eta_u)
 =H(\nu)+\int(\gamma u+W)d\nu+\log Z_u(Y_u)
 \label{an:pythagoras}
\end{equation}
for every finite-entropy second design $\nu$.
\end{theorem}

\begin{proof}
Full-mass finite-energy potentials have $e^{-a u}\in L^1(\sigma)$
for every finite $a$ \citep[Proposition 1.4]{BBEGZ2019}. Entropy duality
then makes all potential pairings in \eqref{an:pythagoras} finite.
The finite-dimensional parameter derivatives are bounded coherent symbols,
so $J_u$ is smooth by dominated differentiation on bounded parameter sets.

The coherent frame in \eqref{an:frame-spectrum} has least eigenvalue
$1/90$. If $\|Y\|_{\rm HS}=1$, then $\int h_Y=0$, $|h_Y|\le1$, and
$\int h_Y^2\ge1/90$. Thus $\int(h_Y)_+\ge1/180$ and
$\sigma\{h_Y\ge1/360\}\ge1/360$. Since $u\le0$,
\begin{equation}
 J_u(Y)\ge\frac{\|Y\|_{\rm HS}}{360}-\|W\|_\infty-\log360.
 \label{an:multiplier-coercivity}
\end{equation}
This proves existence. Its Hessian is
$D^2J_u(Y)[A,B]=\operatorname{Cov}_{\eta_Y}(h_A,h_B)$.
The density is positive almost everywhere; zero variance would make
$h_A$ constant everywhere, hence zero by its Haar mean, and then $A=0$
by frame invertibility. Strict convexity gives uniqueness.
The stationarity condition is $M_2(\eta_u)=I_6/6$.
Exponential integrability gives finite entropy, and
\eqref{an:entropy-to-energy} gives finite energy. Finally substitute
\eqref{an:exponential-family} into relative entropy; the multiplier term
vanishes on the design fibre. Nonnegativity and the equality condition
of relative entropy prove the minimization assertion.
\end{proof}

Comparison with $J_u(0)$ gives the useful explicit bound
\begin{equation}
 \|Y_u\|_{\rm HS}\le360\left(
 2\|W\|_\infty+\log\int e^{-\gamma u}d\sigma+\log360\right).
 \label{an:multiplier-bound}
\end{equation}
This finite covariance Hessian is available to numerical solvers. The
positive repair radius guarantees an interior moment target, while
\eqref{an:multiplier-coercivity} prevents the multiplier from escaping
to infinity.

\subsection{Existence and smoothness of attractive macroscopic equilibria}

Define $\mathcal D_2=\{\mu:M_2(\mu)=I_6/6\}$ and, for
$0\le\gamma<12/7$,
\[
 \mathcal F_{\gamma,W}(\mu)=H(\mu)-\gamma E(\mu)+\int W\,d\mu
\]
on finite-entropy measures in $\mathcal D_2$, with value $+\infty$
elsewhere. The minus sign makes the interaction attractive.

\begin{theorem}[Smooth constrained equilibria]
\label{an:equilibrium}
For $0\le\gamma<12/7$ and $W\in C^\infty(X,\R)$,
$\mathcal F_{\gamma,W}$ attains its minimum. Every minimizer is smooth
and strictly positive. Its Monge--Amp\`ere potential, normalized by
$\sup_Xu=0$, and unique traceless multiplier satisfy
\begin{equation}
 \operatorname{MA}(u)=
 \frac{e^{-\gamma u-W+h_Y}}
 {\int e^{-\gamma u-W+h_Y}d\sigma}\,\sigma.
 \label{an:equilibrium-pde}
\end{equation}
For $\gamma$ in a compact subinterval of $[0,12/7)$ and $W$ in a bounded
subset of $C^\infty$, the potentials, densities, and multipliers have
uniform bounds in every smooth norm, and the densities have a common
positive lower bound. The minimizer itself is not asserted to be unique.
\end{theorem}

\begin{proof}
Choose $q$ strictly between the largest considered $\gamma$ and $12/7$.
Coercivity gives $H-\gamma E\ge(1-\gamma/q)H-\gamma K/q$ on
$\mathcal D_2$. Hence a minimizing sequence has bounded entropy.
Entropy sublevels are compact for weak convergence together with
convergence of $E$ \citep{BBEGZ2019}. The limiting measure remains in
the closed moment fibre, and lower semicontinuity of $H$ proves existence.

Let $u$ be the finite-energy potential of a minimizer $\mu$. Energy
duality gives $E(\nu)\ge\mathcal A(u)-\int u\,d\nu$, with equality at
$\mu$. Therefore
\[
 \mathcal F_{\gamma,W}(\nu)
 \le H(\nu)+\int(\gamma u+W)d\nu-\gamma\mathcal A(u),
\]
again with equality at $\mu$. If $\mu$ differed from the entropy projection
$\eta_u$, identity \eqref{an:pythagoras} would strictly lower this
majorizer, and hence the original functional. This contradiction proves
\eqref{an:equilibrium-pde} without differentiating energy on singular
measures.

For uniform estimates, comparison with $\sigma$ bounds the entropy of
all minimizers. Their normalized potentials form a strongly compact
family. Exponential integrability is uniform on this family, by
\citet[Proposition 1.4]{BBEGZ2019}. Equation
\eqref{an:multiplier-bound} bounds $Y$; its coherent symbol then has
uniform derivatives of every order. The densities in
\eqref{an:equilibrium-pde} are uniformly bounded in every finite $L^p$
and bounded below. The H\"older estimate of
\citet[Theorem A*]{DemaillyEtAl2014} gives a uniform $C^{\alpha_0}$ bound
on $u$, for example with $\alpha_0=1/4$.

To justify higher regularity before applying smooth PDE estimates,
average $u$ against the $SU(3)$ heat kernel, obtaining smooth
$u_\epsilon\to u$ uniformly with $\omega_0+dd^cu_\epsilon\ge0$.
Put $F=h_Y-W$ and solve the smooth auxiliary equation
\[
 \operatorname{MA}(v_\epsilon)=
 \frac{e^{-\gamma u_\epsilon+F}}
 {\int e^{-\gamma u_\epsilon+F}d\sigma}\,\sigma,
 \qquad\sup v_\epsilon=0.
\]
The classical smooth solution and bounded-solution uniqueness are supplied
by \citet[Theorems A and B]{EGZ2009}, specialized to a positive class
and no divisorial singularities. The right sides have uniform positive
lower bounds and H\"older norms. In the Laplacian estimate of
\citet[Theorem B.1]{BBEGZ2019}, use
$\psi^+=F-\log\int e^{-\gamma u_\epsilon+F}d\sigma$ and
$\psi^-=\gamma u_\epsilon$. Their complex Hessians are bounded below,
$\psi^+$ is bounded above, and $e^{-\psi^-}$ has uniform $L^p$ bounds.
It follows that
$\omega_{v_\epsilon}\le C e^{-\gamma u_\epsilon}\omega_0\le C'\omega_0$.
The determinant equation gives the lower ellipticity bound.

The same $L^2$ density estimate gives a uniform $C^0$ bound for
$v_\epsilon$. The $C^{2,\beta}$ theorem of \citet[Theorem 1.1]{TWWY2015} now applies
to the already smooth $v_\epsilon$: their potentials, Laplacians, and
H\"older log densities are uniformly bounded. A convergent subsequence
has limit $u$, by uniqueness of the bounded normalized solution.
Thus $u$ is $C^{2,\beta'}$ for $0<\beta'<\beta$. Differentiating
\eqref{an:equilibrium-pde} gives a uniformly elliptic linear equation
for each first derivative; Schauder bootstrapping yields all higher
bounds. The positive lower bound on the density is retained throughout.
\end{proof}

The same compactness argument proves that the extended functional has
compact sublevels and that a minimizer exists after adding any continuous
functional of $\mu$. Smoothness above requires the displayed smooth
one-body tilt $\int Wd\mu$. There is also an exact potential-side dual:
\begin{equation}
 \inf_{\mathcal D_2}\mathcal F_{\gamma,W}
 =\inf_{u\in\mathcal E^1}
 \left[-\gamma\mathcal A(u)-\min_{Y\in\mathfrak a}J_u(Y)\right].
 \label{an:dual-equilibrium}
\end{equation}
Write $-\gamma E$ as an infimum of affine tests and interchange two
infima; Theorem~\ref{an:entropy-projection} then gives the formula.
The two sides are attained by a minimizing measure and its potential.

\subsection{Continuous energy envelopes and deterministic lower comparison}

The topology used above can be made useful for finite determinant
ensembles. Choose a countable $C^2$-dense sequence of smooth positive
potentials $u_0=0,u_1,\ldots$ and let
\begin{equation}
 E_r(\mu)=\max_{0\le j\le r}
 \left(\mathcal A(u_j)-\int u_j\,d\mu\right).
 \label{an:continuous-envelopes}
\end{equation}
These are continuous, nonnegative, and increase to $E$. Convergence is
uniform on every entropy sublevel by strong compactness and Dini's
theorem. Averaging measures over the unitary heat kernel preserves exact
design moments and decreases both entropy and energy. Lower
semicontinuity then shows that every finite-entropy exact design is
approximable by smooth strictly positive exact designs with convergence
of both functionals.

The smooth-test and finite-energy versions of the dual agree here.
Indeed, after fixing the heat-time normalization,
$\Delta_{\omega_0}u\ge-2$ implies that $S_tu+2t$ decreases to $u$ as
$t\downarrow0$. Its K\"ahler form is strictly positive, since the group
heat kernel has a positive minimum and the Haar average of the current
is $\omega_0$. Monotone continuity of energy and of potential integrals
then identifies the two suprema, including infinite values.

Let $N_k=\binom{3k+2}{2}$ and let $\Delta_k$ be the evaluation
determinant of an orthonormal basis of $H^0(X,\mathcal O(3k))$.
For a particle tuple put $L_k=N_k^{-1}\sum_i\delta_{z_i}$ and
$E_k=-(kN_k)^{-1}\log\|\Delta_k\|^2$.

\begin{theorem}[Lower comparison with explicit finite-test error]
\label{an:lower-comparison}
For each fixed finite collection in \eqref{an:continuous-envelopes}
there is a finite constant $C_r$ such that, on every configuration,
\begin{equation}
 E_k\ge E_r(L_k)-\frac{2\log k+C_r}{k},
 \qquad E_k\ge-\frac{\log N_k}{k}.
 \label{an:LC}
\end{equation}
The inequalities include determinant-zero configurations, where
$E_k=+\infty$.
\end{theorem}

\begin{proof}
Let $G_k(u)$ be the Hilbert matrix with weight $e^{-ku}$. Put
\[
 \mathcal L_k(u)=-\frac{1}{kN_k}\log\det G_k(u).
\]
The exact adjoint
Hilbert comparison is
$\mathcal L_k(u)\ge\mathcal A(u)$. To see the normalization directly,
join $0$ to $u$ by a weak geodesic $u_t$. The adjoint metric on
$kL-K_X$ has nonnegative curvature:
\[
 k(\pi^*\omega_0+dd^cu_t)+\pi^*\operatorname{Ric}(\sigma)\ge0,
 \qquad\operatorname{Ric}(\sigma)=\omega_0.
\]
Thus the maximum principle of
\citet[Proposition 4.1]{geo_zhang2024} gives
$G_k(u)^t\le G_k(u_t)$. Differentiate the logarithmic determinant at
$t=0$. The reference Bergman density is exactly $N_k$, and the
potential energy is affine along the geodesic. This proves the comparison
with no error in its energy coefficient.

The weighted Bergman function satisfies $\sup_XB_{k,u}\le C_uk^2$.
In a holomorphic chart, subtract the constant and real-linear part of the
metric weight; its remainder is bounded by $C|z|^2$. The holomorphic
submean inequality on a radius $k^{-1/2}$ ball gives the bound, uniformly
for potentials in a bounded $C^2$ set. Changing to a weighted orthonormal
basis and applying Hadamard to its evaluation determinant now yields
\[
 E_k\ge\mathcal L_k(u)-\int u\,dL_k
          -k^{-1}\log\sup_X B_{k,u}.
\]
Take a finite maximum. At $u=0$ the exact value $B_{k,0}=N_k$ gives
the second inequality.
\end{proof}

This lower comparison is deterministic, hence can be combined with
ordinary constrained Sanov theory on the repulsive side. Attractive
reweighting instead rewards configurations with very large $E_k$;
its pressure requires upper-tail information beyond the deterministic
inequality. The next section supplies a spectral description of the
coherent kernel, followed by the constrained repulsive Laplace principle.

\section{Directed trees and coherent-state spectral laws}
\label{sec:coherent-spectrum}

A projective coherent state carries considerably more structure than an
isotropic random vector in its ambient Hilbert space.  This distinction
survives when the number of samples and the Hilbert-space dimension grow
at the same rate.  We describe the resulting spectral law by a sum over
set partitions.  Its weights are directed spanning-tree counts.  The
noncrossing partitions familiar from free probability emerge precisely
when the complex dimension tends to infinity.

This is a spatial-kernel moment calculation, in the sense of the
Euclidean random-matrix framework of Bordenave~\cite{prob_bordenave2008}.
The coherent phase is retained: after rescaling it produces a directed,
rather than a symmetric, graph Laplacian.  Lowest-Landau-level random
impurity models provide a related one-complex-dimensional setting
\cite{prob_brezin1984}.  We give a direct proof on projective space,
including arbitrary fixed sampling density.

\subsection{The spectral law at fixed dimension and sampling density}

Fix a positive integer $q$ and write
\[
 D_m=\dim H^0(\mathbb P^q,\mathcal O(m))=\binom{m+q}{q}.
\]
Let $n=n_m$ satisfy $n_m/D_m\longrightarrow\rho\in(0,\infty)$.
Choose independent Fubini--Study distributed rays $z_1,\ldots,z_n$
in $\mathbb P^q$ and arbitrary unit lifts.  Their degree-$m$ coherent
Gram matrix and empirical spectral measure are
\begin{equation}
 B^{(m)}_{ij}=\langle z_i,z_j\rangle^m,
 \qquad
 \nu_m=\frac1n\sum_{j=1}^n\delta_{\lambda_j(B^{(m)})}.
 \label{eq:coherent-gram-definition}
\end{equation}
Changing the phases of the lifts gives a diagonal unitary conjugation.
In particular, the spectrum is intrinsic.  Moreover, $B^{(m)}\geq0$
and $B^{(m)}_{ii}=1$.

For a set partition $\pi$ of $\{1,\ldots,p\}$, identify the vertices
of the directed cycle $1\to2\to\cdots\to p\to1$ according to the
blocks of $\pi$.  Denote the resulting directed multigraph by $G_\pi$.
Parallel edges are retained.  Loops contribute nothing to its
outdegree Laplacian
\[
 L_{aa}=\sum_{b\ne a}n_{ab},\qquad L_{ab}=-n_{ab}\quad(a\ne b).
\]
The graph is connected and Eulerian.  Its rooted cofactor
$\tau(G_\pi)=\det L^{(o)}$ is the number of directed spanning
in-arborescences rooted at $o$.  It is a positive integer and is
independent of $o$.  The one-vertex cofactor is defined to be one.

\begin{theorem}[Coherent spectral law]
\label{thm:coherent-spectral-law}
There is a unique probability measure $\nu_{q,\rho}$ on $[0,\infty)$
whose positive integer moments are
\begin{equation}
 M_p(q,\rho)=
 \sum_{\pi\in\mathfrak P(p)}
       \rho^{|\pi|-1}\tau(G_\pi)^{-q}.
 \label{eq:coherent-tree-moments}
\end{equation}
For every fixed $p\geq1$,
\[
 \mathbb E\!\int x^p\,d\nu_m\longrightarrow M_p(q,\rho),
 \qquad
 \operatorname{Var}\!\left(\int x^p\,d\nu_m\right)=O_p(n_m^{-1}).
\]
Consequently $\nu_m$ converges weakly in probability to
$\nu_{q,\rho}$, with convergence in probability of every fixed
positive integer moment.  If $\sum_m n_m^{-1}<\infty$, these
convergences hold almost surely on any probability space carrying
the rows.  In particular this applies along all degrees when $q\geq2$.

The same assertions hold after conditioning on any fixed finite
collection of bounded empirical moments lying within distance
$\epsilon_m$ of their Fubini--Study expectations, provided
$n_m\epsilon_m^2\to\infty$.
\end{theorem}

The first four limiting moments are
\begin{align}
 M_1&=1,& M_2&=1+\rho,& M_3&=1+3\rho+\rho^2,\notag\\
 M_4&=1+(6+2^{-q})\rho+6\rho^2+\rho^3.
 \label{eq:coherent-first-moments}
\end{align}
At unit sampling density and complex dimension two they are
$1,2,5,57/4$.  The sole crossing partition of four positions,
$\{\{1,3\},\{2,4\}\}$, has two directed spanning trees and
accounts for the excess $2^{-q}\rho$ in the fourth moment.

\fig{fig06_gram_trees}{Directed trees distinguish coherent Gram spectra
from the free limit. (a) Identifying equal-colored vertices of the directed
four-cycle for the crossing partition gives two directed spanning trees.
(b) At unit sampling density the fourth limiting moment is
$14+2^{-d}$; complex dimension two gives $57/4$. (c) Exact integer
cofactor enumeration gives the relative excess $M_p/C_p-1$ over Catalan moments at orders
four through seven. The excess decreases as complex dimension grows,
while higher moments retain more crossing contributions at fixed
dimension. The figure uses exact partition formulas, without sampled
eigenvalue data; $d$ in its axes is the complex dimension denoted $q$ in
the text.}{fig:gram}

\begin{lemma}[The local graph integral]
\label{lem:coherent-graph-integral}
Let $G$ be a fixed connected Eulerian directed multigraph on $r$
vertices, with multiplicities $n_{ab}$.  Then
\begin{equation}
 I_m(G):=\int\prod_{a,b}\langle z_a,z_b\rangle^{m n_{ab}}
                   \,d\sigma(z_1)\cdots d\sigma(z_r)
 =D_m^{-(r-1)}\tau(G)^{-q}(1+o(1)).
 \label{eq:coherent-graph-integral}
\end{equation}
Furthermore $|I_m(G)|\leq C_GD_m^{-(r-1)}$ for all sufficiently
large $m$.
\end{lemma}

\begin{proof}
The Eulerian condition makes the integrand independent of the lifts.
Unitary invariance permits fixing $z_1=[1:0:\cdots:0]$.  The
integrand has modulus at most one, with equality only when all rays
coincide.  Indeed equality forces the endpoints of every edge to
coincide projectively, and the graph is connected.  Away from a
fixed neighborhood of this point, the integral is exponentially
small in $m$.

For the remaining rays use the lifts
$(1,w_a)/(1+|w_a|^2)^{1/2}$, with $w_1=0$.  The integrand near
zero is $\exp(m\Psi(w))$, where
\[
 \Psi(w)=\sum_{a,b}n_{ab}\log(1+\overline w_a\cdot w_b)
       -\sum_a\left(\sum_b n_{ab}\right)\log(1+|w_a|^2).
\]
The Eulerian identity equates indegrees with outdegrees and accounts
for the second sum.  For each of the $q$ complex coordinate
directions the quadratic term is $-\overline w^{\,T}L^{(1)}w$.
The Hermitian part of $L^{(1)}$ is the positive definite pinned
Laplacian with undirected edge weights $(n_{ab}+n_{ba})/2$.
Consequently $\operatorname{Re}\Psi(w)\leq-c\sum_a|w_a|^2$
on a sufficiently small neighborhood.

The affine Fubini--Study density is
$(q!/\pi^q)(1+|w|^2)^{-(q+1)}d^{2q}w$.
After $w_a=m^{-1/2}\xi_a$, Gaussian domination gives
\[
 I_m(G)=(q!)^{r-1}m^{-q(r-1)}\det(L^{(1)})^{-q}(1+o(1)).
\]
Here we used
$\int_{\mathbb C^s}e^{-\overline\xi^{\,T}A\xi}d^{2s}\xi
=\pi^s/\det A$ when the Hermitian part of $A$ is positive
definite.  This identity follows by unitary Schur triangularization
and successive one-variable Gaussian integration.  Since
$D_m\sim m^q/q!$, the asserted asymptotic follows.  The same
Gaussian domination and the exponentially small exterior contribution
give the uniform bound.
\end{proof}

\begin{proof}[Proof of Theorem~\ref{thm:coherent-spectral-law}]
Expanding the trace according to equality patterns of particle
indices gives the exact identity
\begin{equation}
 \mathbb E\frac1n\operatorname{tr}\bigl[(B^{(m)})^p\bigr]
 =\sum_{\pi\in\mathfrak P(p)}
       \frac{(n)_{|\pi|}}n I_m(G_\pi).
 \label{eq:coherent-exact-trace-expansion}
\end{equation}
For a partition with $r$ blocks, Lemma~\ref{lem:coherent-graph-integral}
and $n/D_m\to\rho$ give its limit
$\rho^{r-1}\tau(G_\pi)^{-q}$.

For the variance, expand two trace words.  Disjoint label sets have
zero covariance.  When the words share at least one label, their
combined graph is connected and Eulerian.  If it has $r$ vertices,
its integral is $O(D_m^{-(r-1)})$ and there are $O(n^r)$ labelings.
The two trace normalizations leave $O(n^{-1})$.  Products of the
separate means have the same or smaller order: if the separate words
use $r_1,r_2$ vertices and share $s\geq1$, then
$r=r_1+r_2-s$, giving $O(n^{-s})$.  There are only finitely many
patterns for fixed $p$.

The spectral measures are tight since their first moments equal one.
The bounds on higher moments imply uniform integrability of every
lower moment.  Thus subsequential limits of the expected spectral
measures exist and have all the displayed moments.  Since
$\tau\geq1$,
\[
 M_{2p}(q,\rho)
 \leq \max(1,\rho^{2p-1})(2p)^{2p}.
\]
Carleman's criterion gives uniqueness of the moment problem.
The variance estimate then gives weak and moment convergence in
probability.  Chebyshev's inequality and Borel--Cantelli give the
almost-sure assertion whenever $\sum_m n_m^{-1}<\infty$; no
independence between different rows is required.

For the final assertion let $\Omega_m$ be the prescribed moment
event.  Boundedness and independence give
$\Pr(\Omega_m^c)=O((n_m\epsilon_m^2)^{-1})=o(1)$.
For $T_{m,p}=n_m^{-1}\operatorname{tr}\bigl[(B^{(m)})^p\bigr]$,
\[
 \mathbb E\bigl[|T_{m,p}-\mathbb ET_{m,p}|^2\mid\Omega_m\bigr]
 \leq\frac{\operatorname{Var}(T_{m,p})}{\Pr(\Omega_m)}
 =O_p(n_m^{-1}).
\]
This proves the same concentration, moment convergence, and
summability assertions for the conditioned rows.
\end{proof}

\subsection{Unbounded support and the emergence of noncrossing partitions}

The fixed-dimensional laws are not simply Marchenko--Pastur laws with
a shifted parameter.  Their support is unbounded, even though the
large-dimension limit has compact support.

\begin{lemma}[The noncrossing layer]
\label{lem:tree-one-noncrossing}
For a cycle quotient, $\tau(G_\pi)=1$ if and only if $\pi$ is
noncrossing in the cyclic order.
\end{lemma}

\begin{proof}
A noncrossing partition with more than one block has a proper interval
block.  Its quotient vertex has one nonloop incoming edge and one
nonloop outgoing edge.  Deleting this vertex and joining these edges
is a series reduction with Laplacian diagonal entry one; a Schur
complement shows that a cofactor rooted elsewhere is unchanged.
Induction ends with the one-vertex graph.

Conversely, a crossing supplies blocks $a,b$ occurring in order
$a,\ldots,b,\ldots,a,\ldots,b$.  At least two intervals between
consecutive occurrences of $a$ contain $b$.  Follow each interval
to its first $b$ and erase internal cycles.  This gives two simple
directed paths from $a$ to $b$ with distinct initial edges; parallel
edges are distinguished.  Each extends to an in-arborescence rooted
at $b$, by attaching remaining vertices along paths to the current
tree.  Such paths exist because a connected Eulerian graph is
strongly connected.  The two arborescences choose different outgoing
edges at $a$, so $\tau\geq2$.
\end{proof}

\begin{theorem}[Support, exponential moments, and the free limit]
\label{thm:coherent-free-limit}
For every finite positive integer $q$ and $\rho>0$, the support of
$\nu_{q,\rho}$ is unbounded.  Nevertheless its moment-generating
function is entire, and for $t\geq0$,
\begin{equation}
 \int e^{tx}\,d\nu_{q,\rho}(x)
 \leq1+\frac{\exp\{\rho(e^t-1)\}-1}{\rho}.
 \label{eq:coherent-entire-bound}
\end{equation}
In particular, for $R\geq\rho$,
\begin{equation}
 \nu_{q,\rho}([R,\infty))
 \leq\max(1,\rho^{-1})
       \exp\{-R\log(R/\rho)+R-\rho\}.
 \label{eq:coherent-poisson-tail}
\end{equation}
As $q\to\infty$ through integers, $\nu_{q,\rho}$ converges weakly
and in every fixed positive moment to the mean-one
Marchenko--Pastur law of aspect ratio $\rho$.
\end{theorem}

\begin{proof}
Consider all partitions of $2r$ positions into pairs.  Each quotient
vertex has outdegree at most two, after ignoring loops.  An
in-arborescence chooses one outgoing edge at every nonroot vertex,
so $\tau\leq2^{r-1}$.  There are $(2r-1)!!$ pair partitions.
Equation~\eqref{eq:coherent-tree-moments} therefore gives
\begin{equation}
 M_{2r}(q,\rho)
 \geq (2r-1)!!\left(\frac{\rho}{2^q}\right)^{r-1}.
 \label{eq:coherent-pairing-lower-bound}
\end{equation}
The $2r$-th root of the right side diverges with $r$.  A probability
measure with bounded support could not satisfy this inequality.

For the upper bound let
$T_p(\rho)=\sum_{\pi\in\mathfrak P(p)}\rho^{|\pi|}$
be the Touchard polynomial.  For $p\geq1$,
$M_p(q,\rho)\leq T_p(\rho)/\rho$.
The exponential generating function
$\sum_{p\geq0}T_p(\rho)t^p/p!=\exp\{\rho(e^t-1)\}$,
together with monotone convergence, proves
\eqref{eq:coherent-entire-bound}.  Finiteness for every $t\geq0$
also gives an entire complex moment-generating function.  The right
side of \eqref{eq:coherent-entire-bound} is at most
$\max(1,\rho^{-1})\exp\{\rho(e^t-1)\}$.
Markov's inequality with $t=\log(R/\rho)$ proves
\eqref{eq:coherent-poisson-tail}.

Finally, Lemma~\ref{lem:tree-one-noncrossing} gives
\[
 M_p(q,\rho)=
 \sum_{\pi\in NC(p)}\rho^{|\pi|-1}
 +\sum_{\pi\notin NC(p)}\rho^{|\pi|-1}\tau(G_\pi)^{-q}.
\]
The second sum is at most
$2^{-q}\sum_{\pi\notin NC(p)}\rho^{|\pi|-1}$.
The first is the $p$-th mean-one Marchenko--Pastur moment:
Kreweras complementation interchanges the exponents $|\pi|-1$
and $p-|\pi|$ in the usual noncrossing formula
\cite{prob_nica_speicher2006}.  The uniform Touchard bounds yield
tightness and uniform integrability of fixed moments.  Moment
determinacy of the compactly supported limiting law completes the
proof.
\end{proof}

\subsection{Finite-degree identities and the logarithmic scale}

The first moments admit simple exact checks before taking a limit.
With $D=D_m$ and $D'=D_{2m}$, the identities
$\int P_z^{\otimes m}d\sigma=I/D$ and
$\int|\langle z,w\rangle|^{4m}d\sigma(w)=1/D'$ imply
\begin{align}
 \mathbb E T_{m,2}&=1+\frac{n-1}{D},\notag\\
 \mathbb E T_{m,3}&=1+\frac{3(n-1)}D
                     +\frac{(n-1)(n-2)}{D^2},\notag\\
 \mathbb E T_{m,4}&=1+\frac{6(n-1)}D+\frac{n-1}{D'}
   +\frac{6(n-1)(n-2)}{D^2}
   +\frac{(n-1)(n-2)(n-3)}{D^3}.
 \label{eq:coherent-finite-moments}
\end{align}
The supplemental verifier enumerates cycle partitions, computes all
rooted cofactors by integer elimination, and checks the noncrossing
criterion and the rational moment identities.  These checks concern
finite combinatorics; numerical spectra provide illustrations of the
limiting law.

For anticanonical degree $m=(q+1)k$ and $n=D_m$, let $\Delta_k$
be the Slater determinant of a Fubini--Study orthonormal basis.
The evaluation kernel is $D_m\langle z,w\rangle^m$, so
\begin{equation}
 E_k:=-\frac1{kD_m}\log\|\Delta_k\|^2
 =-\frac1k\int\log\lambda\,d\nu_m(\lambda)
    -\frac{\log D_m}{k}.
 \label{eq:coherent-logarithmic-energy}
\end{equation}
The spectral law consequently gives a precise object for investigating
the collective determinant energy.  Its natural continuation is a
logarithmic-tail theorem on scales $\lambda\asymp e^{-kt}$,
and its analogue under attractive Gibbs reweighting.  Such a theorem
would connect coherent spectral statistics with the moment-reduced
large-deviation program for K\"ahler--Einstein measures
\cite{AndreassonBermanSvensson2026}.  Equation~\eqref{eq:coherent-logarithmic-energy}
also specifies the distinct scales involved: fixed positive eigenvalue
scales contribute only $o(1)$ to this normalized energy.

\section{Torus averaging and optimal transport}
\label{tor:section}

This section establishes mean-energy convergence for the actual complex
Slater determinant under every torus-invariant reference satisfying a
logarithmic moment condition. No negative-temperature large-deviation
principle is assumed. The argument identifies a finite computational
bridge: the phase-averaged logarithm is within an explicit vanishing
error of an assignment problem with relative-entropy cost.

Put $D=q+1$, $X=\mathbb P^q$, $m=Dk$, and
$N=\binom{m+q}{q}$. Let
\[
 \mathcal A_m=\{\alpha\in\mathbb N^D:|\alpha|=m\},\qquad
 \lambda_m=N^{-1}\sum_{\alpha\in\mathcal A_m}\delta_{\alpha/m}.
\]
Write $\Delta_q=\{p\in\mathbb R_+^D:\sum p_j=1\}$ and let
$\lambda$ be normalized Lebesgue measure on this simplex. The coherent
monomial basis, orthonormal for Fubini--Study probability, is
\[
 s_\alpha(z)=\sqrt{N\binom m\alpha}\,z^\alpha,
 \qquad \binom m\alpha=\frac{m!}{\alpha_0!\cdots\alpha_q!},
 \qquad \|z\|=1.
\]
Let $\Delta_k(\boldsymbol z)=\det(s_\alpha(z_i))_{i,\alpha}$ and
$E_k=-(kN)^{-1}\log\|\Delta_k\|^2$. All phases in the statements
below are independent Haar phases, one projective torus per particle.
For $p_i\in\Delta_q^\circ$, write
$z_i=(\sqrt{p_{i0}},\sqrt{p_{i1}}e^{i\theta_{i1}},\ldots,
\sqrt{p_{iq}}e^{i\theta_{iq}})$.

\subsection{An exposed coefficient controls the phase average}

\begin{lemma}[The assignment vertices]
\label{tor:vertices}
Every monomial in $\det(z_i^\alpha)_{i,\alpha\in\mathcal A_m}$
is an exposed vertex of its Newton polytope. Consequently, if
\[
 A_m(\boldsymbol p)=
 \max_{\pi:\{1,\ldots,N\}\simeq\mathcal A_m}
 \prod_{i=1}^N N\binom m{\pi(i)}p_i^{\pi(i)},
\]
then
\begin{equation}
 \log A_m(\boldsymbol p)
 \le \int\log\|\Delta_k(\boldsymbol z)\|^2\,d\boldsymbol\theta
 \le \log A_m(\boldsymbol p)+\log N!.
 \label{tor:phase-squeeze}
\end{equation}
The logarithm is integrable for every such $\boldsymbol p$.
\end{lemma}

\begin{proof}
Delete the zeroth component from each exponent and denote the resulting
integer vector by $\alpha'\in\mathbb N^q$. For a fixed assignment
$\pi$, put $v_i=\pi(i)'$. Its scalar product with an assignment $\tau$
is $\sum_i\pi(i)'\cdot\tau(i)'$. Since both lists enumerate the same
set of distinct vectors,
\[
 \sum_i\|\pi(i)'\|^2-
 \sum_i\pi(i)'\cdot\tau(i)'
 =\frac12\sum_i\|\pi(i)'-\tau(i)'\|^2.
\]
The difference is positive unless $\tau=\pi$, proving exposure.
No two determinant terms have the same exponent vector, so the
coefficient of every term is nonzero.

For completeness, the elementary vertex form of Jensen's inequality
suffices here. For fixed $\boldsymbol\theta$, substitute
$\theta_i\mapsto\theta_i+t v_i$ into the determinant. This is a
one-variable polynomial in $e^{it}$ whose highest-degree term is the
chosen assignment, uniquely. Jensen's formula bounds its logarithmic
circle mean below by the logarithm of the modulus of that leading
coefficient. Integrating $\boldsymbol\theta$ and using invariance of
Haar measure gives the lower bound for the chosen assignment, hence for
the maximum. Logarithms of nonzero trigonometric polynomials are
integrable; one can also obtain this fact by successive one-variable
factorization and Jensen. The upper bound follows from Jensen's
inequality and orthogonality of the distinct torus characters:
\[
 \int\log|\Delta_k|^2\le
 \log\int|\Delta_k|^2
 =\log\sum_\pi\prod_iN\binom m{\pi(i)}p_i^{\pi(i)}
 \le\log A_m+\log N!.
\]
\end{proof}

For $a\in\Delta_q$, $p\in\Delta_q^\circ$, use the conventions
$0\log0=0$ and
\[
 c(a,p)=D_{\rm KL}(a\Vert p)=\sum_{j=0}^q a_j\log(a_j/p_j).
\]
Define the finite assignment cost
\[
 C_m(\boldsymbol p)=\min_\pi\frac1N\sum_{i=1}^N
                         c(\pi(i)/m,p_i).
\]
This is also the Kantorovich transport cost from $\lambda_m$ to
$N^{-1}\sum_i\delta_{p_i}$: both empirical measures have $N$ equally
weighted atoms, and the doubly stochastic polytope has permutation
vertices.

\begin{proposition}[Finite assignment certificate]
\label{tor:finite-certificate}
For every $\boldsymbol p\in(\Delta_q^\circ)^N$,
\begin{equation}
 D C_m(\boldsymbol p)-\frac{\log N}{k}
                       -\frac{\log N!}{kN}
 \le\int E_k(\boldsymbol z)\,d\boldsymbol\theta
 \le D C_m(\boldsymbol p).
 \label{tor:finite-certificate-bound}
\end{equation}
In particular the error is at most $2\log N/k$.
\end{proposition}

\begin{proof}
For $a=\alpha/m$, the multinomial type probability at its own parameter
$a$ is a largest type probability. Indeed, for another type $\beta$,
\[
 \frac{\binom m\beta a^\beta}{\binom m\alpha a^\alpha}
 =\prod_j\frac{\alpha_j!}{\beta_j!}
              \alpha_j^{\beta_j-\alpha_j}\le1.
\]
If $\alpha_j=0<\beta_j$ the numerator probability is zero; otherwise
each displayed factor is at most one, by writing out the consecutive
integer factors. The powers of $m$ cancel. There are $N$ types, so
\[
 N^{-1}\le\binom m\alpha a^\alpha\le1,
 \qquad
 e^{m h(a)}\le N\binom m\alpha\le N e^{m h(a)},
 \quad h(a)=-\sum_j a_j\log a_j.
\]
Therefore
\[
 -mN C_m\le\log A_m\le -mN C_m+N\log N.
\]
Apply Lemma~\ref{tor:vertices} and divide by $-kN$, noting $m/k=D$.
\end{proof}

\fig{fig13_toric_transport}{A finite transport certificate for Slater
energy. (a) In complex dimension $q=2$ at degree $m=6$, the 28 exponent-grid
points $a=\alpha/6$ (open blue circles) are matched to the rational
squared-modulus points $p=(\alpha+\mathbf 1)/9$ (orange). The arrows show
the identity assignment, which is optimal for this deterministic example
by the monotonicity of $\log(\alpha_j+1)$. (b) The matrix of exact costs
$D_{\rm KL}(a\Vert p)$, displayed numerically, with the optimal diagonal
shown in white. The matrix order is the same in the two coordinates.
(c) The guaranteed interval width $\log N_k/k+\log(N_k!)/(kN_k)$ in
\eqref{tor:finite-certificate-bound}, for $q=1,2,3$ and
$N_k=\binom{(q+1)k+q}{q}$. This interval controls the independent
phase average of the actual logarithmic determinant energy for fixed
moduli. It is not a pointwise bound in the phases or a bound on an
attractive exponential moment.}{fig:toric-transport}

\subsection{The transport limit, including the simplex boundary}

\begin{lemma}[Convergence of the unbounded assignment cost]
\label{tor:cost-limit}
Let $\eta$ be a probability on $\Delta_q^\circ$ such that
\begin{equation}
 \int b(p)\,d\eta(p)<\infty,
 \qquad b(p)=\sum_{j=0}^q-\log p_j.
 \label{tor:log-moment}
\end{equation}
For iid $p_1,p_2,\ldots$ with law $\eta$,
\[
 C_m(p_1,\ldots,p_N)\longrightarrow
 C(\eta):=\inf_{\kappa\in\Pi(\lambda,\eta)}
                   \int c(a,p)\,d\kappa(a,p)
\]
almost surely and in $L^1$ as $m\to\infty$ (along $m=Dk$ or all
integer degrees).
\end{lemma}

\begin{proof}
The uniform lattice measures $\lambda_m$ converge weakly to $\lambda$
by Riemann sums on the simplex. The empirical measures $\eta_N$
converge weakly almost surely. For $R>0$, let
\[
 c_R(a,p)=\sum_j a_j\log a_j+
                  \sum_j a_j\min\{-\log p_j,R\}.
\]
Extend the truncated logarithm by $R$ at $p_j=0$. Then $c_R$ is
bounded and continuous on the compact product $\Delta_q^2$.
Its transport cost is continuous in both marginals. One direct proof
couples nearby marginals by couplings converging to the diagonal,
composes these with a fixed optimal coupling, and uses uniform
continuity; compactness of optimal couplings gives the reverse bound.
Thus the $c_R$ costs for $(\lambda_m,\eta_N)$ converge to the $c_R$
cost for $(\lambda,\eta)$.

The truncation error is nonnegative and independent of the coupling:
\[
 0\le c(a,p)-c_R(a,p)
 \le t_R(p):=\sum_j(-\log p_j-R)_+.
\]
Consequently the difference of the two optimal costs is between zero
and $\int t_R\,d\eta_N$, and the analogous assertion holds for
$\eta$. For every integer $R$, the strong law makes
$\int t_R\,d\eta_N\to\int t_R\,d\eta$ simultaneously almost surely.
The latter tends to zero as $R\to\infty$ by
\eqref{tor:log-moment}. This proves almost-sure convergence.
Moreover the expected truncation error equals exactly
$\int t_R\,d\eta$, for every $N$. Bounded convergence for $c_R$ costs,
followed by $R\to\infty$, proves $L^1$ convergence. This argument does
not assume a uniform positive lower bound on the sample coordinates.
\end{proof}

\subsection{Identification with anticanonical pluricomplex energy}

\begin{proposition}[Relative-entropy transport formula]
\label{tor:energy-transport}
Let $\mu$ be invariant under the standard projective torus, and let
$\eta$ be its squared-modulus pushforward. Under
\eqref{tor:log-moment}, its anticanonical pluricomplex energy is
\begin{equation}
 E(\mu)=D C(\eta).
 \label{tor:energy-transport-formula}
\end{equation}
\end{proposition}

\begin{proof}
Averaging any admissible potential over the compact torus increases its
concave Monge--Amp\`ere energy and leaves its integral against $\mu$
unchanged. The energy dual may therefore be restricted to invariant
potentials. In logarithmic coordinates
$y_j=\log(p_j/p_0)$ put
\[
 \phi_0(y)=D\log(1+e^{y_1}+\cdots+e^{y_q}),\qquad
 \phi_0^*(Da')=D\sum_{j=0}^q a_j\log a_j,
\]
where $a'=(a_1,\ldots,a_q)$ and the moment polytope is $D\Delta_q$.
For an invariant positive metric $\phi=\phi_0+u$, normalized
Monge--Amp\`ere energy is
\[
 \mathcal E_L(u)=-\int_{\Delta_q}
                   [\phi^*(Da')-\phi_0^*(Da')]\,d\lambda(a).
\]
This follows first for smooth metrics by differentiating the Legendre
transform and changing variables by the moment map, then for finite
energy by monotone approximation. The energy dual consequently equals
\[
 \int\phi_0(y)\,d\eta_y(y)+\int\phi_0^*(Da')\,d\lambda(a)
 -\inf_\phi\left\{\int\phi(y)\,d\eta_y(y)
                   +\int\phi^*(Da')\,d\lambda(a)\right\}.
\]
The infimum is the maximal coupling integral
$\sup_{\pi\in\Pi(\lambda,\eta_y)}\int D a'\cdot y\,d\pi$.
Here is an explicit approximation argument for this duality.  First
assume that $\eta_y$ is supported on a compact cube.  Compact
Kantorovich duality for the continuous payoff $D a'\cdot y$ supplies
a bounded continuous dual function $h$ on $\Delta_q$.  Its transform
\[
 \phi(y)=\sup_{a\in\Delta_q}\{D a'\cdot y-h(a)\}
\]
has slopes in $D\Delta_q$, and $\phi^*(Da')\le h(a)$.
Fenchel's inequality gives the reverse comparison with every
coupling, so replacing the dual pair by $(\phi,\phi^*)$ attains
the same infimum.  This transform differs by a bounded function from
$D\max(0,y_1,\ldots,y_q)$, hence also from $\phi_0$.
Continuity of $h$ shows that $u=\phi-\phi_0$ extends continuously
across each coordinate face: in the limit only exponents supported
on that face contribute to the supremum.  It therefore defines a
continuous toric positive metric on $\mathbb P^q$.  Unitary heat
averaging, followed by a vanishing convex mixture with the reference
metric, approximates $u$ uniformly by smooth strictly positive toric
metrics.  Uniform convergence gives convergence both of its integrals
and of the Monge--Amp\`ere energy, so these transforms are legitimate
in the energy dual.

For general $\eta_y$, clip $y$ coordinatewise to $[-R,R]^q$.
The logarithmic moment assumption implies $\int|y|\,d\eta_y<\infty$,
so the mean clipping displacement tends to zero.  After normalizing
$\phi(0)=0$, every convex potential with slopes in $D\Delta_q$ is
$D$-Lipschitz.  Its integral changes under clipping by at most
$D\int|y-y^{(R)}|\,d\eta_y$, uniformly over all these potentials.
The same bound holds for every coupling payoff $D a'\cdot y$;
couplings can be pushed forward under clipping or lifted by
disintegration.  Thus both sides of the compact duality converge,
proving it for $\eta_y$ without any boundary-support assumption
beyond \eqref{tor:log-moment}.  The resulting cost is
\[
 \phi_0(y)+\phi_0^*(Da')-Da'\cdot y
 =D\sum_j a_j\log(a_j/p_j).
\]
This proves \eqref{tor:energy-transport-formula}. The Legendre--moment-map
correspondence is the toric optimal-transport description of
\cite{BermanWittNystrom2014}; the energy dual is part of the
variational framework of \cite{BBGZ2013}.
\end{proof}

\begin{theorem}[Actual Slater mean energy for toric reference laws]
\label{tor:mean-energy}
Let $\mu$ satisfy the hypotheses of Proposition~\ref{tor:energy-transport}.
Then, for iid particles with law $\mu$,
\begin{equation}
 \lim_{k\to\infty}\mathbb E_{\mu^{\otimes N_k}}E_k=E(\mu).
 \label{tor:mean-energy-limit}
\end{equation}
Whenever the deterministic continuous-envelope lower comparison is
available for this polarization, $E_k\to E(\mu)$ also in $L^1$ under
these product laws. In particular, for Haar probability on $\mathbb P^q$,
\[
 \mathbb E_{\sigma^{\otimes N_k}}|E_k|\longrightarrow0.
\]
On $\mathbb P^2$ this remains true after conditioning on any fixed-order
projective moment tube of width $N^{-a}$, $0<a<1/2$.
\end{theorem}

\begin{proof}
Torus invariance disintegrates $\mu$ into $\eta$ and independent uniform
angular variables. Integrate \eqref{tor:finite-certificate-bound} and
apply Lemma~\ref{tor:cost-limit} and
Proposition~\ref{tor:energy-transport}. The error tends to zero since
$\log N/k\to0$. All logarithmic determinant integrals are finite, by
the upper bound and the universal lower bound $E_k\ge-\log N/k$.

For the $L^1$ assertion, empirical measures converge to $\mu$, and the
continuous-envelope lower comparison gives
$\Pr(E_k<E(\mu)-\epsilon)\to0$ for every $\epsilon>0$.
The negative deviation $(E(\mu)-E_k)_+$ is bounded by
$E(\mu)+\log N/k$, hence tends to zero in mean. The already proved
convergence of the means then forces the positive deviation to vanish
in mean too.

For Haar probability, $\eta=\lambda$ and the diagonal coupling has
cost zero. The universal lower bound alone therefore proves $L^1$
convergence in every dimension. If $\Omega_N$ is the stated tube,
$\sigma^{\otimes N}(\Omega_N)\to1$ by the finite-moment variance
estimate. The conditional expectation of the nonnegative variable
$E_k+\log N/k$ is at most its unconditioned expectation divided by
$\sigma^{\otimes N}(\Omega_N)$, and hence also tends to zero.
\end{proof}

\paragraph{From mean energy to attractive pressure.}
The finite certificate gives a product-law mean-energy statement for the
actual Slater Hamiltonian using an ordinary assignment solver, known
constants, and explicit logarithmic boundary control. The full attractive
pressure in Section~\ref{att:entropy-section} requires the additional
entropy and conditional-density argument developed there. These are
different levels of control: convergence of the mean alone does not
permit replacing the determinant by its phase-averaged logarithm inside
an attractive partition function.

\section{Moment constraints and repulsive determinant ensembles}
\label{sec:repulsive-constraints}

Positive moment repair also acts directly on particle configurations.
It can enforce a shrinking design constraint by resampling a vanishing
fraction of the particles while controlling both determinant energy
and relative entropy.  This gives a constrained version of the
positive-temperature determinant large-deviation theorem
\cite{prob_berman2017}.

Fix a design order $t\geq1$ on $X=\mathbb P^2$, and set
\[
 T_t(z)=P_z^{\otimes t},\qquad
 H_t=\frac{I_{\Sym^t\C^3}}{\binom{t+2}{2}},\qquad
 \mathcal D_t=\left\{\mu:\int T_t\,d\mu=H_t\right\}.
\]
For $N=N_k=\binom{3k+2}{2}$ let $E_k$ be the anticanonical Slater
energy in \eqref{eq:coherent-logarithmic-energy}, and let $E$ be the
normalized pluricomplex measure energy for $(X,\mathcal O(3))$,
with $E\geq0$ and $E(\sigma)=0$.  Fix $0<a<1/2$
and define
\[
 \Omega_N=\left\{\left\|\frac1N\sum_iT_t(z_i)-H_t\right\|_{\HS}
                      <N^{-a}\right\},
 \qquad Q_N=\sigma^{\otimes N}(\cdot\mid\Omega_N).
\]
The design order $t$ is fixed before $k\to\infty$.

\begin{lemma}[Conditional replacement estimate]
\label{lem:conditional-replacement}
If $\Delta_k(\boldsymbol x)\ne0$ and a set of $m$ particle positions
is replaced independently by laws having densities at most $3/2$
against $\sigma$, then the resulting tuple $\boldsymbol Y$ satisfies
\[
 \Pr\!\left(E_k(\boldsymbol Y)>E_k(\boldsymbol x)
                   +c_*m/N\right)\leq e^{-m},
 \qquad c_*=6(1+\log4).
\]
\end{lemma}

\begin{proof}
For a nonzero section $p\in H^0(\mathbb P^2,\mathcal O(3k))$,
a point $z_0$ with $p(z_0)\ne0$, and $0<s<1/3$, one has
\[
 \int\left(\frac{\|p(z)\|}{\|p(z_0)\|}\right)^{-2s/k}d\sigma(z)
 \leq C(3s),\qquad C(u)=\frac2{(1-u)(2-u)}.
\]
To verify this, it suffices to replace $z_0$ by a maximum point.
Disintegrate Fubini--Study measure along complex lines through it.
On such a line the radial density is $2(1-x)dx\,d\theta/(2\pi)$.
Factor the restriction into $3k$ normalized linear forms and apply
H\"older.  The inverse integral of one unit linear form, whose
squared value at the distinguished point is $b$, is
$2(1-u+ub)/((1-u)(2-u))\leq C(u)$.  The normalizing values of the
individual linear forms contribute factors at most one.

Apply this estimate successively to the determinant in each replaced
coordinate, evaluating it at the old coordinate after each integration.
The restrictions that vanish identically occur on analytic null sets.
At $s=1/6$, $C(3s)=8/3$; the density bounds therefore give
\[
 \mathbb E\exp\left(\frac N6
       [E_k(\boldsymbol Y)-E_k(\boldsymbol x)]\right)\leq4^m.
\]
Exponential Markov inequality proves the claim.
\end{proof}

\begin{proposition}[Entropy-controlled repair of microscopic moments]
\label{prop:microscopic-moment-repair}
Suppose $\delta_N\downarrow0$.  There are integers $m_N=o(N)$,
with $m_N\to\infty$, and measurable replacement kernels
$\mathcal R_N$ on nonzero-determinant configurations of moment defect
at most $\delta_N$, such that, almost surely under the kernel,
\[
 \boldsymbol Y\in\Omega_N,\qquad
 \#\{i:Y_i\ne x_i\}\leq m_N,\qquad
 E_k(\boldsymbol Y)\leq E_k(\boldsymbol x)+c_*m_N/N.
\]
For every input law $R_N$ of finite relative entropy, its repaired law
$S_N=R_N\mathcal R_N$ satisfies
\begin{equation}
 \Ent(S_N\mid\sigma^{\otimes N})
 \leq\Ent(R_N\mid\sigma^{\otimes N})
           +m_N\log(3/2)+2\log2.
 \label{eq:microscopic-repair-entropy}
\end{equation}
The empirical measures before and after repair have bounded-Lipschitz
distance at most $2m_N/N$.
\end{proposition}

\begin{proof}
The coherent covariance map on traceless Hermitian matrices,
\[
 \mathsf S_t(A)=\int\tr(AT_t(z))T_t(z)\,d\sigma(z),
\]
is strictly positive: a coherent symbol vanishing identically is the
zero bihomogeneous polynomial.  Hence it is invertible.  There is a
finite constant $C_t\geq1$ such that, for every sufficiently small
traceless Hermitian $D$, the density
\[
 1+g_D(z),\qquad
 g_D(z)=\tr(\mathsf S_t^{-1}(D)T_t(z)),
\]
has mean one, moment $H_t+D$, and
$\|g_D\|_\infty\leq C_t\|D\|_{\HS}$.  It lies in $[1/2,3/2]$
when $\|D\|\leq1/(2C_t)$.

Choose $\theta_N=m_N/N\to0$ with $m_N\to\infty$ and
$\theta_N\geq4C_t\delta_N$.  A uniformly sampled $m_N$-subset
$I$ has average moment within $1/(4C_t)$ of the full empirical
moment $M_x$ with probability at least $1/2$, for large $N$:
sampling without replacement gives mean squared deviation at most
$1/m_N$.  Condition on these good subsets.  The desired replacement
moment is
\[
 M_\tau=\frac1{m_N}\sum_{i\in I}T_t(x_i)
                      +\frac{H_t-M_x}{\theta_N}.
\]
It is within $1/(2C_t)$ of $H_t$, so a density in $[1/2,3/2]$
realizes it.  Independent replacements from that density give total
expected moment $H_t$ and variance at most $m_N/N^2$.
The probability of violating the tube is at most $N^{2a-1}$.
Lemma~\ref{lem:conditional-replacement} bounds the probability of
violating the energy inequality by $e^{-m_N}$.  Condition on both
requirements; their joint probability is at least $1/2$ for large $N$.

For completeness, the entropy bound includes the adaptive choice of
labels.  Relative to the independent uniform law on subsets, selecting
a good subset costs at most $\log2$ in relative entropy.  Discarding
the old selected coordinates does not increase relative entropy.
Conditional on the retained coordinates and the subset, the new
coordinates have a mixture of densities bounded by
$2(3/2)^{m_N}$.  The entropy chain rule adds at most
$m_N\log(3/2)+\log2$.  Discarding the subset label gives
\eqref{eq:microscopic-repair-entropy}.  All selections and conditional
densities are Borel; changing $m_N$ atoms gives the empirical-distance
bound.
\end{proof}

\begin{theorem}[Shrinking-design pressure and large deviations]
\label{thm:repulsive-design-ldp}
For every $\beta\geq0$ and every continuous
$\Phi:\mathcal P(X)\to\mathbb R$,
\begin{equation}
 \lim_{k\to\infty}\frac1N\log
       \int e^{-N[\beta E_k+\Phi(L_N)]}\,dQ_N
 =-\inf_{\mu\in\mathcal D_t}
       [\Ent(\mu\mid\sigma)+\beta E(\mu)+\Phi(\mu)].
 \label{eq:repulsive-design-laplace}
\end{equation}
The normalized empirical laws satisfy a full speed-$N$ LDP with good
rate function
\[
 I_\beta(\mu)=\Ent(\mu\mid\sigma)+\beta E(\mu)
                         +\iota_{\mathcal D_t}(\mu).
\]
Its unique minimizer is $\sigma$, with minimum zero.  Here
$\iota_{\mathcal D_t}$ is zero on the constraint and infinity outside it.
\end{theorem}

\begin{proof}
First, $\sigma^{\otimes N}(\Omega_N)\to1$ by the variance estimate.
The reference conditioned laws have the constrained Sanov rate
$\Ent(\mu\mid\sigma)+\iota_{\mathcal D_t}$.  The upper bound is
Sanov's upper bound and the closed limiting constraint.  For the lower
bound at a finite-entropy $\mu\in\mathcal D_t$, use $\mu^{\otimes N}$:
its empirical law converges to $\mu$, its moment lies in the tube with
probability tending to one because $N^{1-2a}\to\infty$, and the
log-density law of large numbers gives the entropy cost.

Choose continuous nonnegative envelopes $E_r\uparrow E$.
Theorem~\ref{an:lower-comparison} gives
$E_k\geq E_r(L_N)-o_r(1)$ uniformly for each fixed $r$. The
reference Laplace principle therefore bounds the upper limit of
\eqref{eq:repulsive-design-laplace} by
$-\inf_{\mathcal D_t}(\Ent+\beta E_r+\Phi)$.  Let $r\to\infty$;
monotonicity and compactness of probability space give the desired
upper bound.

For the lower bound at $\beta>0$, use the unrestricted determinant
LDP of \citet[Theorems 1.1 and 5.7]{prob_berman2017}.  At every
$\mu\in\mathcal D_t$ of finite entropy and energy it supplies
shrinking empirical neighborhoods and restricted Gibbs laws $R_N$
with
\[
 \limsup_N\left[
   \frac1N\Ent(R_N\mid\sigma^{\otimes N})
                  +\beta\mathbb E_{R_N}E_k\right]
 \leq \Ent(\mu\mid\sigma)+\beta E(\mu).
\]
Indeed the bracket equals minus the normalized logarithm of the
restricted unnormalized partition function; the LDP lower bound and
a diagonal choice of shrinking neighborhoods give the assertion.
At finite $N$ entropy and energy are finite, since the positive
temperature weight suppresses logarithmic determinant singularities.

The moment defects on these neighborhoods tend to zero uniformly.
Proposition~\ref{prop:microscopic-moment-repair} gives laws $S_N$
supported on $\Omega_N$, changes their empirical measures by $o(1)$,
and increases the displayed mean free energy by at most
\[
 \frac{m_N}{N}[\log(3/2)+\beta c_*]+\frac{2\log2}{N}=o(1).
\]
Furthermore
$\Ent(S_N\mid Q_N)=\Ent(S_N\mid\sigma^{\otimes N})
+\log\sigma^{\otimes N}(\Omega_N)$.
The Gibbs variational principle now proves the pressure lower bound,
including the continuous tilt $\Phi$.  At $\beta=0$ the assertion
is the constrained Sanov theorem.  The Laplace principle for every
continuous $\Phi$ on compact probability space is equivalent to the
full LDP.  Finally $E\geq0$, $E(\sigma)=0$, and entropy has unique
zero at $\sigma$, proving the normalization and uniqueness claims.
\end{proof}

\begin{corollary}
For every fixed $\beta>0$, the unperturbed repulsive ensemble has
$\mathbb E E_k\to0$ and relative entropy per particle tending to
zero, relative to either $Q_N$ or $\sigma^{\otimes N}$.
\end{corollary}
\begin{proof}
The limiting unperturbed pressure is zero at every positive $\beta$.
Convex difference quotients at $\beta\pm h$ imply convergence of
the finite-level derivative $-\mathbb E E_k$ to zero.
The Gibbs variational identity then gives the entropy limit.
\end{proof}

\section{Attractive determinant large deviations}
\label{att:entropy-section}

The argument here concerns the actual Slater determinant on $X=\mathbb
P^2$, with anticanonical degree $3k$, $N=\binom{3k+2}{2}$, and
$E_k=-(kN)^{-1}\log\|\Delta_k\|^2$. Write $\sigma$ for
Fubini--Study probability. The unconditioned attractive ensemble is
\[
 Z_k(\gamma)=\int e^{\gamma N E_k}\,d\sigma^{\otimes N},
 \qquad P_{k,\gamma}=Z_k(\gamma)^{-1}e^{\gamma N E_k}\sigma^{\otimes N}.
\]
The scalar inverse-section estimate ensures finiteness for $0\le\gamma<1/3$.
The proof first obtains the product-law mean energy from the established
repulsive pressure, using a uniform two-sided estimate at zero coupling.
Compact full conditional densities then give the entropy and energy
limits for the actual attractive tilted Gibbs laws. No generic
microscopic compactness hypothesis is imposed on arbitrary symmetric
laws, and no attractive pressure statement is used as an input.

\subsection{Entropy is determined by coarse types when conditionals are compact}

Let $(X,\sigma)$ be a compact probability space with an increasing
sequence $\mathcal A_j$ of finite Borel partitions generating its Borel
sigma field modulo null sets. Cells of positive probability suffice.
Let $\mathbb E_j$ denote conditional expectation onto $\mathcal A_j$.

\begin{lemma}[An exact coarse-graining bound]
\label{att:coarse-entropy}
Let $P_N\ll\sigma^{\otimes N}$ have finite relative entropy, and let
$f_{i,x_{-i}}$ be its full one-coordinate conditional density relative
to $\sigma$. Denote the law of its cell labels by $P_N^{(j)}$ and the
reference cell probabilities by $\sigma_j$. Then
\begin{align}
 0&\le D(P_N\Vert\sigma^{\otimes N})-
               D(P_N^{(j)}\Vert\sigma_j^{\otimes N})\\
 &\le\sum_{i=1}^N\mathbb E_{P_{N,-i}}
    \left[\operatorname{Ent}_\sigma(f_{i,x_{-i}})
              -\operatorname{Ent}_\sigma(\mathbb E_j f_{i,x_{-i}})\right].
 \label{att:coarse-bound}
\end{align}
Here $\operatorname{Ent}_\sigma(f)=\int f\log f\,d\sigma$ for a
probability density $f$.
\end{lemma}

\begin{proof}
Condition on the entire cell-label vector $Y$. Relative entropy chain
rule identifies the difference on the left with
\[
 \mathbb E_Y D\left(P_N(\cdot\mid Y)\,
       \middle\Vert\,\bigotimes_i\sigma(\cdot\mid Y_i)\right).
\]
Tensorize relative entropy inside this product of cells, using full
one-coordinate conditionals. Given $x_{-i}$, the conditional law of
$x_i$ within cell $A$ has density $f_i/\int_Af_i\,d\sigma$ against
$\sigma|_A$. Averaging the relative entropy to normalized $\sigma|_A$
over its cell probabilities gives exactly
\[
 \int f_i\log\frac{f_i}{\mathbb E_j f_i}\,d\sigma
 =\operatorname{Ent}_\sigma(f_i)-
                       \operatorname{Ent}_\sigma(\mathbb E_jf_i).
\]
This proves the upper bound. The lower bound is data processing.
\end{proof}

\begin{lemma}[Entropy chaos from compact conditional densities]
\label{att:entropy-chaos}
Suppose $P_N$ is exchangeable and every full one-coordinate conditional
density belongs to a fixed relatively compact subset $\mathcal F$ of
$L^r(\sigma)$ for some $r>1$. Suppose the empirical laws converge to
$\Gamma\in\mathcal P(\mathcal P(X))$, and partition boundaries have
zero mass for the limiting first marginal. Then
\[
 \lim_N\frac1N D(P_N\Vert\sigma^{\otimes N})
 =\int D(\mu\Vert\sigma)\,d\Gamma(\mu).
\]
In particular, if empirical measures converge in probability to
$\sigma$, the relative entropy per particle tends to zero.
\end{lemma}

\begin{proof}
The operators $\mathbb E_j$ converge strongly to the identity in
$L^r$ and are contractions. Convergence is uniform on the compact
closure of $\mathcal F$. The entropy functional is continuous on
$L^r$-convergent families of densities; uniform $L^r$ control supplies
uniform integrability of $f\log f$. Consequently
\[
 \epsilon_j:=\sup_{f\in\mathcal F}
 [\operatorname{Ent}_\sigma(f)-
       \operatorname{Ent}_\sigma(\mathbb E_jf)]\longrightarrow0.
\]
By Lemma~\ref{att:coarse-entropy}, the fine and coarse entropies per
particle differ by at most $\epsilon_j$.

For a fixed partition with $M$ cells, exchangeability makes the label
law uniform on every type class. If $\nu_N$ is the vector of empirical
cell frequencies, the method-of-types bounds give
\[
 \left|N^{-1}D(P_N^{(j)}\Vert\sigma_j^{\otimes N})
               -\mathbb E D(\nu_N\Vert\sigma_j)\right|
 \le \frac{2M\log(N+1)}N.
\]
One can obtain this directly from the entropy of the type law, which
is at most $M\log(N+1)$, and the corresponding upper and lower bounds
on multinomial type probabilities. The limiting frequencies have law
$(\mu\mapsto\mu_j)_*\Gamma$. Relative entropy on this finite simplex
is continuous and bounded, so the last expectation converges to
$\int D(\mu_j\Vert\sigma_j)\,d\Gamma$.
Let $N\to\infty$ and then $j\to\infty$. The finite-partition
entropies increase to full relative entropy. Tensorization also
bounds the original entropy per particle by
$\sup_{f\in\mathcal F}\operatorname{Ent}_\sigma(f)<\infty$,
so the limiting integral is finite. Monotone convergence completes
the proof.
\end{proof}

\subsection{The conditional compactness supplied by complex geometry}

\begin{lemma}[Compact attractive conditionals]
\label{att:conditional-compactness}
Fix $0\le\gamma<1/3$. The full one-coordinate conditional densities
of $P_{k,\gamma}$, for all $k$, lie in a fixed compact subset of
$L^r(\sigma)$ for some $r>1$.
\end{lemma}

\begin{proof}
For almost every retained tuple the determinant in the remaining
coordinate is a nonzero section $s$ of $\mathcal O(3k)$. Put
\[
 u=k^{-1}\log\|s\|^2-\sup_X k^{-1}\log\|s\|^2.
\]
Then $u$ is $\omega_0$-plurisubharmonic and $\sup u=0$. Its
conditional density is $e^{-\gamma u}/\int e^{-\gamma u}\,d\sigma$.
The normalized $\omega_0$-psh family is compact in $L^1$.
For every $0\le a<1/3$ it satisfies the uniform exponential estimate
\[
 \int e^{-a u}\,d\sigma\le C_a,
 \qquad C_a=\frac2{(1-3a)(2-3a)}.
\]
Choose $1<r<r'$ with $r'\gamma<1/3$. Every sequence of normalized
potentials has an $L^1$-convergent subsequence and, on a further
subsequence, pointwise almost-everywhere convergence. The exponential
bound at $r'\gamma$ and Vitali's theorem imply convergence of
$e^{-\gamma u}$ in $L^r$. Its integral is at least one, and
normalization is continuous in $L^r$. Thus the normalized densities
have compact closure in $L^r$. This argument includes limits with
analytic poles. At $\gamma=0$ the family consists of the constant one.
\end{proof}

The same conclusion holds with a fixed bounded continuous one-body
field, or with a smooth cylindrical empirical tilt whose one-coordinate
increments form a fixed compact bounded family. Such an extension
asserts compactness of conditionals; it does not assert empirical
concentration or identify a pressure for an arbitrary tilt.

\subsection{A uniform quadratic bound near zero attraction}

\begin{lemma}[A small-coupling exponential estimate]
\label{att:quadratic-pressure}
Let $\nu=g\sigma$ be a probability measure with $g\in L^p(\sigma)$
for some $p>1$. There are $t_0>0$ and $C_\nu<\infty$, independent
of $k$, such that
\begin{equation}
 t\,\mathbb E_{\nu^{\otimes N}}E_k
 \le \frac1N\log\int e^{tN E_k}\,d\nu^{\otimes N}
 \le t\,\mathbb E_{\nu^{\otimes N}}E_k+C_\nu t^2,
 \qquad |t|\le t_0.
 \label{att:quadratic-bound}
\end{equation}
The constants can be chosen uniformly when the exponent $p>1$ and the $L^p$ bound of $g$ are fixed. The assertion includes finiteness of the integrals and mean energies.
\end{lemma}

\begin{proof}
Let $p'=p/(p-1)$, and choose $s>0$ with $p's<1/3$. For every
normalized $\omega_0$-psh $u\le0$, H\"older gives
\[
 \int e^{-su}\,d\nu\le
 M:=\|g\|_p C_{p's}^{1/p'}<\infty.
\]
The section inverse estimate can be anchored at any nonzero value,
since its normalized maximum dominates that value. Applying it
successively in all particle coordinates shows that the joint
partition is finite for $t\le s$, at every level. Its logarithmic
moments are finite on compact subintervals below $s$.

For a full conditional at parameter $t$, set $Y=-u\ge0$ and
$F(t)=\log\int e^{tY}\,d\nu$. Its relative entropy to $\nu$ is
$tF'(t)-F(t)$. Jensen and the exponential bound give
$\mathbb E_\nu Y\le s^{-1}\log M$. Hence, for $|t|\le s/2$,
\[
 \int e^{tY}\,d\nu\ge e^{t\mathbb E_\nu Y}\ge M^{-1/2},
 \qquad
 F''(t)\le M^{1/2}\int Y^2e^{sY/2}\,d\nu
                         \le\frac{16M^{3/2}}{s^2}.
\]
In particular
\[
 0\le tF'(t)-F(t)=\int_0^t rF''(r)\,dr
                         \le C_\nu t^2,
 \qquad C_\nu=8M^{3/2}/s^2.
\]
For negative $t$, the integral has negative integrand and reversed
orientation, so it is still nonnegative. Tensorization of relative
entropy for the full Gibbs law yields
$D(P_{k,t}^{\nu}\Vert\nu^{\otimes N})/N\le C_\nu t^2$ for both signs.
Repulsive integrability follows also from the Hadamard lower bound;
attractive integrability on a larger interval was already shown.

Write $p_k(t)=N^{-1}\log Z_k^\nu(t)$. On either side of zero,
\[
 \frac{d}{dt}\frac{p_k(t)}t
 =\frac{t p_k'(t)-p_k(t)}{t^2}
 =\frac{D(P_{k,t}^{\nu}\Vert\nu^{\otimes N})}{Nt^2}
 \in[0,C_\nu].
\]
Integrating between zero and $t$, and multiplying by $t$, gives
\[
 0\le p_k(t)-t\mathbb E_{\nu^{\otimes N}}E_k\le C_\nu t^2
\]
for both signs, since both integration orientation and multiplication
reverse when $t<0$. Differentiability at zero follows from the two-sided
exponential neighborhood. This proves the claimed estimate with
$t_0=s/2$.
\end{proof}

\begin{theorem}[Product-law mean energy from repulsive pressure]
\label{att:general-product-energy}
For every smooth strictly positive probability density $\nu$ on
$\mathbb P^2$,
\[
 \mathbb E_{\nu^{\otimes N_k}}E_k\longrightarrow E(\nu).
\]
Together with the deterministic lower comparison, this convergence
also holds in $L^1$.
\end{theorem}

\begin{proof}
For $a>0$, the established unrestricted repulsive determinant LDP
\cite[Theorems 1.1 and 5.7]{prob_berman2017}, with volume reference
$\nu$ and the original metric, gives
\[
 -\frac1{aN}\log Z_k^\nu(-a)\longrightarrow
 J_\nu(a):=\inf_\mu\left[E(\mu)+\frac1a D(\mu\Vert\nu)\right].
\]
The energy normalization is unchanged by changing the volume reference;
only the entropy reference changes. If the repulsive theorem is stated
using a $\nu$-orthonormal basis, its Gram matrix relative to the fixed
$\sigma$-orthonormal basis has spectrum between $\min g$ and $\max g$.
The resulting additive change of $E_k$ is $O(k^{-1})$, so its limit
is unchanged. Lemma~\ref{att:quadratic-pressure}
gives, for $0<a\le t_0$,
\[
 J_\nu(a)\le\liminf_k\mathbb E_{\nu^N}E_k
 \le\limsup_k\mathbb E_{\nu^N}E_k\le J_\nu(a)+C_\nu a.
\]
We claim $J_\nu(a)\to E(\nu)$ as $a\downarrow0$. The upper bound
uses $\mu=\nu$. Approximate minimizers have relative entropy at most
$a(E(\nu)+o(1))$, since $E\ge0$, so they converge in total variation
to $\nu$. Lower semicontinuity of $E$ gives the reverse limiting
inequality. Letting $a\downarrow0$ proves convergence of the means.
For the $L^1$ statement, empirical convergence to $\nu$ and the
continuous-envelope lower comparison make negative deviations from
$E(\nu)$ tend to zero in probability. They are bounded by
$E(\nu)+\log N/k$. Convergence of the means then forces both deviations
to vanish in mean.
\end{proof}

The theorem is obtained from the repulsive sign, not from an assumed
negative-temperature limit. Its proof applies more generally whenever
the polarized manifold has the stated uniform normalized-section
exponential estimate and the repulsive determinant pressure theorem.
The separate toric assignment formula gives a direct constructive
calculation of the same product mean energy on invariant reference laws.

\subsection{Where the limiting empirical densities lie}

Fix $0<\gamma<1/3$. Let $\Phi$ be Lipschitz on $\mathcal P(X)$ for
the Wasserstein distance of any fixed smooth metric on $X$, and put
\[
 dP_{k,\gamma,\Phi}=Z_{k,\gamma,\Phi}^{-1}
       e^{N[\gamma E_k-\Phi(L_N)]}\,d\sigma^N.
\]
Finiteness and finite relative entropy follow from the strict scalar
integrability margin and boundedness of $\Phi$.

\begin{lemma}[A compact class for all directing densities]
\label{att:directing-class}
For these laws there are $r>1$, $c>0$, and a compact convex set
$\mathcal K\subset L^r(\sigma)$ of probability densities, all at least
$c$, with the following properties. Every full one-coordinate conditional
density belongs to $\mathcal K$. Every subsequential empirical limit
$\Pi$ is supported on the measures $g\sigma$ with $g\in\mathcal K$.
Weak probability topology and $L^r$ topology coincide on this compact set.
\end{lemma}

\begin{proof}
The conditional exponent consists of $-\gamma u(x)$ as in
Lemma~\ref{att:conditional-compactness}, plus
\[
 -V_{k,x_{-i}}(x),\qquad
 V_{k,x_{-i}}(x)=N[\Phi(L_N^x)-\Phi(L_N^{x_*})].
\]
If $L$ is a Lipschitz constant of $\Phi$, these functions have absolute
value at most $L\operatorname{diam}X$ and Lipschitz constant at most
$L$, independently of all parameters. Their closure is compact in
$C(X)$ by Arzel\`a--Ascoli. The psh compactness proof therefore makes
the normalized conditional densities a relatively compact family in
$L^r$, for any $r>1$ with $r\gamma<1/3$. Their normalization is at
most $e^{L\operatorname{diam}X}C_\gamma$, whereas their numerators are
at least $e^{-L\operatorname{diam}X}$. Hence the common lower bound
$c=e^{-2L\operatorname{diam}X}/C_\gamma$ holds.

Take the closed convex hull in $L^r$. It is compact: approximate the
original compact set by a finite epsilon-net; its convex hull lies in
the epsilon-neighborhood of a compact finite-dimensional polytope.
The norm bound, normalization, and positive lower bound persist.

Here is an elementary support argument, avoiding any assumption about
conditional de Finetti limits. For $h\in C(X)$ set
$b_h=\sup_{g\in\mathcal K}\int hg\,d\sigma$. A one-coordinate
conditional given previously revealed coordinates is a mixture of full
conditionals and therefore still belongs to $\mathcal K$. Thus
$\mathbb E[h(X_i)\mid X_1,\ldots,X_{i-1}]\le b_h$. Hoeffding's
bounded-increment supermartingale estimate gives
\[
 P_{k,\gamma,\Phi}(L_N h>b_h+\epsilon)
 \le \exp[-2N\epsilon^2/(\operatorname{osc}h)^2]
\]
when $h$ is nonconstant. Apply this to a countable dense collection of
continuous functions. Separation of a compact convex set of probability
measures shows that every limiting empirical measure belongs to
$\mathcal K$. Finally the continuous injection
$\mathcal K\subset L^r\to\mathcal P(X)$ is a homeomorphism onto
its image, by compactness and the Hausdorff property.
\end{proof}

\begin{lemma}[Uniform smooth reference approximation]
\label{att:smooth-net}
For every $\epsilon>0$, the compact set $\mathcal K$ above has a
finite set of smooth positive reference densities $g_1,\ldots,g_M$
such that every $g\in\mathcal K$ is assigned a reference satisfying
\[
 D(g\sigma\Vert g_a\sigma)<\epsilon,
 \qquad |E(g\sigma)-E(g_a\sigma)|<\epsilon.
\]
The references can always have the same positive lower bound $c$ and
a common $L^r$ bound, independently of $\epsilon$.
\end{lemma}

\begin{proof}
Apply the heat semigroup of a unitary-invariant metric. It preserves
probability, positivity, and the lower bound, and contracts $L^r$.
Strong continuity is uniform on compact sets. Its small-time images,
together with $\mathcal K$, have compact closure in $L^r$.
Relative entropy is jointly continuous on this closure: the map
$g\mapsto\log g$ is continuous into $L^{r'}$ since $g\ge c$, while
logarithmic growth is controlled by arbitrary small positive powers
of $g$. H\"older then treats the cross-entropy term, and
$g\mapsto\int g\log g$ is continuous by $L^r$ integrability.

For energy, normalized psh potentials have a uniform $L^{r'}$ bound
by their exponential integrability. The energy dual consequently gives
\[
 |E(g\sigma)-E(h\sigma)|\le C_{r'}\|g-h\|_r.
\]
This bound follows by taking the difference of the two suprema in the
energy dual and normalizing each test potential to have supremum zero.
For an explicit relative-entropy estimate, choose
$1<p\le\min(r,2)$. The elementary bound
$x\log x-x+1\le C_p|x-1|^p$ for $x\ge0$ gives
\[
 D(g\sigma\Vert h\sigma)
 \le C_p c^{1-p}\|g-h\|_p^p\qquad(h\ge c).
\]
A finite net in the heat-smoothed compact set now has the required
properties.
\end{proof}

\subsection{The actual tilted Gibbs mean energy}

\begin{theorem}[Entropy and energy limits of attractive Gibbs laws]
\label{att:gibbs-energy-limit}
For fixed $0<\gamma<1/3$ and Lipschitz $\Phi$, suppose a subsequence
of the empirical laws of $P_{k,\gamma,\Phi}$ converges to $\Pi$.
Then
\begin{align}
 \frac1N D(P_{k,\gamma,\Phi}\Vert\sigma^N)
       &\longrightarrow\int H(\mu)\,d\Pi(\mu),\\
 \mathbb E_{P_{k,\gamma,\Phi}}E_k
       &\longrightarrow\int E(\mu)\,d\Pi(\mu).
 \label{att:gibbs-energy-limit-eq}
\end{align}
The statements concern the actual attractive Gibbs laws. They do not
assert microscopic energy compactness for all symmetric laws of
bounded entropy and energy.
\end{theorem}

\begin{proof}
Lemma~\ref{att:directing-class} and the entropy-chaos lemma give the
entropy limit. Partition boundaries for the underlying spatial
partitions may be chosen $\sigma$-null, since the limiting first
marginal has an $L^r$ density.

Choose finite smooth references from Lemma~\ref{att:smooth-net}.
Because weak and strong topology agree on $\mathcal K$, finitely many
weak neighborhoods can be chosen so that every point of
$\mathcal K$ in a given neighborhood has the prescribed small relative
entropy and energy error to its reference. Choose radii with
$\Pi$-null boundaries and make the neighborhoods into disjoint cells
$A_1,\ldots,A_M$. Add a remaining cell $A_0$ with reference
$\nu_0=\sigma$; its limiting probability is zero. Put
$p_{k,a}=P_{k,\gamma,\Phi}(L_N\in A_a)$, and let $R_{k,a}$ be the
corresponding conditional law whenever its probability is nonzero.

The disjoint-event entropy identity and change of reference give
\begin{align*}
 \sum_a\frac{p_{k,a}}N D(R_{k,a}\Vert\nu_a^N)
 &=\frac1N D(P_{k,\gamma,\Phi}\Vert\sigma^N)
   +\frac{H((p_{k,a})_a)}N\\
 &\quad-\sum_a p_{k,a}\mathbb E_{R_{k,a}}
                    \int\log g_a\,dL_N.
\end{align*}
For this fixed finite family every $\log g_a$ is bounded continuous.
The entropy limit and convergence on the cells imply
\begin{equation}
 \limsup_k\sum_a\frac{p_{k,a}}N
                D(R_{k,a}\Vert\nu_a^N)
 \le\sum_a\int_{A_a}D(\mu\Vert\nu_a)\,d\Pi(\mu)<\epsilon.
 \label{att:mixture-entropy}
\end{equation}
There is no need for a separate estimate on the exceptional cell:
its nonnegative entropy term is included in this sum.

The two-sided pressure bound and the entropy variational inequality
imply, for either sign and a common $0<t\le t_0$,
\[
 \left|\mathbb E_{R_{k,a}}E_k-
                   \mathbb E_{\nu_a^N}E_k\right|
 \le Ct+\frac{D(R_{k,a}\Vert\nu_a^N)}{Nt}.
\]
The constants $C,t_0$ are uniform over every refinement of the net:
Lemma~\ref{att:quadratic-pressure} uses only the common $L^r$ bound.
Multiply by $p_{k,a}$ and sum. Theorem~\ref{att:general-product-energy}
handles the finitely many smooth references, and
\eqref{att:mixture-entropy} gives
\[
 \limsup_k\left|\mathbb E_{P_{k,\gamma,\Phi}}E_k
                   -\int E(\mu)\,d\Pi(\mu)\right|
 \le Ct+\epsilon/t+\epsilon.
\]
First refine the reference net, then let $t\downarrow0$. This proves
mean-energy convergence. The argument controls collective logarithmic
singularities through entropy and conditional regularity; no spectral
positive-moment bound is substituted for a logarithmic estimate.
\end{proof}

\subsection{The full Laplace principle and shrinking design constraints}

Set $F_\gamma(\mu)=H(\mu)-\gamma E(\mu)$ for finite-entropy
probability measures, and $F_\gamma=+\infty$ otherwise. This convention
never subtracts two infinite quantities: finite entropy implies finite
energy by the scalar exponential estimate. The following argument also
applies on an arbitrary positively polarized compact manifold, replacing
$1/3$ by its normalized-psh integrability threshold $\alpha$.

\begin{lemma}[Good rate and entropy-continuous energy recovery]
\label{att:good-rate}
For $0\le\gamma<\alpha$, $F_\gamma$ is lower semicontinuous and
has compact sublevels. For every $\gamma<a<\alpha$,
\begin{equation}
 F_\gamma(\mu)\ge(1-\gamma/a)H(\mu)
                          -(\gamma/a)\log C_a,
 \label{att:coercivity}
\end{equation}
where $C_a$ bounds $\int e^{-au}\,d\sigma$ over all normalized psh
potentials. If $\mu_j\to\mu$ weakly and
$H(\mu_j)\to H(\mu)<\infty$, then $E(\mu_j)\to E(\mu)$.
\end{lemma}

\begin{proof}
The case $\gamma=0$ for the first assertion is ordinary entropy lower
semicontinuity. Take $0<\gamma<\alpha$. Write $\mathcal A$ for
Aubin--Mabuchi potential energy, extended by $-\infty$ on infinite-energy
normalized psh functions. The measure energy dual gives
\begin{equation}
 F_\gamma(\mu)=
 \inf_{\substack{u\in\operatorname{PSH}(X,\omega_0)\\\sup u=0}}
 \left\{D(\mu\Vert e^{-\gamma u}\sigma)
                       -\gamma\mathcal A(u)\right\}.
 \label{att:compact-potential-dual}
\end{equation}
Here relative entropy against a finite, not necessarily normalized,
measure $\nu$ means $\int\log(d\mu/d\nu)\,d\mu$ and has the
variational formula
\[
 D(\mu\Vert\nu)=\sup_{f\in C(X)}
       \left(\int f\,d\mu-\log\int e^f\,d\nu\right).
\]
For finite $H(\mu)$, the exponential bound and the entropy inequality
make $\int|u|\,d\mu$ finite, uniformly over normalized $u$.
Thus $D(\mu\Vert e^{-\gamma u}\sigma)=H(\mu)+\gamma\int u\,d\mu$,
and \eqref{att:compact-potential-dual} follows from the energy dual.
Potentials with $\mathcal A(u)=-\infty$ contribute $+\infty$.

The equality also holds at infinite entropy. To see this and obtain a
uniform bound, choose $\gamma<a<\alpha$ and put $r=a/\gamma>1$.
For every bounded continuous $f$, H\"older gives
\[
 \int e^{f/r'}e^{-\gamma u}\,d\sigma
 \le \left(\int e^f\,d\sigma\right)^{1/r'} C_a^{1/r}.
\]
Taking the relative-entropy variational supremum therefore yields
\[
 D(\mu\Vert e^{-\gamma u}\sigma)
 \ge (1-\gamma/a)H(\mu)-(\gamma/a)\log C_a.
\]
This is valid with infinite values. Since $\mathcal A(u)\le0$ when
$\sup u=0$, every objective in \eqref{att:compact-potential-dual}
satisfies the same bound. Consequently the infimum is $+\infty$ when
$H(\mu)=+\infty$, and \eqref{att:coercivity} follows in all cases.

The normalized psh family is compact in $L^1$, and $\mathcal A$ is
upper semicontinuous on it, including its value $-\infty$; these are
standard variational pluripotential inputs \cite{BBGZ2013}.
If $u_j\to u$ in this compact family, the strict exponent margin
and Vitali's theorem give
$e^{-\gamma u_j}\to e^{-\gamma u}$ in $L^1(\sigma)$.
The displayed entropy variational formula makes $D(\mu\Vert\nu)$
jointly lower semicontinuous under weak convergence of its two measure
arguments. Hence the objective in \eqref{att:compact-potential-dual}
is jointly lower semicontinuous in $(\mu,u)$. Minimizing over the
compact potential space proves lower semicontinuity of $F_\gamma$:
choose approximate minimizing potentials along any sequence and extract
a convergent subsequence. Compactness of $\mathcal P(X)$ then makes
every free-energy sublevel compact.

Finally choose any fixed $0<\gamma<\alpha$. Under the assumptions
of the last assertion, lower semicontinuity of $F_\gamma$ gives
\[
 H(\mu)-\gamma E(\mu)
 \le H(\mu)-\gamma\limsup_j E(\mu_j).
\]
Thus $\limsup E(\mu_j)\le E(\mu)$; lower semicontinuity of $E$
gives the reverse lower limit. All energies are finite and uniformly
bounded by the entropy/exponential estimate. This proves energy recovery.
\end{proof}

\begin{theorem}[Full attractive determinant Laplace principle]
\label{att:full-laplace}
For every fixed $0\le\gamma<1/3$ and every
$\Phi\in C(\mathcal P(\mathbb P^2))$,
\begin{equation}
 \lim_{k\to\infty}\frac1N\log\int
       e^{N[\gamma E_k-\Phi(L_N)]}\,d\sigma^N
 =-\inf_{\mu\in\mathcal P(X)}[F_\gamma(\mu)+\Phi(\mu)].
 \label{att:full-laplace-eq}
\end{equation}
The empirical laws of the normalized attractive ensemble satisfy a
full speed-$N$ LDP with good rate
$F_\gamma-\inf F_\gamma$.
\end{theorem}

\begin{proof}
The case $\gamma=0$ is Sanov. Assume $0<\gamma<1/3$ and first take
$\Phi$ Lipschitz. Its finite-level Gibbs variational identity is
\[
 N^{-1}\log Z_{k,\gamma,\Phi}
 =\gamma\mathbb E E_k-N^{-1}D(P_{k,\gamma,\Phi}\Vert\sigma^N)
                         -\mathbb E\Phi(L_N).
\]
Every subsequence has an empirically convergent further subsequence.
Theorem~\ref{att:gibbs-energy-limit} identifies its limit as
$\int[\gamma E-H-\Phi]\,d\Pi$, at most
$\sup_\mu[\gamma E(\mu)-H(\mu)-\Phi(\mu)]$.
For the opposite bound, insert a smooth positive product law $\mu^N$
in the variational principle and use
Theorem~\ref{att:general-product-energy}.

Smooth positive references suffice for the variational infimum.
For any finite-entropy $\mu$, unitary heat averaging gives smooth
positive $\mu_\epsilon\to\mu$ with $H(\mu_\epsilon)\to H(\mu)$:
entropy contracts and is lower semicontinuous. Energy also converges,
since its unitary invariance and convexity give
$E(\mu_\epsilon)\le E(\mu)$ and its lower semicontinuity gives the
reverse lower limit. Thus the lower bound extends to all relevant
measures.

Lipschitz functions are uniformly dense in the continuous functions
on the compact metric space $\mathcal P(X)$. Both finite pressure and
the variational limit are 1-Lipschitz in the uniform norm of $\Phi$.
This proves the assertion for every continuous tilt. The full Laplace
principle on a compact probability space is equivalent to the stated
LDP.
\end{proof}

\begin{theorem}[Full attractive LDP on shrinking design tubes]
\label{att:design-laplace}
Fix any design order $t\ge1$ and $0<a<1/2$. Let
\[
 \begin{aligned}
 \mathcal D_t&=\{\mu:M_t(\mu)=I/\tbinom{t+2}{2}\},\\
 \Omega_N&=\{\|M_t(L_N)-I/\tbinom{t+2}{2}\|_{\rm HS}<N^{-a}\},
 \qquad Q_N=\sigma^N(\cdot\mid\Omega_N).
 \end{aligned}
\]
For every fixed $0\le\gamma<1/3$ and every continuous empirical tilt,
\begin{equation}
 \lim_{k\to\infty}\frac1N\log\int
 e^{N[\gamma E_k-\Phi(L_N)]}\,dQ_N
 =-\inf_{\mu\in\mathcal D_t}[H(\mu)-\gamma E(\mu)+\Phi(\mu)].
 \label{att:design-laplace-eq}
\end{equation}
The normalized empirical laws satisfy the full LDP with good rate
$F_\gamma+\iota_{\mathcal D_t}-\inf_{\mathcal D_t}F_\gamma$.
\end{theorem}

\begin{proof}
The reference probability of $\Omega_N$ tends to one by the finite
moment variance bound, since $N^{1-2a}\to\infty$.
For the upper bound, fix a closed moment tube of radius $\delta>0$.
For large $N$ it contains $\Omega_N$. The unrestricted LDP and its
bounded continuous tilt bound the upper pressure by the variational
infimum on that fixed tube. Let $\delta\downarrow0$, using compactness
and lower semicontinuity of $F_\gamma+\Phi$.

For the lower bound take a smooth positive $\mu\in\mathcal D_t$.
Under $\mu^N$, the event $\Omega_N$ has probability tending to one.
Theorem~\ref{att:general-product-energy} and its $L^1$ version show
that conditioning on this event preserves convergence of mean energy
to $E(\mu)$. Its relative entropy per particle to $Q_N$ tends to
$H(\mu)$: write out the conditioned product density, use boundedness
of $\log(d\mu/d\sigma)$, and note that both conditioning logarithms
are $o(N)$. The variational principle gives the desired lower bound.
Finally, unitary heat averaging preserves an exact projective design,
since its moment target is scalar and the heat operation is an average
of simultaneous unitary pushforwards. It approximates any finite-entropy
member of $\mathcal D_t$ with converging entropy and energy as above.
The full continuous-tilt Laplace principle yields the LDP.
\end{proof}

\begin{corollary}[Polarized manifolds and all projective dimensions]
\label{att:polarized-extension}
Let $Y$ be a compact connected complex manifold, let $(L,h)$ be a
positive Hermitian line bundle with curvature $\omega$, and let $\sigma$
be a smooth strictly positive probability volume. For all sufficiently large $k$, choose an $L^2(h^k,\sigma)$
orthonormal basis of $H^0(Y,L^k)$ and let $\Delta_k$ be its determinant,
$N_k=h^0(Y,L^k)$, and $E_k=-(kN_k)^{-1}\log\|\Delta_k\|^2$.
Write $E_\omega$ for the pluricomplex measure energy with potential
energy normalized to vanish at zero, and put
\[
 \alpha(\omega,\sigma)=\sup\left\{a>0:
 \sup_{\substack{u\in\operatorname{PSH}(Y,\omega)\\ \sup u=0}}
             \int e^{-au}\,d\sigma<\infty\right\}.
\]
For every fixed $0\le\gamma<\alpha(\omega,\sigma)$, the unrestricted
attractive determinant ensemble has the full speed-$N_k$ LDP with good
rate
\[
 H_\sigma-\gamma E_\omega-\inf(H_\sigma-\gamma E_\omega).
\]
Its Laplace principle holds for every continuous empirical-law tilt.
For $Y=\mathbb P^q$, $q\ge1$, $L=\mathcal O(q+1)$, and Fubini--Study probability,
the admissible range is $0\le\gamma<1/(q+1)$; in that range the same
shrinking-tube conclusion holds for every fixed projective design order,
with target $I/\binom{q+t}{q}$.
\end{corollary}

\begin{proof}
The proof above uses only compactness of normalized psh functions,
a uniform exponential estimate below $\alpha$, and the repulsive determinant
pressure theorem. These hold for the stated polarized manifold; the
repulsive input is \cite[Theorem 5.7]{prob_berman2017}. The good-rate
argument of Lemma~\ref{att:good-rate} uses only the same integrability
margin and the basic potential-energy duality of \cite{BBGZ2013}.
For clarity, there is no hidden homogeneity assumption in the reference
approximation. Use a smooth uniformly elliptic diffusion whose invariant
volume is $\sigma$; its heat semigroup contracts all $L^p(\sigma)$,
preserves positive lower bounds, and has a positive smooth kernel.
It therefore supplies the finite smooth reference nets. For an arbitrary
finite-entropy density its heat regularizations have converging entropy,
and their energies converge by Lemma~\ref{att:good-rate}. This replaces
the unitary averaging argument without imposing homogeneity on $Y$. The lower
Hadamard bound is $E_k\ge-o(1)$, since the smooth positive Bergman kernel
has polynomial growth in $k$. Every other step is unchanged. This proves
the unrestricted assertion; no moment constraint on a general manifold
is assumed here.

On projective space put $D=q+1$. The normalized psh threshold is
$1/D$, as can also be seen directly. Choose a maximum point of $u$ and
disintegrate Haar measure along projective lines through it. The
conditional radial density on a line is $q(1-x)^{q-1}dx$, bounded by
$q$ times line Haar density. The Riesz representation of $u/D$ on the
line, normalized at the maximum point, and Jensen's inequality give
\[
 \int_{\mathbb P^q}e^{-au}\,d\sigma\le\frac{q}{1-Da},
 \qquad 0\le a<1/D.
\]
Indeed a unit linear form has squared modulus uniform on $[0,1]$ under
line Haar measure. Conversely $u=D\log|z_0|^2$ has supremum zero and
fails exponential integrability at $a=1/D$, proving sharpness of this
one-coordinate threshold. Finally unitary heat averaging on
$\mathbb P^q$ preserves each exact projective design. The proof of
Theorem~\ref{att:design-laplace} applies verbatim in every dimension.
\end{proof}

\begin{corollary}[Zero pressure and the unique Haar equilibrium]
\label{att:haar-minimizer}
On $\mathbb P^q$ with anticanonical polarization, let
$0\le\gamma<1/(q+1)$. Then
\[
 \lim_k N_k^{-1}\log Z_{k,\gamma}=0,
 \qquad H(\mu)-\gamma E(\mu)\ge0,
\]
and equality in the second assertion holds exactly for $\mu=\sigma$.
The unconditioned attractive rate is therefore $F_\gamma$ itself.
For every fixed-order shrinking design tube the rate is
$F_\gamma+\iota_{\mathcal D_t}$, and its unnormalized pressure is also
zero.
\end{corollary}

\begin{proof}
Corollary~\ref{ac:tubes} gives empirical convergence to $\sigma$.
The entropy and energy limit in Theorem~\ref{att:gibbs-energy-limit},
with the same proof on projective space of every dimension as explained
in Corollary~\ref{att:polarized-extension}, then gives
\[
 N^{-1}D(P_{k,\gamma}\Vert\sigma^N)\longrightarrow0,
 \qquad \mathbb E_{P_{k,\gamma}}E_k\longrightarrow E(\sigma)=0.
\]
The finite Gibbs identity yields zero pressure. The full Laplace principle
therefore identifies $\inf F_\gamma=0$, attained by $\sigma$.
For $\gamma=0$ these assertions are also immediate from entropy.

To prove uniqueness, let $\mu\ne\sigma$. Choose a real smooth $f$
with $\mu(f)\ne\sigma(f)$, and a weak neighborhood $U$ of $\mu$
on which $|\nu(f)-\sigma(f)|>s$ for some $s>0$. The exponential
bound of Theorem~\ref{ac:concentration} gives
$P_{k,\gamma}(L_N\in U)\le2e^{-cN}$ with $c>0$.
The LDP lower bound on this open set gives
$-\inf_U F_\gamma\le-c$. In particular $F_\gamma(\mu)\ge c>0$.
The use of the exponential estimate is essential: convergence in
probability alone would not imply strict positivity of the rate away
from $\sigma$. Finally $\sigma$ is an exact design of every order,
so restriction to a design fibre preserves the unique zero and the
normalization in Theorem~\ref{att:design-laplace}.
\end{proof}

\paragraph{Range and the next problem.}
The proven interval is a fixed positive attractive interval,
$0\le\gamma<1/3$, for the natural complex determinant and for every
fixed-order shrinking design constraint. It includes all continuous
empirical-law tilts. The scalar section threshold is used twice:
for integrability and compactness of attractive conditional densities.
A result near $\gamma=1$ requires a replacement for that conditional
compactness mechanism; first-moment macroscopic coercivity alone does
not provide it. The toric assignment formula separately makes the product energy
computable by finite optimal assignments, while the polarized extension
shows that the entropy-and-conditional-compactness mechanism is not
specific to a homogeneous projective space.

\section{A nonlinear Green formula from coherent moments}
\label{cy:section}

The product-law energy limit has a differential consequence. It
recovers the Calabi--Yau potential itself from random interpolation
sections, rather than only recovering its variational energy. The
representation below is the determinant formula proposed in
\cite[Conjecture 6.1]{prob_berman2017}. Its finite-level convexity
mechanism reduces a two-coordinate determinant to the projective line.
The underlying product-energy theorem addresses the determinant instance
of \cite[Problem 4.10]{prob_berman2017}, rather than arbitrary singular
Hamiltonians. A later conditional result obtained the same potential
formula from a uniform complex zero-free neighborhood of the partition
functions \cite[Theorems 3.2--3.3]{Berman2023zeta}. The argument here
uses the real-parameter energy estimate and finite-level convexity;
it makes no assertion about the complex zeros of those functions.
In complex dimension one the representation was already unconditional
by bosonization, as recalled in \cite[Remark 3.4]{Berman2023zeta}.
The argument below applies in every complex dimension under the stated
smooth positive hypotheses.

Let $X$ be a compact connected complex manifold of dimension $q\ge1$,
and let $(L,h)$ be a positive Hermitian line bundle with curvature
$\omega=dd^c\phi_0$, where $dd^c=i\partial\bar\partial/(2\pi)$.
Put $V=\int_X\omega^q$. For all sufficiently large $k$, let
$N=N_k=h^0(X,L^k)$ and choose an orthonormal basis for a fixed smooth
positive reference volume $\sigma$. Denote its evaluation determinant
by $\Delta_k$ and keep this basis fixed when varying a probability
measure $\mu$. As before,
\[
 E_k=-\frac1{kN}\log\|\Delta_k\|^2,
 \qquad \overline E_k(\mu)=\int E_k\,d\mu^{\otimes N}.
\]
All norms use the fixed metric $h$, in every determinant coordinate.

\begin{lemma}[Convexity of the product mean energy]
\label{cy:finite-convexity}
For every sufficiently large fixed $k$, the functional
$\overline E_k$ is convex on smooth strictly positive probability
volumes. If $\eta$ is a smooth signed measure of total mass zero and
$\mu_t=\mu+t\eta$ remains strictly positive, then
\begin{equation}
 \frac{d}{dt}\overline E_k(\mu_t)\bigg|_{t=0}
 =-\int_X U_{k,\mu}(x)\,d\eta(x),\qquad
 U_{k,\mu}(x)=\frac1k\int_{X^{N-1}}
        \log\|\Delta_k(x,\mathbf y)\|^2\,d\mu^{N-1}(\mathbf y).
 \label{cy:first-variation}
\end{equation}
\end{lemma}

\begin{proof}
The logarithm of a nonzero holomorphic section is locally integrable.
On the compact product manifold this makes the determinant logarithm
absolutely integrable against every product of smooth volume densities.
Thus expansion of $(\mu+t\eta)^N$ gives an ordinary polynomial in
$t$, and differentiation gives \eqref{cy:first-variation} and
\[
 \frac{d^2}{dt^2}\overline E_k(\mu_t)
 =-\frac{N-1}{k}\int_{X^{N-2}}
   \left[\int_{X^2}\log\|\Delta_k(x,y,\mathbf z)\|^2
                      \,d\eta(x)d\eta(y)\right]d\mu_t^{N-2}(\mathbf z).
\]
For almost every retained tuple $\mathbf z$, its evaluation rows have
rank $N-2$. Indeed, a nonzero full determinant gives one such tuple,
and rank failure is a proper analytic locus. The two-dimensional
annihilator of these rows is a pencil of sections $a,b$ of $L^k$.
Expansion of the determinant gives a nonzero constant times
$a(x)b(y)-b(x)a(y)$, with the retained norm factors constant in $x,y$.
Set
\[
 r(x)=\|a(x)\|^2+\|b(x)\|^2,\qquad
 P(x)=P_{[a(x):b(x)]}\in\mathbb P^1.
\]
The common zero set has volume zero. Off it,
\[
 \log\|\Delta_k(x,y,\mathbf z)\|^2
 =c_{\mathbf z}+\log r(x)+\log r(y)
                     +\log[1-\operatorname{tr}(P(x)P(y))].
\]
Since $\eta(X)=0$, the first three terms contribute zero to the signed
double integral. The remaining kernel is positive definite after
changing its sign:
\begin{equation}
 K(P,Q)=-\log[1-\operatorname{tr}(PQ)]
   =\sum_{j\ge1}\frac1j
            \operatorname{tr}(P^{\otimes j}Q^{\otimes j}).
 \label{cy:projective-kernel}
\end{equation}
Each finite sum has nonnegative quadratic form because
\[
 \iint\operatorname{tr}(P(x)^{\otimes j}P(y)^{\otimes j})
                         \,d\eta(x)d\eta(y)
 =\left\|\int P(x)^{\otimes j}\,d\eta(x)\right\|_{\rm HS}^2\ge0.
\]
To justify the infinite sum with a signed measure, first apply Fubini
to the absolute determinant logarithm. For almost every $\mathbf z$
its double integral against $|\eta|\otimes|\eta|$ is finite. The
functions $\log r$ are also integrable: they are bounded above, and a
nonzero component section bounds their negative parts. The displayed
kernel identity therefore makes $K$ integrable against
$|\eta|\otimes|\eta|$. Its increasing finite sums are bounded by $K$,
so dominated convergence proves nonnegativity of the full quadratic
form. More precisely, if $P_{\mathbf z}$ denotes the pencil map for
these retained rows, the computation proves the exact hierarchy
\begin{equation}
 \frac{d^2}{dt^2}\overline E_k(\mu_t)
 =\frac{N-1}{k}\int_{X^{N-2}}
   \sum_{j\ge1}\frac1j
   \left\|\int_X P_{\mathbf z}(x)^{\otimes j}\,d\eta(x)
          \right\|_{\rm HS}^2
                  \,d\mu_t^{N-2}(\mathbf z)\ge0.
 \label{cy:hessian-hierarchy}
\end{equation}
Tonelli justifies averaging the nonnegative terms over $\mathbf z$.
Thus coherent moment discrepancies form the finite-level Hessian.
No convergence of these Hessians is required for the result below.
\end{proof}

\begin{theorem}[Nonlinear Green representation]
\label{cy:green-formula}
Let $\mu$ be any smooth strictly positive probability volume on $X$.
There is a unique smooth potential $\phi_\mu$, normalized by
$\int_X\phi_\mu\,d\mu=0$, such that
\[
 \frac{(\omega+dd^c\phi_\mu)^q}{V}=\mu,
 \qquad \omega+dd^c\phi_\mu>0.
\]
The mean-zero determinant potentials
\begin{equation}
 \phi_{k,\mu}(x)=\frac1k\int_{X^{N_k-1}}
       \log\|\Delta_k(x,\mathbf y)\|^2\,d\mu^{N_k-1}(\mathbf y)
      -C_{k,\mu},
 \quad
 C_{k,\mu}=\frac1k\int_{X^{N_k}}\log\|\Delta_k\|^2\,d\mu^{N_k}
 \label{cy:potential-formula}
\end{equation}
converge to $\phi_\mu$ in $L^p(X,\sigma)$, equivalently in
$L^p(X,\mu)$, for every fixed $1\le p<\infty$. The formula is independent of the chosen basis after
the displayed centering.
\end{theorem}

\begin{proof}
Existence and uniqueness of the smooth solution are Yau's theorem
\cite{Yau1978}. The energy dual is attained at that solution:
\[
 E_\omega(\mu)=\mathcal A(\phi_\mu)-\int\phi_\mu\,d\mu.
\]
For a smooth signed variation $\eta$ of mass zero, the normalized
solutions for $\mu+t\eta$ vary smoothly for small $t$. This follows
from the elliptic implicit function theorem: the linearization of the
normalized Monge--Amp\`ere map is the Laplacian of the positive metric,
an isomorphism modulo constants onto zero-integral variations.
Differentiating the last energy identity and using
$d\mathcal A_\phi(v)=\int v\,\operatorname{MA}_\omega(\phi)$ gives
\begin{equation}
 \frac{d}{dt}E_\omega(\mu+t\eta)\bigg|_{t=0}
                 =-\int\phi_\mu\,d\eta.
 \label{cy:energy-derivative}
\end{equation}

The product mean-energy theorem
\ref{att:general-product-energy}, in its polarized form explained in
Corollary~\ref{att:polarized-extension}, gives
$\overline E_k(\mu+t\eta)\to E_\omega(\mu+t\eta)$ for every
fixed sufficiently small $t$. The functions of $t$ on the left are
convex by Lemma~\ref{cy:finite-convexity}. Hence the secant inequalities
\[
 \frac{\overline E_k(\mu)-\overline E_k(\mu-\varepsilon\eta)}\varepsilon
 \le \frac{d}{dt}\overline E_k(\mu+t\eta)\bigg|_{t=0}
 \le
 \frac{\overline E_k(\mu+\varepsilon\eta)-\overline E_k(\mu)}\varepsilon
\]
show, by first taking $k\to\infty$ and then $\varepsilon\downarrow0$,
that their derivatives converge to \eqref{cy:energy-derivative}.
Because constants pair to zero with $\eta$, this says
\[
 \int\phi_{k,\mu}\,d\eta\longrightarrow\int\phi_\mu\,d\eta
 \quad\text{for every smooth zero-mass signed }\eta.
\]

For completeness, these weak identities determine an $L^1$ limit.
The averaged logarithm in \eqref{cy:potential-formula} defines an
$L^1$ function, by absolute logarithmic integrability and Fubini.
In distributions it satisfies
$dd^c\phi_{k,\mu}\ge-\omega$, since each determinant slice is a
section of $L^k$. Its psh representative is therefore an
$\omega$-psh potential. Centering makes its $\mu$ mean zero.
The normalized-psh exponential estimate gives a uniform constant $B_\mu$
such that every $\omega$-psh $v$ with $\sup v=0$ satisfies
$\int(-v)\,d\mu\le B_\mu$. Consequently
\[
 0\le\sup_X\phi_{k,\mu}\le B_\mu.
\]
Compactness of normalized psh potentials makes $\phi_{k,\mu}$ relatively
compact in $L^1(\sigma)$. Every cluster point has mean zero against
$\mu$, and its integrals against all smooth zero-mass signed measures
agree with those of $\phi_\mu$. Their difference is constant as a
distribution, and normalization makes that constant zero. Every cluster
point is therefore $\phi_\mu$, proving $L^1$ convergence of the full
sequence. The same uniform exponential estimate, together with the
uniform upper bounds, makes $|\phi_{k,\mu}|^p$ uniformly integrable
for every finite $p$. Vitali therefore upgrades the convergence to
$L^p$. A fixed basis change multiplies the determinant by a nonzero constant;
its logarithm cancels exactly in \eqref{cy:potential-formula}.
\end{proof}

\begin{corollary}[Random interpolation divisors and sampling error]
\label{cy:interpolation-sampling}
Draw $\mathbf Y=(Y_2,\ldots,Y_N)$ independently from $\mu$.
Almost surely the section
$s_{\mathbf Y}(x)=\Delta_k(x,\mathbf Y)$ is nonzero and is the unique
section vanishing at these $N-1$ points, up to scale. Define its
mean-zero potential
\[
 u_{k,\mathbf Y}(x)=\frac1k\log\|s_{\mathbf Y}(x)\|^2
                 -\int_X\frac1k\log\|s_{\mathbf Y}\|^2\,d\mu.
\]
Then
\begin{equation}
 \mathbb E\frac1k[Z_{s_{\mathbf Y}}]
       =\omega+dd^c\phi_{k,\mu}
       \longrightarrow\omega+dd^c\phi_\mu
 \label{cy:divisor-limit}
\end{equation}
weakly as currents. Moreover, if $u^{(1)},\ldots,u^{(M)}$ are independent
copies of $u_{k,\mathbf Y}$, then, for every $a>0$ with
\[
 C_a(\mu):=\sup_{\substack{v\in\operatorname{PSH}(X,\omega)\\\sup v=0}}
                      \int e^{-av}\,d\mu<\infty,
\]
one has the uniform sampling estimate
\begin{equation}
 \mathbb E\left\|\frac1M\sum_{j=1}^M u^{(j)}-\phi_{k,\mu}
             \right\|_{L^2(\mu)}^2
       \le\frac{2C_a(\mu)}{a^2M}.
 \label{cy:sampling-bound}
\end{equation}
On the anticanonical projective plane, if $\mu=g\sigma$, one may take
$a=1/6$, obtaining the explicit bound $192\|g\|_\infty/M$.

Exact centering integrals can also be replaced by sampling. For each
section independently draw $J\ge1$ calibration points from $\mu$
and subtract their mean log squared norm in place of its $\mu$ integral.
Then the left side of \eqref{cy:sampling-bound} is bounded by
\begin{equation}
 \frac{2C_a(\mu)}{a^2M}\left(1+\frac1J\right).
 \label{cy:sample-centered-bound}
\end{equation}
\end{corollary}

\begin{proof}
Independence of $N-1$ evaluation rows fails only on a proper analytic
subset, so the interpolation assertion holds almost surely. The
Poincar\'e--Lelong formula in our convention is
$k^{-1}[Z_s]=\omega+dd^c(k^{-1}\log\|s\|^2)$.
Taking expectations and using \eqref{cy:potential-formula} proves the
identity of currents; $L^1$ convergence of potentials proves its limit.

For the quantitative statement subtract the supremum from the section
potential and write $v\le0$ for the result. Centering at its $\mu$ mean
gives $u=v-\mu(v)$, so
\[
 \|u\|_{L^2(\mu)}^2\le\int v^2\,d\mu
              \le\frac2{a^2}\int e^{-av}\,d\mu
              \le\frac{2C_a(\mu)}{a^2}.
\]
Here $x^2\le2a^{-2}e^{ax}$ for $x\ge0$ follows from the exponential
series. The random functions form a square-integrable random variable
in the separable Hilbert space $L^2(\mu)$, and their mean is
$\phi_{k,\mu}$. The elementary variance identity for independent
Hilbert-space samples gives \eqref{cy:sampling-bound}.
On $\mathbb P^2$, the scalar bound gives
$C_{1/6}(\mu)\le(8/3)\|g\|_\infty$, yielding the stated constant.

For the sample-centered version, write the calibration error for section
$j$ as $e_j=J^{-1}\sum_{\ell=1}^J u^{(j)}(Z_{j\ell})$.
Conditionally on the interpolation sections it has mean zero and
variance at most $2C_a(\mu)/(a^2J)$, and the $e_j$ are conditionally
independent. The final estimator is
$M^{-1}\sum_j u^{(j)}-M^{-1}\sum_j e_j$.
The second term is a random constant of squared $L^2(\mu)$ norm
$|M^{-1}\sum_j e_j|^2$. Its cross term with the first centered
estimation error has conditional mean zero. Adding variances gives
\eqref{cy:sample-centered-bound}.
\end{proof}

\paragraph{Computational access and scope.}
A direct sampling procedure takes a basis of $H^0(X,L^k)$ that can be
evaluated, the Hermitian norm, and independent sample access to $\mu$.
For each repetition, draw $N_k-1$ interpolation points, form their
evaluation matrix, and choose a nonzero null vector. Evaluate the
resulting section, subtract its mean log squared norm on $J$ fresh calibration
points, divide by $k$, and average over the $M$ repetitions. Changing
the normalization of any null vector cancels in this subtraction.
Thus no exact integration oracle is needed. The theorem specifies the
sampling error of this ideal arithmetic procedure, without an assertion
of efficient sample access or finite-precision conditioning.

The sampling bound controls the error at a fixed polynomial level.
Theorem~\ref{cy:green-formula} separately proves convergence of the
level-dependent bias, without an effective rate in $k$. Taking
$k\to\infty$ and $M\to\infty$, even with $J=1$, gives convergence
in expected squared $L^2(\mu)$ to the Calabi--Yau potential: the error
is the sampling variance plus $\|\phi_{k,\mu}-\phi_\mu\|_2^2$.
No uniform-norm or higher-derivative convergence is asserted.
If the canonical bundle is trivial and $\mu$ is the normalized volume
of a nowhere-vanishing holomorphic top form, the limiting K\"ahler
form is the Ricci-flat metric in $c_1(L)$. In particular the K3 surface
of Theorem~\ref{geo:k3-lattice}, equipped with any ample polarization,
is an algebraic test case for the construction. This application uses
its Calabi--Yau structure; it does not assert an effective polynomial-
level bias bound or a practical sampling complexity for that surface.

\section{Curvature, concentration, and exact susceptibilities}
\label{ac:section}

The attractive determinant measure has a geometric positivity property
before any mean-field limit is taken.  Its density is the canonical
volume of a positively curved metric on the anticanonical line bundle
of the full configuration space.  The resulting estimate controls
empirical observables uniformly in the particle number, including
shrinking projective-moment tubes.

Write $X=\mathbb P^q$, $D=q+1$, $L=-K_X=\mathcal O(D)$,
$N=N_k=\binom{Dk+q}{q}$, and let $\sigma$ be Fubini--Study
probability.  Give $L$ its invariant metric $h_0=e^{-\phi_0}$, and
choose an orthonormal basis of $H^0(X,L^k)$ for $(h_0^k,\sigma)$.
Its evaluation determinant is a nonzero section $\Delta_k$ of
$\boxtimes_{i=1}^N L^k$ on $M=X^N$.  Put
\[
 w_{k,\gamma}=\|\Delta_k\|_{h_0^{k\boxtimes N}}^{-2\gamma/k},
 \qquad
 Z_{k,\gamma}=\int_M w_{k,\gamma}\,d\sigma^{\otimes N},
 \qquad
 dP_{k,\gamma}=Z_{k,\gamma}^{-1}w_{k,\gamma}
                                      \,d\sigma^{\otimes N}.
\]
The value on the zero divisor is immaterial whenever the partition
function is finite.  Additive constants in $\phi_0$ or a constant
rescaling of $\Delta_k$ do not change $P_{k,\gamma}$.

The metric convention in the estimates below is explicit.  Set
\[
 G_X=i\partial\bar\partial\phi_0=2\pi\omega_0,
 \qquad \omega_0=D\omega_{\rm FS},
 \qquad G_M=\sum_{i=1}^N\pi_i^*G_X.
\]
Here $dd^c=i\partial\bar\partial/(2\pi)$.  All covector norms
and Hessian comparisons in this section use the Hermitian metric with
K\"ahler form $G_X$ or $G_M$, rather than $\omega_0$.

\subsection{The anticanonical variance estimate}

\begin{lemma}[Canonical-volume Poincar\'e inequality]
\label{ac:canonical-poincare}
Let $(M,G)$ be a connected compact K\"ahler manifold and let
$e^{-\Phi}$ be a smooth metric on $-K_M$ satisfying
$i\partial\bar\partial\Phi\ge aG$, where $a>0$.
Let $\mu_\Phi$ be its canonical volume, normalized to mass one.
Then every smooth complex-valued function $F$ satisfies
\begin{equation}
 \int_M|F-\mu_\Phi(F)|^2\,d\mu_\Phi
 \le \frac1a\int_M|\bar\partial F|_G^2\,d\mu_\Phi.
 \label{ac:canonical-poincare-eq}
\end{equation}
\end{lemma}

\begin{proof}
Regard $F$ as a section of
$K_M\otimes(-K_M)\simeq\mathcal O_M$.  In a holomorphic coordinate
chart, its squared $L^2$ norm as a $(-K_M)$-valued top-degree form is
$|F|^2e^{-\Phi}$ times coordinate volume: the metric factors from
$K_M$ and the auxiliary volume of $G$ cancel.  The analogous norm
of $\bar\partial F$ is
$|\bar\partial F|_G^2e^{-\Phi}$ times coordinate volume.

The Bochner--Kodaira identity on $(-K_M)$-valued $(\dim M,1)$-forms
has curvature term $[i\Theta(-K_M),\Lambda_G]\ge aI$.
Consequently the $L^2$ estimate for the minimal solution of
$\bar\partial U=\bar\partial F$ gives
\[
 \|U\|_{L^2(e^{-\Phi})}^2
 \le a^{-1}\|\bar\partial F\|_{L^2(G,e^{-\Phi})}^2.
\]
This is the compact K\"ahler form of H\"ormander's estimate
\cite{Hormander1965}; it also follows directly from the preceding
curvature lower bound and the Hilbert-space minimal-solution argument.
The holomorphic sections of $K_M\otimes(-K_M)$ are exactly the
constant functions.  Thus the minimal solution is
$U=F-\mu_\Phi(F)$.  Dividing both sides by the mass of the canonical
volume proves the assertion.
\end{proof}

\begin{proposition}[Attractive covariance and exponential tilts]
\label{ac:variance}
Suppose $0\le\gamma<1$ and $Z_{k,\gamma}<\infty$.
Then, for every smooth function $F$ on $X^N$,
\begin{equation}
 \operatorname{Var}_{P_{k,\gamma}}(F)
 \le\frac1{1-\gamma}
       \int|\bar\partial F|_{G_M}^2\,dP_{k,\gamma}.
 \label{ac:variance-eq}
\end{equation}
For a real $f\in C^\infty(X)$, put
\[
 A_f=\sup_X|\bar\partial f|_{G_X}^2,
 \qquad
 -C_fG_X\le i\partial\bar\partial f\le C_fG_X.
\]
Here $C_f$ is any nonnegative constant satisfying the displayed
inequalities.  If
$dP_{k,\gamma,t,f}\propto e^{t\sum_i f(x_i)}dP_{k,\gamma}$ and
$|t|C_f<1-\gamma$, then
\begin{equation}
 \operatorname{Var}_{P_{k,\gamma,t,f}}
          \left(\sum_{i=1}^N f(x_i)\right)
 \le\frac{NA_f}{1-\gamma-|t|C_f}.
 \label{ac:tilted-variance}
\end{equation}
Every one-particle marginal of $P_{k,\gamma}$ is exactly $\sigma$.
\end{proposition}

\begin{proof}
Write $\Phi_0=\sum_i\phi_0(x_i)$ in a local frame of $-K_M$ and
write $s$ for the corresponding holomorphic representative of
$\Delta_k$.  The unnormalized density is the canonical volume of
the singular metric
\[
 \Phi_\gamma=(1-\gamma)\Phi_0+\frac\gamma k\log|s|^2.
\]
For $\varepsilon>0$ use the smooth metric
\begin{equation}
 \Phi_{\gamma,\varepsilon}
  =(1-\gamma)\Phi_0+\frac\gamma k
                 \log\bigl(|s|^2+\varepsilon e^{k\Phi_0}\bigr).
 \label{ac:regularization}
\end{equation}
Under a change of line-bundle frame the expression in parentheses
changes by the appropriate squared transition factor, so this formula
defines a global metric.  The logarithm is plurisubharmonic: it is
the logarithmic sum of exponentials of the plurisubharmonic functions
$\log|s|^2$ and $\log\varepsilon+k\Phi_0$.  It is smooth even at
$s=0$.  Hence
\[
 i\partial\bar\partial\Phi_{\gamma,\varepsilon}
                          \ge(1-\gamma)G_M.
\]
Its canonical density, relative to $\sigma^{\otimes N}$ up to an
irrelevant constant, is
$(\|\Delta_k\|^2+\varepsilon)^{-\gamma/k}$.
As $\varepsilon\downarrow0$, these densities increase to
$w_{k,\gamma}$, and their integrals increase to $Z_{k,\gamma}<\infty$.
The resulting probability measures therefore converge in total
variation.  Apply Lemma~\ref{ac:canonical-poincare} to the regularized
metrics and pass to the limit.  Both the integrands on the right and
the first two moments on the left are bounded for a fixed smooth $F$.
This proves \eqref{ac:variance-eq} without applying a smooth theorem
directly to a singular metric.

Multiplication of the density by $e^{t\sum_i f(x_i)}$ changes the
metric in \eqref{ac:regularization} to
$\Phi_{\gamma,\varepsilon}-t\sum_i f(x_i)$.  Its curvature is at
least $(1-\gamma-|t|C_f)G_M$.  In the product metric,
\[
 \left|\bar\partial\sum_i f(x_i)\right|_{G_M}^2
             =\sum_i|\bar\partial f(x_i)|_{G_X}^2\le NA_f.
\]
The same approximation proves \eqref{ac:tilted-variance}.

Finally, the diagonal action of $SU(D)$ preserves the product reference
measure and the squared determinant norm.  Indeed, it acts unitarily
on $H^0(X,L^k)$, so the evaluation determinant changes only by the
determinant of a unitary basis-change matrix.  Its modulus is unchanged.
Every one-particle marginal is therefore $SU(D)$-invariant, and the
unique invariant probability on $X$ is $\sigma$.
\end{proof}

\subsection{An exact projective Ward identity}

\begin{proposition}[Sharp covariance of the first moment]
\label{ac:ward}
Suppose $\gamma<1$ and $Z_{k,\gamma}<\infty$, using the same
density formula also for negative $\gamma$.
For traceless Hermitian $D\times D$ matrices $H,K$, let
$f_H(z)=\operatorname{tr}(HP_z)$ and $S_H=\sum_i f_H(z_i)$.
Then
\begin{equation}
 \operatorname{Cov}_{P_{k,\gamma}}(S_H,S_K)
 =\frac{N\operatorname{tr}(HK)}{D(D+1)(1-\gamma)}.
 \label{ac:ward-covariance}
\end{equation}
In the attractive range $0\le\gamma<1$, the constant
$1/(1-\gamma)$ in \eqref{ac:variance-eq} is therefore optimal.
For two distinct labels,
\begin{equation}
 \mathbb E_{P_{k,\gamma}}\operatorname{tr}(P_{z_1}P_{z_2})
 =\frac1D+
   \frac{(D-1)\gamma}{D(N-1)(1-\gamma)}.
 \label{ac:pair-overlap}
\end{equation}
Equivalently, with $M_1=N^{-1}\sum_iP_{z_i}$,
\begin{equation}
 \mathbb E_{P_{k,\gamma}}
      \left\|M_1-\frac{I}{D}\right\|_{\rm HS}^2
       =\frac{D-1}{DN(1-\gamma)}.
 \label{ac:first-moment-variance}
\end{equation}
Consequently finiteness of the partition function in this range
requires $\gamma<1-1/N$.
\end{proposition}

\begin{proof}
For $A\in SL(D,\mathbb C)$ the projective map
$g_A(z)=Az/\|Az\|$ has Jacobian $\|Az\|^{-2D}$ relative to
$\sigma$.  The determinant of the $SL(D,\mathbb C)$ representation
on $\operatorname{Sym}^{Dk}\mathbb C^D$ is one.  Multihomogeneity
of the evaluation determinant therefore gives
\[
 \|\Delta_k(g_Az_1,\ldots,g_Az_N)\|^2
 =\|\Delta_k(z_1,\ldots,z_N)\|^2
                  \prod_i\|Az_i\|^{-2Dk}.
\]
Change variables by $g_A$ in the finite partition function.  Taking
$A=e^{tH/2}$ gives the exact identity
\begin{equation}
 \mathbb E_{P_{k,\gamma}}
 \exp\!\left[-D(1-\gamma)
       \sum_i\log(z_i^*e^{tH}z_i)\right]=1.
 \label{ac:ward-identity}
\end{equation}
For $t$ in a compact interval the exponent and its first two
derivatives are bounded, so differentiation is justified by the
already assumed finiteness of $Z_{k,\gamma}$.
The first derivative vanishes, and the second gives
\[
 D(1-\gamma)\mathbb E S_H^2
 =\mathbb E\sum_i
       \left[\operatorname{tr}(H^2P_{z_i})-f_H(z_i)^2\right].
\]
Every one-particle marginal is Haar.  Its second-moment identities are
\[
 \mathbb E_\sigma\operatorname{tr}(H^2P)=\frac{\operatorname{tr}H^2}{D},
 \qquad
 \mathbb E_\sigma f_H^2=\frac{\operatorname{tr}H^2}{D(D+1)}.
\]
Substitution and polarization prove \eqref{ac:ward-covariance}.
Direct differentiation of $f_H$ in the invariant metric gives
\[
 |\bar\partial f_H|_{G_X}^2
 =\frac1D\left[\operatorname{tr}(H^2P)-f_H^2\right].
\]
In the attractive range, equality therefore holds in
\eqref{ac:variance-eq} for $F=S_H$ whenever $H\ne0$, proving
sharpness.  The Ward-identity calculation itself is unchanged for
negative $\gamma$, where the partition function is finite because
the density is bounded on the compact configuration space.

Permutation symmetry and the known one-particle covariance imply
\[
 \mathbb E f_H(z_1)f_K(z_2)
 =\frac{\gamma\operatorname{tr}(HK)}
 {(N-1)(1-\gamma)D(D+1)}.
\]
Sum over a Hilbert--Schmidt orthonormal basis of traceless Hermitian
matrices and use
$\operatorname{tr}(PQ)=D^{-1}+\sum_a
\operatorname{tr}(H_aP)\operatorname{tr}(H_aQ)$ to obtain
\eqref{ac:pair-overlap}.  Summing \eqref{ac:ward-covariance} in the
same basis and dividing by $N^2$ proves
\eqref{ac:first-moment-variance}.  For two distinct coordinate labels
under a measure absolutely continuous with respect to
$\sigma^{\otimes N}$, one has
$\operatorname{tr}(P_{z_1}P_{z_2})<1$ almost surely.
The right side of \eqref{ac:pair-overlap} must therefore be strictly
less than one, which is equivalent to $\gamma<1-1/N$.
\end{proof}

\subsection{An unconditional attractive range}

\begin{lemma}[Integrability by monomial degeneration]
\label{ac:integrability}
For every $k\ge1$ the partition function is finite whenever
$0\le\gamma<1/D$.  More generally, every nonzero section of
$\mathcal O(m,\ldots,m)$ on $(\mathbb P^q)^N$ has complex
singularity exponent at least $1/m$.
\end{lemma}

\begin{proof}
Expand a nonzero section in the multihomogeneous monomial basis.
Choose integer weights on the homogeneous coordinates, independently
in each factor, such that one monomial in its finite support uniquely
minimizes the total weight.  Pullback by the corresponding toric
one-parameter subgroup, followed by multiplication by the scalar
power of the parameter that sets this minimum weight to zero, gives
a holomorphic family whose limit is that single monomial.  For a
nonzero parameter the global complex singularity exponent is unchanged:
the transformation is an automorphism and the scalar is nonzero.

The divisor of the limiting monomial is a sum of the coordinate
hyperplanes, each with multiplicity at most $m$.  These hyperplanes
form a simple normal-crossing divisor: in each projective factor at
most $q$ of its $q+1$ coordinate hyperplanes meet.  In local coordinates
integrability of a monomial is the product of the conditions
$c\alpha_j<1$; its global exponent is therefore at least $1/m$.
Lower semicontinuity of complex singularity exponents
\cite{geo_demailly_kollar2001} shows that the original section has
exponent at least that of this limit.  Smooth positive metric and
volume factors do not affect local integrability.  Taking $m=Dk$
and $c=\gamma/k$ proves the first assertion.
\end{proof}

\begin{theorem}[Attractive empirical concentration]
\label{ac:concentration}
Suppose $0\le\gamma<1$, and suppose $Z_{k,\gamma}<\infty$ for the
values of $k$ under consideration.  For $f\in C^\infty(X,\mathbb R)$
and $|t|C_f<1-\gamma$,
\begin{equation}
 \log\mathbb E_{P_{k,\gamma}}
       \exp\!\left(t\sum_{i=1}^N[f(x_i)-\sigma(f)]\right)
 \le\frac{NA_ft^2}{2(1-\gamma-|t|C_f)}.
 \label{ac:mgf}
\end{equation}
In particular, if $A_f>0$, then for every $s>0$,
\begin{equation}
 P_{k,\gamma}\!\left(
    \left|\frac1N\sum_i f(x_i)-\sigma(f)\right|\ge s\right)
 \le 2\exp\!\left(
    -\frac{N(1-\gamma)s^2}{2(A_f+C_fs)}\right).
 \label{ac:tail}
\end{equation}
If $A_f=0$, then $f$ is constant and the deviation probability is
zero.  These conclusions are unconditional for $0\le\gamma<1/D$.
\end{theorem}

\begin{proof}
Let $S_f=\sum_i f(x_i)$ and
$\psi(t)=\log\mathbb E_{P_{k,\gamma}}e^{tS_f}$.
The random variable $S_f$ is bounded for fixed $N$, so differentiation
under the integral gives
$\psi''(t)=\operatorname{Var}_{P_{k,\gamma,t,f}}(S_f)$.
By Proposition~\ref{ac:variance}, $\psi'(0)=N\sigma(f)$, and
\eqref{ac:tilted-variance} bounds the second derivative on the segment
between $0$ and $t$ by $NA_f/(1-\gamma-|t|C_f)$.
Integrating twice proves \eqref{ac:mgf}.

For the upper tail in \eqref{ac:tail}, apply Chernoff's inequality
with
$t=(1-\gamma)s/(A_f+C_fs)$, which lies in the allowed interval.
Substitution in \eqref{ac:mgf} yields the exponent in
\eqref{ac:tail}.  Apply the same argument to $-f$ for the lower tail
and add the two bounds.  Lemma~\ref{ac:integrability} supplies the
unconditional coupling range.
\end{proof}

\begin{corollary}[Moment tubes and propagation of chaos]
\label{ac:tubes}
Fix real smooth functions $f_1,\ldots,f_r$ and set
$F=(f_1,\ldots,f_r)$,
$A=\max_j A_{f_j}$, and $C=\max_j C_{f_j}$.
Under the assumptions of Theorem~\ref{ac:concentration},
\begin{equation}
 P_{k,\gamma}\!\left(
 \left\|\frac1N\sum_iF(x_i)-\sigma(F)\right\|_2\ge s\right)
 \le 2r\exp\!\left(
 -\frac{N(1-\gamma)s^2}{2(rA+\sqrt r\,Cs)}\right)
 \label{ac:vector-tail}
\end{equation}
whenever $A>0$.  Thus every fixed-order projective moment tube of
width $N^{-a}$, $0<a<1/2$, has probability at least
$1-O(e^{-cN^{1-2a}})$, with constants independent of $k$ for fixed
$q,\gamma$, and moment order.

As $k\to\infty$, the empirical measure converges in probability,
in the weak topology, to $\sigma$.  For every fixed $\ell$, the
$\ell$-particle marginal converges weakly to $\sigma^{\otimes\ell}$.
\end{corollary}

\begin{proof}
If the norm in \eqref{ac:vector-tail} is at least $s$, at least one
coordinate has absolute value at least $s/\sqrt r$.  The union bound
and \eqref{ac:tail} give the displayed estimate.
For projective moments use the real coordinates
$f_j(z)=\operatorname{tr}(H_jP_z^{\otimes t})$ in a Hilbert--Schmidt
orthonormal basis of Hermitian operators on
$\operatorname{Sym}^t\mathbb C^D$, omitting the fixed trace coordinate.
These are a fixed finite family of smooth functions, and their
Euclidean norm is the Hilbert--Schmidt moment error.  Substitute
$s=N^{-a}$.

Smooth functions are uniformly dense in the continuous functions on
$X$, so the scalar bounds imply weak convergence of the empirical
measure in probability.  The measures $P_{k,\gamma}$ are permutation
invariant.  The difference between an expectation of a fixed product
$\prod_{j=1}^{\ell}f_j(x_j)$ and the corresponding expectation under
the $\ell$th power of the empirical measure is $O(\ell^2/N)$ for
bounded $f_j$: it is precisely the error from sampling indices with
replacement rather than without replacement.  Empirical convergence
and boundedness then give the product limit.  Density of finite sums
of products in $C(X^\ell)$ proves weak convergence of the marginals.
\end{proof}

\paragraph{Consequences for constrained partitions.}
Let $\Omega_N$ be one of the preceding shrinking moment tubes and
let $Q_N=\sigma^{\otimes N}(\,\cdot\mid\Omega_N)$.
Its attractive partition function satisfies the exact relation
\[
 \int e^{N\gamma E_k}\,dQ_N
 =Z_{k,\gamma}
    \frac{P_{k,\gamma}(\Omega_N)}
         {\sigma^{\otimes N}(\Omega_N)},
 \qquad
 E_k=-\frac1{kN}\log\|\Delta_k\|^2.
\]
Both probabilities tend to one at a stretched-exponential rate, by
Corollary~\ref{ac:tubes} for $\gamma$ and for $\gamma=0$.
Consequently conditioning on these tubes does not change any existing
speed-$N$ pressure limit in this coupling range.

\paragraph{Scope of the curvature argument.}
The curvature argument supplies finite-particle variance estimates,
exponential concentration of empirical observables, and propagation
of chaos for the actual attractive determinant law.  The logarithmic
energy control needed for pressure is provided by the compact-conditional
entropy method in Theorem~\ref{att:gibbs-energy-limit}.
Theorems~\ref{att:full-laplace} and \ref{att:design-laplace} establish
the full attractive Laplace principle on $\mathbb P^2$ for
$0\le\gamma<1/3$, and Corollary~\ref{att:polarized-extension}
extends the unrestricted statement to positive polarized manifolds
below their scalar section-integrability threshold.  The concentration
constants above degenerate at $\gamma=1$; extending the logarithmic
energy method toward that coupling is a separate question.



\section{Ordered logarithms and positive modular calculus}
\label{sec:positive-calculus}

The logarithm is where commutative and noncommutative expansions part company.
For a scalar gas it selects connected clusters. For an operator gas it must
also remember the order of mutually incompatible factors. Heaps retain that
order economically. Their coefficients explain both the cancellations that
make a logarithm local in a formal sense and the obstruction to estimating it
by absolute scalar cluster weights. A different representation, based on
positive random times, then gives an operator approximation with a useful
error certificate.

\subsection{Heaps, order polynomials, and the logarithm}

Fix a finite alphabet with a symmetric reflexive dependence relation. Two
distinct letters commute precisely when they are independent. A word modulo
these commutations is a \emph{heap}: its occurrences form a partially ordered
set, with the order generated by dependent letters appearing in word order.
This is the Cartier--Foata--Viennot description of partial commutation
\citep{CartierFoata1969,Viennot1986}. It records the information that would be
lost by replacing ordered products by commutative monomials.

Let $M=\sum_C(-1)^{|C|}C$ and $Z=\sum_CC$, where $C$ runs over cliques of
mutually independent letters, including the empty clique. Both are formal
series in the completed heap algebra. For a nonempty heap $H$, write
$\Omega_H(t)$ for the polynomial counting weakly order-preserving maps from
$H$ to $\{1,\ldots,t\}$.

\begin{proposition}[The order-polynomial logarithm]
\label{op:heap}
For every nonempty heap,
\begin{equation}
 [H](-\log M)=\Omega_H'(0),\qquad
 [H]\log Z=(-1)^{|H|+1}\Omega_H'(0).
 \label{op:heap-coefficient}
\end{equation}
Disconnected heaps have zero coefficient. For a natural labeling of a heap
with $n$ elements, the first coefficient also equals
\begin{equation}
 \Omega_H'(0)=\sum_{\pi\in\mathcal L(H)}
 \frac{(-1)^{\operatorname{des}(\pi)}}
 {n\binom{n-1}{\operatorname{des}(\pi)}},
 \label{op:eulerian}
\end{equation}
where $\mathcal L(H)$ denotes its linear extensions.
\end{proposition}

\begin{proof}
For a positive integer $t$, the coefficient of $H$ in $M^t$ is
$(-1)^n\overline\Omega_H(t)$: choosing the clique in each factor is a strict
order-preserving map. Polynomial reciprocity gives
$\overline\Omega_H(t)=(-1)^n\Omega_H(-t)$. Differentiating at $t=0$ yields
the first identity; substituting the negative of each letter gives the
second. A disjoint union has a product of order polynomials, each vanishing
at zero, so its derivative vanishes. Finally, the standard linear-extension
expansion
\[
 \Omega_H(t)=\sum_{\pi\in\mathcal L(H)}
 \binom{t+n-1-\operatorname{des}(\pi)}n
\]
gives \eqref{op:eulerian} by differentiation.
\end{proof}

There are two useful interpretations of the coefficient. If
$\mathcal O(H)\subset[0,1]^H$ is Stanley's order polytope, then
$\Omega_H(t)=L_{\mathcal O(H)}(t-1)$, where $L$ is its Ehrhart polynomial
\citep{Stanley1986}. Thus the logarithm evaluates a shifted lattice-point
polynomial at a negative argument. Formula \eqref{op:eulerian}, on the other
hand, is the Eulerian-idempotent expression familiar in Lie-theoretic
logarithms and their pre-Lie refinements \citep{BandieraSchaetz2017}. These
are structural identifications of the coefficient, rather than positivity
statements about it.

For example, let $H_m$ consist of $m$ independent leaves below one root.
Then $\Omega_{H_m}(t)=\sum_{j=1}^tj^m$ and
\begin{equation}
 \Omega_{H_m}'(0)=B_m(1),
 \label{op:bernoulli}
\end{equation}
with the Bernoulli-polynomial convention $B_1(1)=1/2$. The coefficients
change sign and vanish for odd $m>1$. A connected heap need not make a
positive contribution to an ordered logarithm.

\fig{fig07_ordered_heaps}{Ordered logarithms retain signed combinatorics.
(a) A star heap consists of independent leaves below a common root; its order
polynomial is a sum of powers. (b) The derivatives at zero are Bernoulli
values; positive and negative nonzero values are distinguished by color,
and gray crosses indicate zeros. The vertical axis records the logarithm of
the magnitude only for nonzero coefficients. (c) The root-linear logarithmic
response of the positive operator gas in the text has Taylor radius
$2\sin(\pi/m)$ for $m\ge7$. It eventually lies below the fixed leaf activity
$1/16$, despite a uniform scalar cluster criterion.}{fig:heaps}

\paragraph{A positive gas for which absolute ordered expansion fails.}
On $m$ qubits put $P_i=|1\rangle\langle1|_i$ and
$X=|0^m\rangle\langle1^m|+|1^m\rangle\langle0^m|$. Consider
\[
 Z_m=\prod_{i=1}^m(I+a_iP_i)+bX,
 \qquad a_i=a=\frac1{16},\qquad b=\frac1{8\,4^m}.
\]
Its least eigenvalue is positive: the diagonal product is at least $I$ and
$\|bX\|=b<1$. Treating each singleton and the full support as polymers,
the Koteck\'y--Preiss condition with size cost $|S|\log2$ remains valid even
after the additional support tilt $2^{|S|}$. Indeed, the incompatible sum at
a leaf is $4a+4^mb=3/8\le\log2$, and at the root is
$4ma+4^mb=m/4+1/8\le m\log2$.

Nevertheless the $(0^m,1^m)$ matrix element of the derivative of $\log Z_m$
at $b=0$, with all leaf activities equal, is
\begin{equation}
 f_m(a)=\frac{m\log(1+a)}{(1+a)^m-1}.
 \label{op:root-response}
\end{equation}
This is the divided difference of the scalar logarithm between the diagonal
eigenvalues $1$ and $(1+a)^m$. For $m\ge7$ its nearest nonremovable poles
are at $a=e^{\pm2\pi i/m}-1$, with modulus $2\sin(\pi/m)<1$.
At $m=128$ this is smaller than $1/16$. Thus even the root-linear part of
the absolute multivariate Taylor expansion cannot converge at the displayed
positive activities: absolute convergence would imply convergence on the
diagonal. The scalar criterion controls its own commutative expansion; the
ordered logarithm requires a cancellation-aware representation.

\subsection{The normalization of a quantum response}
\label{op:normalization}

Let $\rho>0$ be a density matrix, $K=\log\rho$, and
$\dot\rho=\dot\rho^*$ a trace-zero tangent. Two natural Hermitian inputs are
\[
 A=\rho^{-1/2}\dot\rho\rho^{-1/2},\qquad
 H=\frac12(\dot\rho\rho^{-1}+\rho^{-1}\dot\rho).
\]
They are related but not interchangeable. Write $D=\ad_K$. Divided
differences in an eigenbasis of $\rho$ give
\begin{equation}
 D\log_\rho(\dot\rho)
 =\frac{D/2}{\sinh(D/2)}A
 =\frac D{\sinh D}H.
 \label{op:log-response}
\end{equation}
For instance, the first multiplier at eigenvalues $p_i,p_j$ is
$\sqrt{p_ip_j}(\log p_i-\log p_j)/(p_i-p_j)$, with diagonal value one.
The half-argument in the first expression is forced by the relative
normalization. It is relevant when translating between modular response,
quantum scores, and numerical implementations.

The functions in \eqref{op:log-response} are characteristic functions.
With the Fourier convention $f(x)=\int e^{itx}p(t)\,dt$,
\begin{equation}
 \frac{x/2}{\sinh(x/2)}
 =\int_{\R}e^{itx}\frac{\pi}{2\cosh^2(\pi t)}\,dt,
 \qquad
 \frac x{\sinh x}
 =\int_{\R}e^{itx}\frac{\pi}{4\cosh^2(\pi t/2)}\,dt.
 \label{op:fourier-log}
\end{equation}
Consequently each multiplier is a positive average of unitary
conjugations $A\mapsto e^{itK}Ae^{-itK}$. This replaces ordered Taylor
coefficients by an exact completely positive integral. The scalar infinite
divisibility has a classical probabilistic theory \citep{PitmanYor2003};
positive-definite modular multipliers also belong to the theory of
quantum information metrics \citep{HiaiKosakiPetzRuskai2013}.

\subsection{Gamma clocks and errors measured on the observable}

For a Hermitian $G$, let $L=\ad_G^2$. The heat map $e^{-sL}$ is the average
of $A\mapsto e^{itG}Ae^{-itG}$ for a centered Gaussian time of variance
$2s$. It is unital, trace preserving, completely positive, and an operator
norm contraction. We use this representation even though $L$ is a
superoperator rather than a positive map.

\begin{theorem}[Positive hyperbolic clock and local error]
\label{op:hyperbolic-clock}
For $\tau\ge0$,
\begin{equation}
 T_\tau=\left(\frac{\ad_G}{\sinh\ad_G}\right)^\tau
 =\mathbb E e^{-S_\tau L},\qquad
 S_\tau=\sum_{j\ge1}Y_j,\quad
 Y_j\sim\operatorname{Gamma}\left(\tau,\text{rate }\pi^2j^2\right)
 \label{op:gamma-clock}
\end{equation}
with independent summands. It is a random-unitary semigroup. Put
$S_{\tau,R}=\sum_{j\le R}Y_j$ and
$a_R=\mathbb E(S_\tau-S_{\tau,R})$. The finite and mean-corrected channels
obey
\begin{align}
 \|\mathbb Ee^{-S_{\tau,R}L}A-T_\tau A\|
 &\le\frac{\tau}{\pi^2R}\|\ad_G^2A\|,\label{op:clock-first}\\
 \|\mathbb Ee^{-(S_{\tau,R}+a_R)L}A-T_\tau A\|
 &\le\frac{\tau}{6\pi^4R^3}\|\ad_G^4A\|.
 \label{op:clock-mean}
\end{align}
\end{theorem}

\begin{proof}
Euler's product for $\sinh$ and the Laplace transform of a Gamma variable
give the identity. The sum exists almost surely because its mean is
$\tau\sum_j(\pi j)^{-2}=\tau/6$. For the first bound, integrate the
generator along the omitted nonnegative time and use contraction. For the
second, write the omitted time as $U$. Taylor's theorem with integral
remainder around $\mathbb EU$ is valid on every segment between these
nonnegative times. Its linear term has mean zero, and its remainder is at
most $\tfrac12\operatorname{Var}(U)\|L^2A\|$. The independent Gamma
variances give
\[
 \operatorname{Var}(U)=\tau\sum_{j>R}(\pi j)^{-4}
 \le\frac{\tau}{3\pi^4R^3}.
\]
The retained heat channel is a contraction, which completes both estimates.
\end{proof}

These bounds involve the supplied observable and its actual commutators.
Replacing them by $\|G\|$ would discard useful spatial information and
could introduce the total system size. Conversely, the relative input $A$
in \eqref{op:log-response} need not have the same support as $\dot\rho$.
The theorem does not assert spatial locality of modular evolution.

\fig{fig08_modular_clocks}{Positive modular filters and finite random-time
approximations. (a) The normalized power-response multiplier
$\sinh(rx)/(r\sinh x)$ approaches $x/\sinh x$ as $r\downarrow0$.
(b) Its positive Fourier density is
$\sin(\pi r)/[2r(\cosh(\pi t)+\cos(\pi r))]$; each curve integrates to one.
(c) At scalar frequency $4$ and semigroup time $1$, retaining $R$ Euler
factors has an $R^{-1}$ error bound, while replacing the omitted clock by
its mean gives an $R^{-3}$ bound. Solid curves are exact scalar errors and
dashed curves are the stated analytic bounds; this panel tests the scalar
multiplier, not a many-body interaction norm.}{fig:clocks}

\subsection{Bessel functions connect the clock to quantum information geometry}
\label{op:bessel}

The hyperbolic example belongs to a larger family. For $\nu>-1$, define
\[
 \mathcal I_\nu(x)=\Gamma(\nu+1)(2/x)^\nu I_\nu(x)
 ={}_0F_1(\nu+1;x^2/4),\qquad \phi_\nu(x)=\mathcal I_\nu(x)^{-1}.
\]
The value at zero is one. If $j_{\nu,k}$ are the positive zeros of $J_\nu$,
then the Bessel product gives
\[
 \phi_\nu(x)=\prod_{k\ge1}(1+\lambda_kx^2)^{-1},\qquad
 \lambda_k=j_{\nu,k}^{-2},\qquad
 \sum_k\lambda_k=\frac1{4(\nu+1)}.
\]
Thus $\phi_\nu(\ad_G)^\tau$ is a random-unitary semigroup with clock
$\sum_k\operatorname{Gamma}(\tau,\text{scale }\lambda_k)$. At $\tau=1$
this is one half of the exit time from radius one of a Bessel process of
dimension $2\nu+2$, started at zero. The reciprocal-product and probability
connections are classical; finite moment-optimal Gamma approximations are
developed by Clark \citep{Clark2026}. Here their operator use is governed by
the commutator estimates below.

To connect with a quantum statistical model, set $G=K/2$ and
\begin{equation}
 m_\nu(t)=\sqrt t\,\mathcal I_\nu(\tfrac12\log t).
 \label{op:bessel-mean}
\end{equation}
A symmetric normalized operator-monotone function $m$ defines the Petz
information metric through
\[
 \mathcal J_{m,\rho}=R_\rho m(L_\rho R_\rho^{-1}),\qquad
 g_{m,\rho}(\dot\rho_i,\dot\rho_j)
 =\Re\tr\bigl(\dot\rho_i\mathcal J_{m,\rho}^{-1}(\dot\rho_j)\bigr),
\]
where $L_\rho$ and $R_\rho$ mean left and right multiplication
\citep{Petz1996,HiaiKosakiPetzRuskai2013}. Operator monotonicity ensures
that information cannot increase after a quantum channel.

\begin{proposition}[Exact metric range]
\label{op:metric-range}
The function \eqref{op:bessel-mean} defines a symmetric normalized
operator-monotone mean exactly for $\nu\ge-1/2$. In that range,
\begin{equation}
 \mathcal J_{m_\nu,\rho}^{-1}(\dot\rho)
 =\phi_\nu(\ad_K/2)A.
 \label{op:bessel-score}
\end{equation}
The endpoint is the symmetric logarithmic derivative (SLD) metric;
$\nu=1/2$ is the Bogoliubov--Kubo--Mori (BKM) metric; the limit as
$\nu\to\infty$ is the geometric-mean metric.
\end{proposition}

\begin{proof}
For $\nu>-1/2$ the Bessel integral representation becomes
\[
 m_\nu(t)=\mathbb E t^U,\qquad
 U\sim\operatorname{Beta}(\nu+\tfrac12,\nu+\tfrac12).
\]
It is a positive mixture of the operator-monotone functions $t^u$,
$0\le u\le1$; the symmetric beta law gives $tm_\nu(1/t)=m_\nu(t)$.
The endpoint law puts mass $1/2$ at each of $0,1$. Accordingly,
\[
 m_{-1/2}(t)=\frac{1+t}2,\qquad
 m_{1/2}(t)=\frac{t-1}{\log t},\qquad
 m_\infty(t)=\sqrt t.
\]
For $-1<\nu<-1/2$ direct differentiation gives
$m_\nu''(1)=-(2\nu+1)/[8(\nu+1)]>0$, contradicting the scalar concavity
of an operator-monotone function. Finally, comparing entries in an
eigenbasis of $\rho$ proves \eqref{op:bessel-score}.
\end{proof}

This has an operational distinction worth preserving. SLD Fisher
information describes local parameter estimation; the BKM metric is the
affine Hessian of $\tr(\rho\log\rho)$ and the quadratic term of quantum
relative entropy. The remaining means interpolate between information
geometries. Complete positivity of their normalized score filters holds
even below $\nu=-1/2$, but monotonicity of the corresponding metric does
not. For example, $\nu=-3/4$, $\rho=\diag(4/5,1/5)$ and
$\dot\rho=X/10$ give information approximately $0.033246$, whereas
measurement in the $X$ basis gives classical Fisher information $0.04$.
This numerical illustration is supplementary; the exact second-derivative
test already proves the boundary of the metric family.

\subsection{Positive quadrature produces information-matrix brackets}
\label{op:quadrature}

An infinite clock is not needed for computation. Let
$\mu=\sum_k\lambda_k\delta_{\lambda_k}$, a finite positive measure on a
bounded interval, and write
\[
 g(L)=\int F_\lambda(L)\,d\mu(\lambda),\qquad
 F_\lambda(L)=\frac{\log(I+\lambda L)}\lambda
 =\int_0^1 L(I+u\lambda L)^{-1}\,du.
\]
At $\lambda=0$ the continuous value is $L$. Approximating this generator
by positive quadrature retains the entire channel semigroup. It also makes
two complementary error directions available.

\begin{theorem}[Finite positive quadrature with local certificates]
\label{op:quadrature-theorem}
Let $\mu$ be a finite positive measure of bounded support in $[0,\infty)$,
with the nondegenerate moment matrices needed for the following rules.
For its $r$-node Gaussian rule $(c_i,w_i)$, put
\[
 T_r^G(\tau)=\prod_{i=1}^r(I+c_iL)^{-\tau w_i/c_i},\qquad
 h_r=\int p_r(\lambda)^2\,d\mu(\lambda),
\]
where $p_r$ is the monic orthogonal polynomial. Then
\begin{equation}
 \|(T_r^G(\tau)-e^{-\tau g(L)})A\|
 \le\frac{\tau h_r}{2r+1}\|\ad_G^{4r+2}A\|.
 \label{op:gauss-error}
\end{equation}
For the Gaussian rule $(c_i,v_i)$ of $\lambda\,d\mu$, set
$w_i=v_i/c_i$ and $w_0=\mu([0,\infty))-\sum_iw_i\ge0$. The Radau rule
\[
 T_r^R(\tau)=e^{-\tau w_0L}
 \prod_{i=1}^r(I+c_iL)^{-\tau w_i/c_i}
\]
satisfies, with $\widetilde p_r$ monic orthogonal for $\lambda\,d\mu$,
\begin{equation}
 \|(T_r^R(\tau)-e^{-\tau g(L)})A\|
 \le\frac{\tau\widetilde h_r}{2r+2}\|\ad_G^{4r+4}A\|,
 \quad
 \widetilde h_r=\int\lambda\widetilde p_r(\lambda)^2\,d\mu.
 \label{op:radau-error}
\end{equation}
All maps displayed are random-unitary channels for $\tau\ge0$. Their
scalar multipliers obey $T_r^R(\tau;s)\le e^{-\tau g(s)}\le T_r^G(\tau;s)$
for $s\ge0$.
\end{theorem}

\begin{proof}
Differentiating the resolvent representation gives
\[
 \partial_\lambda^kF_\lambda(L)
 =(-1)^kk!\int_0^1u^k L^{k+1}(I+u\lambda L)^{-k-1}\,du.
\]
Resolvent powers are positive averages of heat channels and commute with
$L$. Thus their action on $A$ has norm at most
$k!\|L^{k+1}A\|/(k+1)$. The Gaussian Hermite-interpolation remainder has
factor $p_r^2$ and derivative order $2r$; the Radau remainder has factor
$\lambda\widetilde p_r^2$ and derivative order $2r+1$.
Integration therefore bounds the generator errors by the right sides of
\eqref{op:gauss-error}--\eqref{op:radau-error} without $\tau$.
For the operator-valued remainder this statement follows from the positive
Peano-kernel integral for Hermite interpolation, followed by the norm bound;
equivalently, scalarize by norm-one functionals and apply Hahn--Banach.
The signed scalar derivative formula gives the stated multiplier directions.
Commuting-semigroup Duhamel and contraction multiply the generator bounds
by $\tau$. Gaussian and Radau weights are positive, so each approximant
is a finite Gamma clock, with an additional deterministic heat time for
Radau. Degenerate finite-support cases terminate exactly.
An atom at zero contributes its continuous heat factor
$\exp(-\tau w_0L)$ rather than a Gamma factor with zero scale.
\end{proof}

For Bessel clocks the Rayleigh sums satisfy
\begin{equation}
 s_1=\frac1{4(\nu+1)},\qquad
 (\nu+n)s_n=\sum_{i=1}^{n-1}s_i s_{n-i},\quad n\ge2.
 \label{op:rayleigh}
\end{equation}
This follows by inserting the logarithmic derivative of the Bessel product
into its differential equation. Rational $\nu$ gives exact rational moments
and hence algebraically specified quadrature nodes and weights. For the
hyperbolic clock $\nu=1/2$, one Gaussian node already gives
\begin{equation}
 T_1^G(\tau)=(I+L/15)^{-5\tau/2},\qquad
 \|\mathrm{error}(A)\|\le\frac{\tau}{9450}\|\ad_G^6A\|.
 \label{op:one-clock}
\end{equation}
The supplemental implementation isolates the quadrature nodes by rational
Sturm intervals and encloses the weights independently of the floating
eigensolver. A separate rounding budget is added before a numerical rule
is called certified.

For a statistical model, let
$\langle A,B\rangle_{\rho,1/2}=\Re\tr(\rho^{1/2}A^*\rho^{1/2}B)$.
Since $\ad_K^2$ is self-adjoint and nonnegative for this inner product,
the scalar multiplier inequalities become finite information-matrix bounds:
\begin{equation}
 \mathsf G_r^R\preceq\mathsf G_\nu\preceq\mathsf G_r^G,
 \qquad
 (\mathsf G_r^*)_{ij}=
 \langle A_i,T_r^*(1)A_j\rangle_{\rho,1/2}.
 \label{op:information-bracket}
\end{equation}
This is Loewner order on matrices of parameter information. It is not
positive-map order: the difference of distinct unital maps cannot be
positive, since a positive map vanishing at the identity vanishes everywhere.
Nor is each finite approximation necessarily a monotone metric for all
states. Equation \eqref{op:information-bracket} is the fixed-state promise
needed for certified statistical calculations.

\fig{fig12_information_clocks}{Bessel clocks connect positive channels to
quantum information metrics. (a) The symmetric beta laws yield the exact
operator means $m_\nu(t)=\mathbb E[t^U]$; the displayed densities have
$\nu=0,1/2,2$. (b) The relative-score filters include the symmetric
logarithmic derivative (SLD), the Bogoliubov--Kubo--Mori (BKM) score, and
the limit associated with the geometric mean. (c) One positive Gaussian
quadrature node gives the upper scalar multiplier and one Radau node plus
heat drift gives the lower multiplier, enclosing $x/\sinh x$. The shaded
interval gives fixed-state information-matrix bounds through functional
calculus; it does not express an order between completely positive maps.
All curves are evaluated from the stated formulas.}{fig:information-clocks}

\paragraph{A positive empirical output.}
Here $G=K/2$, so modular unitary conjugation is an isometry for
$\langle\cdot,\cdot\rangle_{\rho,1/2}$.
Let $B_i=T_r^*(1/2)A_i$. The desired information matrix is the Gram matrix
of the $B_i$. Averaging $M$ independent half-time unitary features first,
then taking their Gram matrix, gives a positive semidefinite estimator.
Within each sampled tuple the same unitary is applied to every parameter
index $i$; different tuples use independent draws.
If $S=\sum_i\|A_i\|_{\rho,1/2}^2$, its operator error is at most
$2\sqrt S\,e_M+e_M^2$ with probability at least $1-\delta$, where
\[
 e_M=\sqrt{S/M}\bigl(1+\sqrt{2\log(1/\delta)}\bigr).
\]
Indeed, the expected squared feature error is at most $S/M$, and bounded
differences apply to its Hilbert norm. The expectation of the Gram estimator
is $(1-1/M)\mathsf G_r+\mathsf G_0/M$, where
$(\mathsf G_0)_{ij}=\langle A_i,A_j\rangle_{\rho,1/2}$. Its bias is explicit. This protocol
assumes access to features or their inner products; it does not replace the
measurement cost of obtaining them on quantum hardware.

\subsection{From Lee--Yang zeros to reusable channel constructions}

The preceding construction depends on a positive zero measure, not on the
Bessel differential equation. Suppose an even real entire function has
the normalized Laguerre--P\'olya factorization
\[
 F(z)=e^{-az^2}\prod_k(1-z^2/r_k^2),\qquad
 a\ge0,\quad r_k>0,\quad\sum_kr_k^{-2}<\infty.
\]
Then
\[
 \frac1{F(ix)}=e^{-ax^2}\prod_k(1+x^2/r_k^2)^{-1}
\]
is an infinitely divisible characteristic function. Substitution of
$\ad_G$ gives a random-unitary semigroup with heat drift $a$ and Gamma
scales $r_k^{-2}$. The same quadrature theorem compiles it. When $F$ is
itself a characteristic function, $(F,1/F(i\,\cdot))$ is a van Dantzig
pair; the relation to Lee--Yang functions is developed in
\citet{KonstantopoulosPatieSarkar2024}.

This supplies a precise route from a zero-location theorem to an operator
algorithm: real zeros determine positive scales, the scales determine a
clock, and quadrature yields a short positive approximation with a local
commutator error. It invokes no unproved assertion about the zeros of a
particular function. The zero-location hypothesis is a named input to the
construction and can be certified independently in finite examples.

% Self-contained journal-style fragments. Labels and citation keys use q:.
% Required packages: amsmath, amssymb, amsthm, mathtools, hyperref.
% Required theorem environments: theorem, lemma, proposition, corollary, remark.

\section{Completely positive transport and tensor-network gaps}
\label{q:sec-networks}

A transfer operator is usually used to estimate the decay of correlations
along a one-dimensional tensor network.  Here it has a second role: it
controls the angle between two local ground spaces meeting at a vertex
of an arbitrary graph.  The useful comparison is an inequality between
completely positive maps.  Such an inequality remains valid when a wire
is attached to an arbitrarily entangled positive core.  Moreover, two
different edge constraints can interact through that core only when both
of their wire maps depart from their rank-one limits.  This gives a
quadratic, rather than linear, error bound.

The use of transfer operators and ground-space angles for decorated
AKLT models originates in \cite{q:ARLLNY2020}; the extension to bounded-degree
graphs with a sufficient decoration linear in the degree is
\cite{LuciaYoung2023}.  The argument below gives a degree-independent
two-port angle estimate, leading to logarithmic sufficient decoration.
The final conversion from a blocked projector parent to a specified
microscopic Hamiltonian is stated separately, with its normalization
and finite-patch constants explicit.

\subsection{Finite local data}

Let $\mathcal G=(V,E)$ be a finite simple graph without isolated
vertices.  Open regions may also have dangling boundary ports; these
are included in every product over ports below, but carry no internal
edge constraint.  Every incidence $(v,e)$ carries a finite-dimensional exposed
virtual space.  Every edge carries a fixed nonzero vector of full Schmidt
rank on its two exposed spaces.  At each vertex choose an arbitrary
nonzero core map $A_v$, and attach one matrix-product half-wire to
each of its input legs.  Write $T_v$ for the resulting map from the
exposed virtual spaces to the physical block.  The block consists of
the core and its incident half-wires.

Put $C_v=A_v^*A_v$.  Contraction of the half-wires gives the exact identity
\begin{equation}
 T_v^*T_v=\left(\bigotimes_{e\ni v}\Phi_{v,e}\right)(C_v),
 \label{q:gram-channel}
\end{equation}
where each $\Phi_{v,e}$ is a completely positive (CP) map.  One may fix
the convention directly: if $W_{p,a,b}$ is the amplitude of a half-wire
with physical index $p$, core index $a$, and exposed index $b$, its
Kraus matrices are $(K_p)_{b,a}=\overline{W_{p,a,b}}$.  Thus
\[
 (T^*T)_{b',b}
 =\sum_{p,a',a}\overline{W_{p,a',b'}}(C_v)_{a',a}W_{p,a,b}
 =\left(\sum_p K_pC_vK_p^*\right)_{b',b}.
\]
In particular, no partial transpose of an entangled core is involved.

Assume that, for every incidence, there is a rank-one CP map
\[
 \Psi_{v,e}(X)=\operatorname{tr}(\tau_{v,e}X)R_{v,e},
 \qquad \tau_{v,e}>0,\quad R_{v,e}>0,
\]
and a number $0\leq\delta_{v,e}<1$ such that
\begin{equation}
 (1-\delta_{v,e})\Psi_{v,e}
 \ \leq_{\mathrm{CP}}\ \Phi_{v,e}
 \ \leq_{\mathrm{CP}}\ (1+\delta_{v,e})\Psi_{v,e}.
 \label{q:cp-sandwich}
\end{equation}
Here $\Phi\leq_{\mathrm{CP}}\widetilde\Phi$ means that
$\widetilde\Phi-\Phi$ is completely positive.  This is equivalent to a
positive-semidefinite comparison of their Choi matrices
\cite{q:Choi1975}.  The reference
map need not preserve the trace.  A sandwich
$a\Psi\leq_{\mathrm{CP}}\Phi\leq_{\mathrm{CP}}b\Psi$, with $0<a\leq b$,
is centered by replacing $\Psi$ with $(a+b)\Psi/2$ and setting
$\delta=(b-a)/(b+a)$.

Tensor products preserve CP order.  Consequently
\begin{equation}
 T_v^*T_v\geq
 \left[\prod_{e\ni v}(1-\delta_{v,e})\right]
 \operatorname{tr}\!\left[
   \left(\bigotimes_{e\ni v}\tau_{v,e}\right)C_v\right]
 \bigotimes_{e\ni v}R_{v,e}>0.
 \label{q:injectivity}
\end{equation}
Thus every block map is injective, even when the original core has rank
one.  Their tensor product is injective as well: it sends the nonzero
product of internal edge vectors and any nonzero boundary vector to a
nonzero physical state.  This proves frustration-freeness without an
assumption that a general tensor contraction avoids cancellation.
All estimates below are independent of the condition number of $C_v$.

For $e=uv$, contract the fixed edge vector between $T_u$ and $T_v$,
leaving their other exposed legs free, and denote the resulting regional
image by $\mathcal R_e$.  Let $G_e$ be its orthogonal projector in the
two physical blocks, extended by the identity elsewhere.  Define the
norm-one-projector parent
\begin{equation}
 H_{\mathrm{block}}=\sum_{e\in E}(I-G_e).
 \label{q:block-parent}
\end{equation}
Its ground space retains all outer virtual degrees of freedom when the
graph describes an open region.

\begin{theorem}[Two-port CP gap criterion]
\label{q:thm-wire-gap}
Under \eqref{q:cp-sandwich}, set
\begin{equation}
 \mathcal B(x,y)=\frac{2xy(1+xy)}{(1-x)^2(1-y)^2}.
 \label{q:Bxy}
\end{equation}
Let $C$ be the symmetric matrix indexed by edges, with zero diagonal,
with $C_{ef}=0$ for disjoint edges, and with
$C_{ef}=\mathcal B(\delta_{v,e},\delta_{v,f})$ when $e,f$ meet at $v$.
If its spectral radius satisfies $\rho(C)<1$, then
\begin{equation}
 H_{\mathrm{block}}\geq
 [1-\rho(C)](I-G_{\mathrm{all}}),
 \label{q:block-gap}
\end{equation}
where $G_{\mathrm{all}}$ projects onto its full ground space.
The same statement holds uniformly for a family of finite graphs if
their matrices $C$ have a common spectral-radius upper bound below one.
\end{theorem}

The simpler sufficient condition is the row bound

\fig{fig09_quantum_networks}{Completely positive wire estimates give
finite gap certificates. (a) A schematic network has arbitrary nonzero
cores and mixing half-wires; four sites per original edge are drawn only
to illustrate the construction. (b) The exact two-port estimate
$\mathcal B(\delta_i,\delta_j)$ is quadratic at the origin, reflecting
cancellation when either wire is exactly rank one. (c) Substitution of
$\delta_\ell=2/(3^\ell+(-1)^\ell)$ into the degree bound gives
$1-2(\Delta-1)\mathcal B(\delta_\ell,\delta_\ell)$ for a norm-one
blocked parent. Negative lower bounds are plotted as zero, meaning that
this criterion provides no positive certificate there. A microscopic
AKLT gap additionally includes the finite-patch gap and coverage factors
proved in the text.}{fig:networks}
\begin{equation}
 q=\max_{e=uv}\left[
 \sum_{\substack{f\ni u\\f\ne e}}
       \mathcal B(\delta_{u,e},\delta_{u,f})+
 \sum_{\substack{f\ni v\\f\ne e}}
       \mathcal B(\delta_{v,e},\delta_{v,f})\right]<1,
 \label{q:row-criterion}
\end{equation}
which gives gap at least $1-q$.  Rational positive-vector certificates
$Cw\leq rw$, $w>0$, $r<1$, can improve the row bound while avoiding any
floating-point eigenvalue calculation.  Finite edge types suffice when
their incidence counts give a uniform finite comparison matrix.

\subsection{The quadratic cancellation}

For subspaces $M,N$, their Friedrichs cosine is the largest overlap of
unit vectors in $M\ominus(M\cap N)$ and $N\ominus(M\cap N)$.
Removing the entire intersection is essential for degenerate local
ground spaces.

\begin{lemma}[A quotient-metric estimate]
\label{q:lem-schur}
Let $p,q$ be commuting orthogonal projections.  Suppose
\[
 G=I+A+B+D>0,\qquad [A,q]=[B,p]=0,
\]
\[
 \|A\|\leq x,\quad \|B\|\leq y,\quad \|D\|\leq xy,
 \qquad G\geq mI,\quad m=(1-x)(1-y)>0.
\]
The Friedrichs cosine between $G^{1/2}\operatorname{ran}p$ and
$G^{1/2}\operatorname{ran}q$ is at most $\mathcal B(x,y)$.
\end{lemma}

\begin{proof}
Put $W=\operatorname{ran}(pq)$,
$U=\operatorname{ran}(p(I-q))$, and
$V=\operatorname{ran}((I-p)q)$.  They are mutually orthogonal.  The
quotient by the common image of $W$ is represented by minimizing over
the $W$ coordinate in the $G$-metric.  Its $U,V$ cross block is therefore
\[
 S_{UV}=G_{UV}-G_{UW}G_{WW}^{-1}G_{WV}.
\]
Only $W$ is eliminated: the sector
$\operatorname{ran}((I-p)(I-q))$ is not part of the intersection.
The commutation assumptions imply
\[
 G_{UV}=D_{UV},\qquad
 G_{UW}=B_{UW}+D_{UW},\qquad
 G_{WV}=A_{WV}+D_{WV}.
\]
Hence
\[
 \|S_{UV}\|\leq xy+\frac{(y+xy)(x+xy)}m.
\]
For $u\in U$, the diagonal quotient metric obeys
\[
 \inf_{w\in W}\langle u+w,G(u+w)\rangle
 \geq m\inf_{w\in W}(\|u\|^2+\|w\|^2)=m\|u\|^2,
\]
and similarly on $V$.  The cosine equals
$\|S_{UU}^{-1/2}S_{UV}S_{VV}^{-1/2}\|$ and is at most
\[
 \frac{xy}{m}+
 \frac{xy(1+x)(1+y)}{m^2}
 =\frac{2xy(1+xy)}{(1-x)^2(1-y)^2}.
\]
\end{proof}

\begin{proof}[Proof of Theorem~\ref{q:thm-wire-gap}]
First restrict each block to its image under $T_v$ and remove its
isometric polar factor.  This identifies the network with independent
edge vectors acted on by invertible positive block filters.  An invertible
map preserves intersections of subspaces; hence the common kernel in
these block images is exactly the image of the product of fixed internal
edge vectors, with all dangling boundary factors arbitrary.  Consider
two edges meeting at $v$, and label their ports $i,j$.  Keep the actual
attachment map $\mathcal R$ on every unselected port.  With
$D_i=\Phi_i-\Psi_i$ and $D_j=\Phi_j-\Psi_j$, take the reference metric
\[
 G_0=(\Psi_i\otimes\Psi_j\otimes\mathcal R)(C_v).
\]
It is faithful by \eqref{q:injectivity}, and has product form
$R_i\otimes R_j\otimes R_{\mathrm{rest}}$.  Whitening therefore
acts separately on the selected ports and on the remaining ones.
The filters at the two outside endpoints can be absorbed into two
reference subspaces supported on disjoint sets of factors.  Their
orthogonal projectors $p,q$ commute.  The factor acting on the retained
ports preserves both subspaces.

After whitening the exact metric has the expansion
\[
 G=G_0^{-1/2}(\Phi_i\otimes\Phi_j\otimes\mathcal R)(C_v)G_0^{-1/2}
   =I+A+B+D.
\]
Here $A$ uses $D_i\otimes\Psi_j$, $B$ uses
$\Psi_i\otimes D_j$, and $D$ uses $D_i\otimes D_j$.
The first two terms have identity on the opposite selected port, so
$[A,q]=[B,p]=0$.  The single-port sandwiches imply
$\|A\|\leq x$, $\|B\|\leq y$, for $x=\delta_i$, $y=\delta_j$.
For the mixed term use the exact CP inequalities
\[
 -xy\,\Psi_i\otimes\Psi_j
 \leq_{\mathrm{CP}}D_i\otimes D_j
 \leq_{\mathrm{CP}}xy\,\Psi_i\otimes\Psi_j.
\]
For example, the upper inequality is the average of the two CP maps
\[
 (x\Psi_i+D_i)\otimes(y\Psi_j-D_j),\qquad
 (x\Psi_i-D_i)\otimes(y\Psi_j+D_j).
\]
The other pair of signs gives the lower inequality.  Tensoring with
$\mathcal R$, applying the result to $C_v$, and whitening gives
$\|D\|\leq xy$.  Tensoring the lower single-port bounds also gives
$G\geq(1-x)(1-y)I$.  Lemma~\ref{q:lem-schur} now bounds the
Friedrichs cosine by $C_{ef}$.

Put $h_e=I-G_e$.  The two-projection decomposition implies
\begin{equation}
 \operatorname{Re}\langle h_e\xi,h_f\xi\rangle
 \geq-C_{ef}\|h_e\xi\|\|h_f\xi\|.
 \label{q:projection-cross}
\end{equation}
For completeness, the components of $h_e\xi$ and $h_f\xi$ in their
common range coincide.  Removing this component leaves vectors whose
overlap is controlled by the nontrivial principal angles.  These angles
are the same for a pair of projections and for their complements.
The removed common component contributes a nonnegative squared norm.
This proves \eqref{q:projection-cross}; commuting disjoint constraints
have zero cosine and a nonnegative cross term.

With $z_e=\|h_e\xi\|$, summing yields
\[
 \|H_{\mathrm{block}}\xi\|^2
 \geq z^{\mathsf T}(I-C)z
 \geq[1-\rho(C)]\langle\xi,H_{\mathrm{block}}\xi\rangle.
\]
The spectral theorem gives the desired gap on the product of block
images.  In the full physical space, every one-block image projector
commutes with each $G_e$.  A sector outside any block image has an
incident term equal to the identity, and hence energy at least one.
This proves the stated bound on the full Hilbert space.
\end{proof}

The matrix step is a form of the established subspace-angle method
\cite{q:Kassabov2011}.  Its representation-theoretic interpretation is
useful: averaging over a finite subgroup projects onto its fixed vectors,
and a matrix of pairwise angles controls the gap of the sum of the
corresponding averaging Laplacians.  Here CP transport supplies concrete
entries of that matrix for tensor-network subspaces.

\subsection{Microscopic normalization and decorated AKLT models}

\begin{lemma}[Finite-patch comparison]
\label{q:lem-seed}
Let $H_{\mathrm{micro}}=\sum_a h_a$ be a frustration-free microscopic
Hamiltonian with $h_a\geq0$.  For every coarse edge $e$, let $P_e$ be
its two-block patch and let $H_{P_e}$ be the sum of microscopic terms
contained in it.  Assume
\[
 \ker H_{P_e}=\mathcal R_e,\qquad
 H_{P_e}\geq g_* (I-G_e),\qquad g_*>0.
\]
Suppose every microscopic term lies in at least one and at most $m$
patches.  Then the two Hamiltonians have the same full kernel and
\begin{equation}
 \operatorname{gap}H_{\mathrm{micro}}
 \geq \frac{g_*}{m}\operatorname{gap}H_{\mathrm{block}}.
 \label{q:seed-comparison}
\end{equation}
\end{lemma}

\begin{proof}
Positivity and patch coverage give
$\ker H_{\mathrm{micro}}=\bigcap_e\ker H_{P_e}
=\bigcap_e\mathcal R_e=\ker H_{\mathrm{block}}$.
The operator inequality
$mH_{\mathrm{micro}}\geq\sum_eH_{P_e}\geq g_*H_{\mathrm{block}}$
then proves the claim above their common kernel.
\end{proof}

In the standard AKLT model \cite{q:AKLT1988}, a vertex of full graph degree $D_v$ carries
spin $D_v/2$, and every interaction is the norm-one projector onto the
maximal total spin of its two endpoints.  The spin-one wire channel is
\[
 \Phi(X)=\frac13\sum_{a=1}^3\sigma_aX\sigma_a
 =\frac23\operatorname{tr}(X)I-\frac13X.
\]
Let $\Psi(X)=\operatorname{tr}(X)I/2$.  Then
\[
 \Phi^\ell=\Psi+\beta_\ell(\mathrm{id}-\Psi),\qquad
 \beta_\ell=(-1/3)^\ell.
\]
Relative to the faithful Choi metric of $\Psi$, the Choi eigenvalues
are $1+3\beta_\ell$ once and $1-\beta_\ell$ three times.  Centering
them gives
\begin{equation}
 \delta_\ell=\frac{2}{3^\ell+(-1)^\ell},\qquad\ell\geq2.
 \label{q:aklt-radius}
\end{equation}

\begin{corollary}[Logarithmic sufficient decoration]
\label{q:cor-aklt}
Let $\mathcal G$ be a simple graph of maximum degree $\Delta\geq2$.
Replace every edge by a chain with $n=2\ell$ additional spin-one
vertices, and give each original degree-$d$ vertex spin $d/2$.
For connected finite regions made of whole half-wire blocks, with at
least two blocks, suppose
\[
 2(\Delta-1)\mathcal B(\delta_\ell,\delta_\ell)<1.
\]
Then the standard microscopic AKLT Hamiltonians have the uniform bound
\begin{equation}
 \operatorname{gap}H_{\mathrm{AKLT}}
 \geq\frac{g_*(\Delta,\ell)}\Delta
       \left[1-2(\Delta-1)
                    \mathcal B(\delta_\ell,\delta_\ell)\right]>0,
 \label{q:aklt-gap}
\end{equation}
where $g_*(\Delta,\ell)>0$ is the minimum positive gap of the finitely
many two-block patch Hamiltonians with these parameters.  In particular,
\begin{equation}
 \ell=\left\lceil\log_9(144\Delta)\right\rceil,
 \qquad n=2\ell
 \quad\Longrightarrow\quad
 \operatorname{gap}H_{\mathrm{AKLT}}
 \geq\frac{g_*(\Delta,\ell)}{2\Delta}>0.
 \label{q:log-decoration}
\end{equation}
Thus $n<\log_3(144\Delta)+2$ is a sufficient even decoration count.
\end{corollary}

\begin{proof}
The only microscopic input not yet verified is the regional kernel.
Realize spin $D_v/2$ as the homogeneous binary polynomials of degree
$D_v$ in $(x_v,y_v)$.  For an included edge $vw$, the kernel of the
maximal-spin projector consists precisely of polynomials divisible by
\[
 b_{vw}=x_vy_w-y_vx_w.
\]
Indeed, multiplication by this invariant bracket is injective and
$SU(2)$-equivariant from degrees $(D_v-1,D_w-1)$.  Its image has dimension
$D_vD_w$ and contains no maximal-spin summand.  This equals the
dimension of the complement of that summand in
$(D_v+1)(D_w+1)$ dimensions.  Distinct edge brackets are nonassociate
irreducibles in a polynomial unique factorization domain.  Simultaneous
divisibility is therefore divisibility by their product.  The remaining
degrees are exactly the free boundary virtual degrees, proving the
regional tensor-image identity.

On a half-wire patch, every boundary wire vertex has one omitted edge
and one free virtual leg.  In particular no boundary corner imposes an
additional symmetrization kernel.  The local degree and wire-length
bounds give only finitely many two-block patch types; each finite
positive matrix has a positive least nonzero eigenvalue.  Their minimum
$g_*(\Delta,\ell)$ is positive.  Internal-block microscopic terms occur
in at most $\Delta$ two-block patches, and a term crossing a block
interface occurs once.  Hence Lemma~\ref{q:lem-seed} applies with
$m\leq\Delta$.

Every edge meets at most $2(\Delta-1)$ others, proving
\eqref{q:aklt-gap}.  For the explicit choice
\eqref{q:log-decoration}, one has $\ell\geq3$,
$\delta_\ell\leq3^{1-\ell}\leq1/9$, and
\[
 \mathcal B(\delta_\ell,\delta_\ell)
 \leq4\delta_\ell^2.
\]
The last inequality follows by monotonicity and
$2(1+(1/9)^2)/(1-1/9)^4=3321/1024<4$.  Therefore
\[
 2(\Delta-1)\mathcal B(\delta_\ell,\delta_\ell)
 \leq\frac{72(\Delta-1)}{9^\ell}<\frac12.
\]
\end{proof}

The coefficient $g_*(\Delta,\ell)$ fixes the microscopic energy
normalization.  The result asserts a positive constant for each fixed
degree and decoration; it does not assert a degree-uniform physical
gap.  At small degree one can use the exact criterion rather than the
closed bound: $\ell=2$ gives $\delta=1/5$, $\mathcal B=13/64$, and
coarse gap at least $3/16$ for $\Delta=3$.  At $\ell=3$,
$\mathcal B=85/5184$, and the criterion holds through $\Delta=31$.
These values are sufficient bounds, not optimal-decoration claims.

For any fixed primitive finite-dimensional wire channel, normalized
powers converge to a faithful rank-one map.  Since its limiting Choi
matrix is positive definite, there are $C_0>0$ and $0<\theta<1$ such
that $\delta_\ell\leq C_0\theta^\ell$ for all sufficiently large
$\ell$; a Jordan polynomial is absorbed by increasing $\theta$.
The same theorem then gives sufficient total wire length
\[
 n=\frac{\log\Delta}{|\log\theta|}+O(1).
\]
This conclusion concerns the canonical blocked parent for arbitrary
nonzero cores.  A specified microscopic interaction is covered whenever
the finite regional-kernel and seed-gap hypotheses of
Lemma~\ref{q:lem-seed} hold.

\section{Compatible idempotents as local gap certificates}
\label{q:sec-oblique}

The preceding construction starts from independent edge constraints and
changes their metric.  A related algebraic description starts from
commuting, generally nonorthogonal idempotents.  Orthogonalizing their
ranges produces the physical parent Hamiltonian.  Exact compatibility
is useful because a small commutator of arbitrary orthogonal projections
does not distinguish nearly perpendicular subspaces from nearly parallel
ones.

\begin{theorem}[A compatible-idempotent certificate]
\label{q:thm-oblique}
Let $P_i^2=P_i$ and $[P_i,P_j]=0$ on a finite-dimensional Hilbert space.
Let $\Pi_i$ be the orthogonal projector onto $\operatorname{ran}P_i$,
$H=\sum_i(I-\Pi_i)$, and set
\[
 a_i=\|P_i-P_i^*\|,\qquad
 \kappa_i=\sqrt{1+a_i^2},\qquad
 L_i=\kappa_i+2\kappa_i^2a_i,\qquad
 \epsilon_{ij}=\|[P_i,P_j^*]\|.
\]
For $i\ne j$ define
\begin{equation}
 b_{ij}=\min\left\{
 \frac{a_i}{\kappa_i},\frac{a_j}{\kappa_j},
 2\min(\kappa_i,\kappa_j)L_iL_j\epsilon_{ij}\right\},
 \qquad b_{ii}=0.
 \label{q:oblique-b}
\end{equation}
If the symmetric matrix $B=(b_{ij})$ has $\rho(B)<1$, then
\[
 H\geq[1-\rho(B)](I-\Pi_{\mathcal K}),\qquad
 \mathcal K=\bigcap_i\operatorname{ran}P_i.
\]
Only the individual nonorthogonality and pairwise adjoint-commutator
norms enter the certificate.
\end{theorem}

\begin{proof}
Relative to its range and its orthogonal complement, an idempotent has
the form $P=\left(\begin{smallmatrix}I&A\\0&0\end{smallmatrix}\right)$.
Direct multiplication gives
\begin{equation}
 \Pi(P)=PP^*[I-(P-P^*)^2]^{-1},\qquad
 \|P\|\leq\sqrt{1+\|P-P^*\|^2}.
 \label{q:orthogonalizer}
\end{equation}
The denominator is at least $I$.  For any operator $T$, product and
inverse commutator identities consequently give
\begin{equation}
 \|[\Pi(P),T]\|
 \leq (\kappa+2\kappa^2a)
          \bigl(\|[P,T]\|+\|[P^*,T]\|\bigr).
 \label{q:ortho-comm}
\end{equation}
Indeed the numerator costs $\kappa$ times the bracketed sum, and the
commutator of $I-(P-P^*)^2$ costs $2a$ times that sum.  Applying
\eqref{q:ortho-comm} twice and using exact compatibility gives
\[
 \|[\Pi_i,\Pi_j]\|\leq2L_iL_j\epsilon_{ij}.
\]

To control the nearly parallel possibility, quotient by
$W=\operatorname{ran}P_i\cap\operatorname{ran}P_j$.  The two idempotents
descend with norms no larger than $\kappa_i,\kappa_j$, and their product
is zero.  For vectors $u,v$ in their respective quotient ranges,
\[
 \|u\|\leq\kappa_i\|u-tv\|\qquad(t\in\mathbb C).
\]
Minimizing in $t$ bounds their cosine by
$\sqrt{1-\kappa_i^{-2}}=a_i/\kappa_i$, and symmetrically by
$a_j/\kappa_j$.  On a nontrivial principal-angle block with cosine $c_r$,
the orthogonal-projection commutator has norm
$c_r\sqrt{1-c_r^2}$.  The preceding bounds imply
$\sqrt{1-c_r^2}\geq1/\min(\kappa_i,\kappa_j)$.
Taking the maximal cosine therefore gives
\[
 c_{ij}\leq\min(\kappa_i,\kappa_j)\|[\Pi_i,\Pi_j]\|
 \leq2\min(\kappa_i,\kappa_j)L_iL_j\epsilon_{ij}.
\]
This proves $c_{ij}\leq b_{ij}$.  The argument
\eqref{q:projection-cross} now gives $H^2\geq[1-\rho(B)]H$ and the
claimed spectral consequence.
\end{proof}

\begin{corollary}[Filtered commuting-projector models]
\label{q:cor-filtered}
Let $E_i$ be commuting local orthogonal projections and let
$L=\bigotimes_x L_x$ be any product of invertible site filters.  Then
$P_i=LE_iL^{-1}$ are local compatible idempotents.  Theorem
\ref{q:thm-oblique} applies to the physical Hamiltonian
\[
 H_L=\sum_i\bigl(I-\Pi_{L\operatorname{ran}E_i}\bigr),
 \qquad
 \ker H_L=L\bigcap_i\operatorname{ran}E_i.
\]
In particular, suppose every reference support has at most $s$ sites,
meets at most $M$ other supports, and $L_x=e^{K_x}$ with
$\|K_x\|\leq\varepsilon$.  If
\[
 2M(e^{2s\varepsilon}-1)<1,
\]
then $\operatorname{gap}H_L\geq1-2M(e^{2s\varepsilon}-1)$.
\end{corollary}

\begin{proof}
Factors of $L$ outside a reference support cancel in $P_i$, proving
locality.  Similarity preserves idempotence, commutativity, and the
intersection of the ranges.  On a support of at most $s$ sites,
integration of the similarity derivative gives
$\|P_i-E_i\|\leq e^{2s\varepsilon}-1=:\delta$.
Thus $a_i\leq2\delta$, and \eqref{q:oblique-b} gives
$b_{ij}\leq2\delta$.  Disjoint supports give $b_{ij}=0$, so the
maximum row sum is at most $2M\delta$.
\end{proof}

The common ground space in this corollary may be degenerate, including
commuting-projector code spaces.  The argument estimates a concrete
finite family of local matrices; it does not require a bound on the
condition number of the global product $L$.

\subsection{A parent from integer modular conjugates}

Let $\rho>0$ on $\bigotimes_x\mathbb C^{d_x}$ and use fixed product
bases for vectorization.  On two copies of the system let $E_x$ be the
normalized Bell projector at site $x$, and set
\[
 P_x=(\rho\otimes I)E_x(\rho^{-1}\otimes I).
\]
They are compatible idempotents with common range
$\mathbb C|\rho\rangle\!\rangle$.  For a unitary Weyl basis
$\{U_{x,a}\}_{a=1}^{d_x^2}$,
\begin{equation}
 P_x=\frac1{d_x^2}\sum_a
       (\rho U_{x,a}\rho^{-1})\otimes\overline{U_{x,a}}.
 \label{q:integer-parent}
\end{equation}
This is an exact finite-dimensional identity, so local estimates of the
integer modular conjugates give the entries of
\eqref{q:oblique-b} directly.  For example, if
$\sigma_x=I_x/d_x\otimes\operatorname{tr}_x\rho$ and
$Y_x=\rho\sigma_x^{-1}$, then for every one-site $O_x$,
\begin{equation}
 \rho O_x\rho^{-1}=Y_xO_xY_x^{-1},
 \label{q:ratio-modular}
\end{equation}
because $\sigma_x$ commutes with $O_x$.

The normalized vector $|\rho\rangle\!\rangle$ purifies
$\rho^2/\operatorname{tr}(\rho^2)$.  For a generic $\rho$ the operators
in \eqref{q:integer-parent} can have large support; locality follows
when the supplied modular conjugates are local or admit certified
decaying approximations.  This construction is related to canonical
purified parents and their mixing criteria \cite{q:LPP2025}, and to
factorizations of exact KMS parent Hamiltonians used in quantum
algorithms \cite{q:LJL2026}.  The additional input here is the explicit
compatible-idempotent certificate \eqref{q:oblique-b}.

\section{Positive modular response calculus}
\label{q:sec-kms}

For a differentiable positive matrix path, its right logarithmic
derivative need not be Hermitian.  Keeping its Hermitian and
anti-Hermitian parts together reveals a useful cancellation.  The
functional calculus below uses positive-definite scalar kernels, a
standard source of operator inequalities
\cite{q:HK1999,q:BP2000}.  We give the kernels explicitly so that their
norms, ancilla extensions, and composition rules are available without
an implicit functional-calculus convention.

Let $\rho_\lambda>0$, $K=\log\rho$, and
\[
 Q=\dot\rho\rho^{-1}=H+A=H+iJ,
 \qquad H^*=H,\quad A^*=-A,\quad J^*=J.
\]
Throughout this section powers of $\rho$ are not renormalized to have
trace one.  Renormalization adds the scalar derivative of their
normalizing factor.

\begin{theorem}[Exact response of fractional powers]
\label{q:thm-powers}
For $0\leq r\leq1$, let
$Q_r=(\partial_\lambda\rho^r)\rho^{-r}=H_r+A_r$ be its Hermitian and
anti-Hermitian decomposition.  With $\operatorname{ad}_K(O)=[K,O]$,
\begin{align}
 H_r&=\psi_r(\operatorname{ad}_K)H,
 &\psi_r(x)&=\frac{\sinh(rx)}{\sinh x},\quad\psi_r(0)=r,
 \label{q:power-H}\\
 A_r&=\chi_r(\operatorname{ad}_K)A,
 &\chi_r(x)&=\left(\frac{\sinh(rx/2)}{\sinh(x/2)}\right)^2,
 \quad\chi_r(0)=r^2.
 \label{q:power-A}
\end{align}
Both maps are completely positive, with completely bounded
operator norms $r$ and $r^2$, respectively.  Consequently
\[
 \|H_r\|\leq r\|H\|,\qquad
 \|A_r\|\leq r^2\|A\|.
\]
The estimates hold with an arbitrary ancillary system.  Moreover,
\begin{equation}
 \lim_{r\downarrow0}\frac{Q_r}{r}
 =\dot K=f(\operatorname{ad}_K)H,
 \qquad f(x)=\frac{x}{\sinh x},\quad f(0)=1.
 \label{q:log-response}
\end{equation}
\end{theorem}

\begin{proof}
Differentiation of the exponential gives
\[
 Q_r=\int_0^r e^{sK}\dot K e^{-sK}\,ds
     =\frac{e^{r\operatorname{ad}_K}-I}
                  {\operatorname{ad}_K}\dot K.
\]
Taking Hermitian and anti-Hermitian parts yields
\[
 H_r=\frac{\sinh(r\operatorname{ad}_K)}{\operatorname{ad}_K}\dot K,
 \qquad
 A_r=\frac{\cosh(r\operatorname{ad}_K)-I}
                    {\operatorname{ad}_K}\dot K.
\]
Comparison with $r=1$ proves the formulas, using
$\cosh(rx)-1=2\sinh^2(rx/2)$.  All quotients have their continuous
value at zero.

For $0<r<1$, the positive density
\begin{equation}
 p_r(t)=\frac{\sin(\pi r)}
              {2[\cosh(\pi t)+\cos(\pi r)]}
 \label{q:pr-kernel}
\end{equation}
satisfies
\begin{equation}
 \psi_r(x)=\int_{\mathbb R}p_r(t)e^{itx}\,dt,
 \qquad \int_{\mathbb R}p_r(t)\,dt=r.
 \label{q:pr-transform}
\end{equation}
One direct verification integrates
$e^{ikz}/(\cosh z+\cos\alpha)$ around a rectangle of height $2\pi i$.
The residues at $i(\pi-\alpha)$ and $i(\pi+\alpha)$ give
\[
 \int_{\mathbb R}\frac{e^{iku}}{\cosh u+\cos\alpha}\,du
 =\frac{2\pi}{\sin\alpha}
      \frac{\sinh(\alpha k)}{\sinh(\pi k)}.
\]
Evenness and continuity include $k\leq0$.  Taking $\alpha=\pi r$,
$u=\pi t$, and $k=x/\pi$ proves \eqref{q:pr-transform}.
Thus
\[
 \psi_r(\operatorname{ad}_K)(O)
 =\int_{\mathbb R}p_r(t)e^{itK}Oe^{-itK}\,dt
\]
is a positive mixture of unitary conjugations of mass $r$.  It is
completely positive with completely bounded norm $r$, attained on $I$.
Since $\chi_r(x)=\psi_r(x/2)^2$, its positive Fourier density is
$(\widetilde p_r*\widetilde p_r)(t)$, where
$\widetilde p_r(t)=2p_r(2t)$, and its total mass is $r^2$.
The endpoints are the zero and identity maps.  Finally
$p_r/r$ converges in $L^1$ to
\[
 p_0(t)=\frac\pi4\operatorname{sech}^2(\pi t/2),
\]
a probability density with Fourier transform $f$.  Alternatively the
limit in \eqref{q:log-response} follows directly by differentiating
$\rho^r=I+rK+O(r^2)$.
\end{proof}

The full response obeys the exact KMS compatibility relation
\begin{equation}
 Q^*=e^{-\operatorname{ad}_K}Q,
 \qquad A=\tanh(\operatorname{ad}_K/2)H.
 \label{q:kms-compatibility}
\end{equation}
Thus its anti-Hermitian part is constrained by its Hermitian part and
the modular generator.  It is not an independent input.  The
fractional-power maps compose coherently:
\[
 \psi_s(rx)\psi_r(x)=\psi_{rs}(x),\qquad
 \chi_s(rx)\chi_r(x)=\chi_{rs}(x).
\]
In particular, square roots give
\[
 H_{1/2}=\tfrac12\operatorname{sech}(\operatorname{ad}_K/2)H,
 \qquad
 A_{1/2}=\tfrac14\operatorname{sech}^2(\operatorname{ad}_K/4)A.
\]

\begin{proposition}[A positive rotation generator]
\label{q:prop-rotation}
Define the even function
\[
 h(x)=\frac{1-x/\sinh x}{x\tanh(x/2)},\qquad h(0)=\frac13,
 \qquad C=h(\operatorname{ad}_K)J.
\]
Then $C=C^*$ and
\begin{equation}
 \dot K=H+i[C,K],\qquad \|C\|\leq\frac13\|J\|.
 \label{q:rotation}
\end{equation}
The map $h(\operatorname{ad}_K)$ is completely positive with completely
bounded norm $1/3$.  Its Fourier density has finite exponentially
weighted mass for every weight exponent $v<2\pi$.
\end{proposition}

\begin{proof}
Differentiating the Mittag--Leffler expansion of $x\coth(x/2)$ gives
\[
 h(x)=\frac1x\frac{d}{dx}[x\coth(x/2)]
     =8\sum_{n\geq1}
       \frac{(2\pi n)^2}{[x^2+(2\pi n)^2]^2}.
\]
Its nonnegative Fourier density is
\begin{align*}
 k(t)
 &=\sum_{n\geq1}\left(\frac1{\pi n}+2|t|\right)e^{-2\pi n|t|}\\
 &=-\frac1\pi\log(1-e^{-2\pi|t|})
       +\frac{2|t|}{e^{2\pi|t|}-1}.
\end{align*}
The logarithmic singularity at zero is integrable, and its total mass
is $h(0)=1/3$.  The positive-mixture argument proves the CP and norm
claims.  For $0\leq v<2\pi$ its weighted mass is
\[
 2\sum_{n\geq1}\left[
 \frac1{\pi n(2\pi n-v)}+
 \frac2{(2\pi n-v)^2}\right]<\infty.
\]
Equation \eqref{q:kms-compatibility} and the scalar identity
$xh(x)\tanh(x/2)=1-f(x)$ give
$i[C,K]=-[1-f(\operatorname{ad}_K)]H$, proving
\eqref{q:rotation}.
\end{proof}

If $\dot U=iCU$ and $U(0)=I$, then
\[
 \frac{d}{d\lambda}(U^*KU)=U^*HU.
\]
This separates a positive path into Hermitian forcing and unitary
rotation.  The operator-norm estimates are dimension independent.
Spatial locality, however, depends on the support growth of the
conjugations $e^{itK}(\cdot)e^{-itK}$.  The next two propositions
describe precisely why ordinary norm contraction alone does not settle
that spatial question.

\section{Locality distinctions for conditional ratios and logarithms}
\label{q:sec-obstructions}

\subsection{Local conditional ratios with a three-body logarithm}

An algebraic conditional ratio and a Hermitian conditional logarithm
encode different forms of locality.  Their distinction is related to
Belavkin--Staszewski quantum Markov chains, whose structure and recovery
maps are studied in \cite{q:BCCRJ2026}.  The following small example
quantifies the distinction and persists arbitrarily close to a product
state.

\begin{proposition}[Exact local ratios, nonlocal log difference]
\label{q:prop-ratio-obstruction}
There are faithful three-qubit states with every one-site conditional
ratio supported on the closed star of that site in the path $A-B-C$,
arbitrarily close to the maximally mixed state in operator norm, such that
$\log\rho-\log\mathbb E_A\rho$ has a nonzero $ABC$ Pauli coefficient.
Here $\mathbb E_A(T)=I_A/2\otimes\operatorname{tr}_A T$.
\end{proposition}

\begin{proof}
For $0<a,b,t<1$, take
\[
 \sigma_B=\frac{I+tX_B}{2},\quad
 U_{AB}=I+aZ_AZ_B,\quad V_{BC}=I+bZ_BZ_C,
 \quad
 \rho=\frac14\sigma_B^{1/2}U_{AB}V_{BC}\sigma_B^{1/2}.
\]
The middle factors commute and are positive, so $\rho$ is faithful.
Partial traces give $\rho_B=\sigma_B$ and
\[
 \rho_{AB}=\tfrac12\sigma_B^{1/2}U_{AB}\sigma_B^{1/2},\qquad
 \rho_{BC}=\tfrac12\sigma_B^{1/2}V_{BC}\sigma_B^{1/2}.
\]
Consequently
\begin{equation}
 \rho=\rho_{AB}\rho_B^{-1}\rho_{BC},\qquad
 \rho(\mathbb E_A\rho)^{-1}
       =\sigma_B^{1/2}U_{AB}\sigma_B^{-1/2}.
 \label{q:bs-recovery}
\end{equation}
The ratio on the right, its inverse, and its matrix logarithm act on
$AB$.  The analogous ratio at $C$ acts on $BC$.

Set $a=b=s$ and put
\[
 g(u)=\frac{\operatorname{atanh}u}{u}
      =\sum_{k\geq0}\frac{u^{2k}}{2k+1},\qquad
 v=\sqrt{t^2+\frac{4s^2(1-t^2)}{(1+s^2)^2}}.
\]
In the $Z_A=z,Z_C=w$ block, $8\rho$ equals
\[
 (1+s^2zw)I+(1+s^2zw)tX_B
           +s(z+w)\sqrt{1-t^2}\,Z_B.
\]
For a positive $2\times2$ matrix $cI+u\cdot\sigma$, the vector
coefficient of its logarithm is
$[\operatorname{atanh}(|u|/c)/|u|]u$.  Hence its $X_B$ coefficient is
$tg(v)$ in the two blocks $z=w$ and $tg(t)$ in the other two.  For
$P=Z_AX_BZ_C$, the signed average gives
\begin{equation}
 \frac18\operatorname{tr}(P\log\rho)
 =\frac t2[g(v)-g(t)]
 \geq\frac{2ts^2(1-t^2)}{3(1+s^2)^2}>0.
 \label{q:three-body-coefficient}
\end{equation}
The logarithm of $\mathbb E_A\rho$ acts on $BC$, so it has zero
$P$ coefficient.

The endpoint ratio satisfies
\[
 \|\rho(\mathbb E_A\rho)^{-1}-I\|
 =a\sqrt{\frac{1+t}{1-t}},
\]
and the middle ratio obeys
\[
 \|\rho(\mathbb E_B\rho)^{-1}-I\|
 \leq t+\frac{\sqrt{1-t^2}(a+b)}{1-ab}.
\]
Taking $a=b=t\downarrow0$ proves the smallness assertion while
\eqref{q:three-body-coefficient} stays strictly positive.
\end{proof}

The same construction on qutrits uses
$Z=\operatorname{diag}(1,-1,0)$,
$X=|0\rangle\langle1|+|1\rangle\langle0|$,
$\sigma_B=(I_3+tX_B)/3$, and prefactor $1/9$ in $\rho$.
It is real symmetric and invariant under spin-one time reversal in a
Cartesian basis.  Restricting the endpoint eigenvalues to $\pm1$ gives
the same nonzero normalized three-site coefficient.  Thus the
distinction is also present in the time-reversal-invariant spin-one
setting.  It concerns strict support; the generated distance-two term
is fully consistent with an exponentially decaying support expansion.

\subsection{A CP contraction need not contract an interaction norm}

For $n$ sites and weights $w(m)\geq1$, $w(1)=1$, define the intrinsic
interaction seminorm, modulo scalars, by
\begin{equation}
 \|O\|_w=
 \inf_{O=cI+\sum_{\varnothing\ne S\subseteq[n]}O_S}
 \sup_x\sum_{S\ni x}w(|S|)\|O_S\|,
 \label{q:intrinsic-norm}
\end{equation}
where $O_S$ acts on $S$ and as identity elsewhere.  The infimum makes
this a property of the operator, rather than of a particular expansion.

\begin{proposition}[No generic same-norm modular bound]
\label{q:prop-filter-obstruction}
For every $a,b>0$ and every such sequence of weights, there are
Hermitian $K_n,H_n$ with
\[
 \|K_n\|_w\leq a+b/n,\qquad\|H_n\|_w=1,
\]
but
\begin{equation}
 \|f(\operatorname{ad}_{K_n})H_n\|_w
 \geq n\frac{ab}{a^2+b^2}[1-f(2b)]\longrightarrow\infty,
 \qquad f(x)=x/\sinh x.
 \label{q:interaction-blowup}
\end{equation}
In particular, no bound
$\|f(\operatorname{ad}_K)H\|_w\leq C(\|K\|_w)\|H\|_w$
with locally bounded $C$ holds on this full class, even near $K=0$.
\end{proposition}

\begin{proof}
Set
\[
 D_n=\sum_{j=1}^n Z_j,\qquad V_n=X_1\cdots X_n,\qquad
 e_n=\frac a{w(n)},\qquad
 K_n=\frac bnD_n+e_nV_n,
 \qquad H_n=D_n.
\]
These displayed decompositions give the upper bounds.  For $H_n$, its
spectral diameter $2n$ gives the reverse bound one, since for any
decomposition in \eqref{q:intrinsic-norm},
\[
 \tfrac12\operatorname{diam}\operatorname{spec}(O)
 \leq\sum_S\|O_S\|
 \leq n\sup_x\sum_{S\ni x}w(|S|)\|O_S\|.
\]
The subspace spanned by $|0^n\rangle,|1^n\rangle$ is reducing for
$K_n,H_n$, on which they equal $bZ+e_nX$ and $nZ$.  Put
$r_n=\sqrt{b^2+e_n^2}$.  The modular filter fixes the Bloch component
parallel to $(e_n,0,b)$ and multiplies the perpendicular component by
$f(2r_n)$.  Hence
\[
 \langle0^n|f(\operatorname{ad}_{K_n})H_n|1^n\rangle
 =\frac{nbe_n}{r_n^2}[1-f(2r_n)].
\]
Every proper-support operator has zero matrix element between these
two product vectors.  Thus every decomposition in
\eqref{q:intrinsic-norm} has a full-support term of norm at least this
matrix element.  Multiplying by $w(n)$ and using
$b\leq r_n\leq\sqrt{a^2+b^2}$ and monotonic decrease of $f$ on
$(0,\infty)$ proves \eqref{q:interaction-blowup}.
\end{proof}

This example also explains the role of KMS compatibility.  A Hermitian
input $H$ determines the compatible tangent uniquely: in a $\rho$
eigenbasis with eigenvalues $\lambda_i$,
\[
 \dot\rho_{ij}=\frac{2\lambda_i\lambda_j}{\lambda_i+\lambda_j}H_{ij}.
\]
The resulting full response is
$Q=[I+\tanh(\operatorname{ad}_K/2)]H$.  In the reducing two-dimensional
block above, its additional anti-Hermitian part is
\[
 -i\frac{ne_n}{r_n}\tanh(r_n)Y.
\]
Therefore $\|Q_n\|_w\geq na\tanh(r_n)/r_n$, which also diverges.
Proposition~\ref{q:prop-filter-obstruction} concerns an arbitrary
Hermitian input, and does not rule out spatial estimates that use a
separately controlled full compatible response.  Together with
Theorem~\ref{q:thm-powers}, it identifies this full response as the
appropriate datum for positive-path reconstruction.

\subsection{Computational verification and formalization interfaces}

The proofs above are analytic finite-dimensional arguments.  Their
supplementary computations serve two complementary roles.  Exact
rational Choi comparisons and positive-vector inequalities are
checkable certificates for concrete channels and graph types.  Floating
point calculations independently test tensor conventions and matrix
identities; they are not used to establish the universal bounds.

The supplied two-port tests compare the quotient Schur formula with
orthogonal projectors computed from the square-root metric.  In 600
deterministically seeded examples, including nearly rank-one and
ill-conditioned cores, their largest discrepancy was
$2.75\times10^{-14}$.  Separately, direct contraction of physical AKLT
star tensors at degrees $2,3,4$ and half-wire lengths $1,2$ agreed with
\eqref{q:gram-channel} within $5\times10^{-16}$ relative error.
The modular checks comprise 20 Fourier-integral comparisons, 28 matrix
response comparisons, and 20 obstruction checks.  The supplement records
the seeds, dimensions, formulas, and complete outputs.

For formal verification, the main network proof separates into CP-order
tensoring, a Schur-complement lemma, the two-projection angle inequality,
and a polynomial divisibility proof for the AKLT regional kernel.  The
modular proof separates scalar Fourier identities from finite-dimensional
functional calculus.  This division exposes a small set of reusable
lemmas; the numerical checks do not discharge their universally
quantified proof obligations.


\section{Haar measurement simulation and the geometry of contact}
\label{geo:qutrit}

The same projective moments have an operational use beyond ensemble
certification. A \emph{parent measurement} is a fixed measurement from
which a family of measurements can be obtained by classical
postprocessing. In dimension \(d\), the unitary-invariant pure-state
parent is
\[
 G(dz)=dP_z\,d\sigma(z).
\]
For \(m\) labelled outcomes its postprocessing body is
\[
 \mathcal H_{d,m}=
 \left\{\left(d\int p_i(z)P_z\,d\sigma(z)\right)_{i=1}^m:
      p_i\ge0,\quad\sum_i p_i=1\ {\rm a.e.}\right\}.
\]
It is compact and convex by weak-star compactness of bounded response
functions. Its support function,
\[
 h(W)=d\int\max_i\operatorname{tr}(W_iP_z)\,d\sigma(z),
 \label{geo:haar-support}
\]
turns simulation into a comparison between a classical maximum and
quantum measurement optimization. This convex-body description belongs
to the established theory of steering and measurement compatibility
\cite{geo_nguyen2018}.

Let
\[
 \Phi_\eta(A)=\eta A+\frac{1-\eta}{d}\operatorname{tr}(A)I,\qquad
 \eta_d=\frac{H_d-1}{d-1},\qquad c_d=\frac{1-\eta_d}{d}.
\]
The notation \(d\) in this section denotes Hilbert-space dimension.
The number \(\eta_d\) is fixed by the exact maximum-coordinate cells
\(R_a=\{z:|z_a|^2\ge|z_b|^2\text{ for every }b\}\):
\begin{equation}
 d\int_{R_a}P_z\,d\sigma(z)=\eta_d P_a+c_dI.
 \label{geo:haar-cell}
\end{equation}
Indeed, \((|z_1|^2,\ldots,|z_d|^2)\) is uniform on the simplex and
\(\mathbb E\max_a|z_a|^2=H_d/d\).
Thus the same expected-maximum computation governing the entropy
threshold also determines the noise level at which every basis
measurement is a postprocessing of one Haar parent.

\fig{fig10_measurement_cells}{The Haar parent at qutrit visibility $5/12$.
(a) Squared coordinate moduli of a Haar qutrit are uniform on the simplex.
Assigning each point to its largest coordinate partitions it into three
cells; phase integration removes off-diagonal matrix entries. (b) The three
postprocessed effects have eigenvalues $11/18$ in their selected direction
and $7/36$ in the other two directions. Segment heights are the three
diagonal eigenvalues, colored by coordinate direction; each complete column
sums to one. Hence each equals
$(5/12)P_a+(7/36)I$. The neighborhood theorem extends this exact basis-cell
construction to the specified classes of nearby measurements; the figure
itself depicts the basis case.}{fig:measurement}

Write \(\mathsf Q_{d,m}\) for the compact set of labelled \(m\)-outcome
POVMs, and \(\mathsf P_{d,m}\) for the convex hull of its commuting
members. This is the standard projective-simulable set
\cite{geo_oszmaniec2017}.
For qutrits, let \(\mathsf R_m\) consist of all POVMs with a common
nontrivial reducing subspace, and set
\(\mathsf K_m=\operatorname{conv}\mathsf R_m\). A qutrit reducing
decomposition is a two-dimensional block plus an orthogonal flag.
The decomposition can vary between convex-combination terms.

\begin{theorem}[Uniform operational neighborhoods]
\label{geo:operational}
For every fixed \(d\ge2\) and finite \(m\), there exists \(r_{d,m}>0\)
such that
\[
 E\in\mathsf Q_{d,m},\quad
 \operatorname{dist}(E,\mathsf P_{d,m})<r_{d,m}
 \quad\Longrightarrow\quad
 \Phi_{\eta_d}(E)\in\mathcal H_{d,m}.
 \tag{i}
\]
For \(d=3\), there exists \(r_m>0\) with the stronger conclusion
\[
 E\in\mathsf Q_{3,m},\quad
 \operatorname{dist}(E,\mathsf K_m)<r_m
 \quad\Longrightarrow\quad
 \Phi_{5/12}(E)\in\mathcal H_{3,m}.
 \tag{ii}
\]
Any fixed norm on tuples can define distance. Zero labelled effects
are allowed, and nearby POVMs may be irreducible and have arbitrary
effect ranks. The radii are qualitative; no value uniform in
unbounded \(m\) is asserted.
\end{theorem}

In particular, every reducible qutrit POVM and every convex mixture of
such measurements is simulated by the same Haar parent at \(5/12\).
Every qutrit POVM with at most three outcomes lies in this class.
To check the last statement, two effects with overlapping support
admit opposite small Hermitian perturbations supported in their
intersection, contradicting extremality. In dimension three, an
extreme ternary POVM therefore has either three rank-one effects,
hence is projective, or rank pattern \((2,1,1)\), with smaller cases
allowed. In the latter case choose a ray \(f\) orthogonal to the
supports of both rank-one effects. Those effects annihilate \(f\),
and the remaining effect, being their complement in \(I\), satisfies
\(E_1f=f\). This is a common reducing flag. The finite-dimensional
convex decomposition into extreme POVMs finishes the assertion.

\begin{remark}
The theorem is an open-neighborhood result around a specified compact
class. It does not assert that all irreducible qutrit POVMs are
simulable at \(5/12\). The contact theorem below is fully anisotropic
on the flagged stratum: its restriction is geometric support, not a
restriction on the spectra of the contact matrices.
\end{remark}

\subsection{Contact inequalities and martingale transport}
\label{geo:contact}

Quantum discrimination gives the semidefinite duality
\[
 q(W):=\max_{E\in\mathsf Q_{d,m}}
              \sum_i\operatorname{tr}(\Phi_{\eta_d}(W_i)E_i)
       =\min_{Y\succeq\Phi_{\eta_d}(W_i)}\operatorname{tr}Y.
\]
For a dual optimum put \(C_i=Y-\Phi_{\eta_d}(W_i)\succeq0\).
Complementary slackness says that an active rank-one effect lies in
\(\ker C_i\); its weighted rays are balanced. Define
\[
 F_d(C)=\int\min_i
       \{z^*C_i z-c_d\operatorname{tr}C_i\}\,d\sigma(z).
\]
Since
\(\Phi_{\eta_d}^{-1}(C)=
\eta_d^{-1}(C-c_d\operatorname{tr}(C)I)\), one has the exact identity
\begin{equation}
 h(W)=\operatorname{tr}Y-\frac d{\eta_d}F_d(C).
 \label{geo:contact-identity}
\end{equation}
Thus a nonpositive contact integral controls the support comparison.
The general contact formulation appears in
\cite[Corollary 2, Eq.~(35)]{geo_yzhang2026}; the identity above also
provides its direct derivation in the present normalization.

There is a second useful interpretation. For a rank-one POVM
\(E_i=t_iP_i\), let \(w_i=t_i/d\) and
\(\rho_i=\eta P_i+(1-\eta)I/d\).
A Haar postprocessing of \(\Phi_\eta(E)\) is exactly a joint law of
the label \(i\) and a Haar pure projector \(P\) satisfying
\[
 \mathbb P(i)=w_i,\qquad \mathbb E[P\mid i]=\rho_i.
\]
Consequently, by the classical martingale coupling theorem
\cite{geo_strassen1965}, simulation is equivalent to
\[
 \sum_iw_i\delta_{\rho_i}\preceq_{\rm cx}\sigma
\]
in the real affine space of trace-one Hermitian matrices.
This is ordinary classical convex order on a matrix-valued state
space. The contact functions are specific convex-envelope tests for
that order; testing only covariance or quadratic polynomials is weaker.

\subsection{The anisotropic qutrit flag theorem}
\label{geo:flag}

\begin{theorem}[Arbitrary contacts on a flagged qutrit frame]
\label{geo:flag-theorem}
Let \(\mathbb C^3=H\oplus\mathbb Cf\), \(\dim H=2\), and let
\(P_i\) be finitely many rays in \(H\) with
\(\sum_i t_iP_i=I_H\), \(t_i\ge0\). Include \(P_f\).
For every collection of positive matrices \(C_i,C_f\) with
\(C_iP_i=0\) and \(C_fP_f=0\), one has
\[
 F_3(C)\le0,\qquad \eta_3=\frac5{12},\quad c_3=\frac7{36}.
\]
\end{theorem}

\begin{proof}
Average the contacts under the flag phase
\(I_H\oplus e^{i\chi}\). Concavity of \(F_3\) and Haar invariance
show that this can only increase \(F_3\), so it suffices to consider
\[
 C_i=a_i(I_H-P_i)+b_iP_f,\quad a_i,b_i\ge0,\qquad
 C_f\succeq0\text{ on }H.
\]
In Bloch notation \(P_i=(I_H+n_i\cdot\sigma)/2\), balance means
\(0\in\operatorname{conv}\{n_i\}\). Carath\'eodory reduces to at most
four rays: deleting contacts increases the pointwise minimum.
One or two rays give a basis case by \eqref{geo:haar-cell}.

For three selected rays \(i,j,k\), put
\[
 w=a_in_i+a_jn_j-2a_kn_k,\quad
 v=a_in_i-a_jn_j,\quad A=a_i+a_j+2a_k.
\]
Choose a qubit basis with Bloch axis \(n\). On its first
maximum-coordinate cell assign the better of contacts \(i,j\), on
the second assign \(k\), and on the flag cell assign \(f\).
The exact diagonal moments cancel the flag coefficients. Averaging
the absolute value of the pair difference over its relative phase
gives
\begin{equation}
 F_3(C)\le\frac5{144}
 \left(A-n\cdot w-t_*|v\times n|\right),
 \quad
 t_*=\frac{36K_0}{5},\quad
 K_0=\frac4\pi\int_{R_1}\sqrt{p_1p_2}\,d\sigma.
 \label{geo:pair-cell}
\end{equation}
The phase inequality is
\(\mathbb E_\phi|b+a\cos\phi|\ge2|a|/\pi\).
An elementary one-dimensional integration gives
\[
 \int_{R_1}\sqrt{p_1p_2}\,d\sigma
 =\frac23\int_{1/2}^1\sqrt{u(1-u)}
                       \bigl(1-(1+u)^{-3}\bigr)\,du
 \ge\frac{19\pi}{648}+\frac{37}{3888}.
\]
Using \(\pi<22/7\),
\begin{equation}
 t_*\ge\frac{2767}{2970}>\frac9{10},
 \qquad t_*-\frac9{10}\ge\frac{47}{1485}.
 \label{geo:strict-reserve}
\end{equation}
The needed geometric statement is that among any at most four
balanced Bloch rays, for arbitrary \(a_i\ge0\), some triple satisfies
\begin{equation}
 |w|^2+\frac{81}{100}|v|^2+
             \frac95|w\times v|\ge A^2.
 \label{geo:four-selection}
\end{equation}
Its proof is given below, including the nearest-facet estimates.
Finally,
\[
 \max_{|n|=1}\{n\cdot w+t|v\times n|\}
 =\sqrt{|w|^2+t^2|v|^2+2t|w\times v|}.
\]
Equations \eqref{geo:four-selection} and \eqref{geo:strict-reserve}
make the right side of \eqref{geo:pair-cell} nonpositive.
\end{proof}

% The geometric proof may instead be moved to a technical appendix.

For a reducible POVM, pinching any separating witness to the same
two-plus-one block leaves its score unchanged and decreases Haar
support. The corresponding quantum SDP splits into a qubit measurement
and assignment of the flag. An extreme qubit optimizer has at most
four rank-one effects, with a balanced set of Bloch rays.
Complementary slackness gives exactly the contacts in
Theorem~\ref{geo:flag-theorem}; \eqref{geo:contact-identity} rules out
separation. Convexity then proves
\(\Phi_{5/12}(\mathsf K_m)\subseteq\mathcal H_{3,m}\).

\subsection{Why local contact control implies an operational neighborhood}
\label{geo:op-proof}

The main distinction is between the rays of an external POVM and
the rays of an optimizing dual contact frame. A compactness argument
must bridge them. The following estimate supplies the required control.

\begin{lemma}[A negative cusp for pinched clusters]
\label{geo:cusp}
Suppose a balanced \(m\)-ray frame is partitioned into \(d\) nonempty
groups, with actual maximum within-group projector distance
\(\delta\le1/(2d)\). Suppose
\[
 a(I-P_i)\preceq C_i\preceq b(I-P_i),\qquad 0<a\le b.
\]
For a constant \(\kappa_d>0\) depending only on dimension,
\begin{equation}
 F_d(C)\le9\eta_db\delta^2-\frac{\kappa_da}{m}\delta.
 \label{geo:negative-cusp}
\end{equation}
One may take \(\kappa_3=325/536313\).
\end{lemma}

\begin{proof}
For group \(I_a\), form its weighted mean
\(A_a=T_a^{-1}\sum_{i\in I_a}t_iP_i\). Its largest eigenvalue is
at least \(1-\delta^2\), and its remaining spectral mass at most
\(\delta^2\). If \(Q_a=u_au_a^*\) is its principal projector,
\(\|P_i-Q_a\|\le\sqrt2\delta\).
Balance gives
\(\sum_aT_a\lambda_aQ_a=I-S\), where \(S\succeq0\) and
\(\operatorname{tr}S\le d\delta^2\). No group can have zero total
weight, since \(I-S\) is invertible.
The polar decomposition of the matrix with columns
\(\sqrt{T_a\lambda_a}u_a\) gives an orthonormal basis \(e_a\)
with \(\|P_i-P_{e_a}\|<3\delta\) in group \(a\).
Consequently \(e_a^*C_ie_a\le9b\delta^2\).

On basis cell \(R_a\), bound the minimum by the uniform average of
its group's contacts. Equation \eqref{geo:haar-cell} bounds the
total integrated average by \(9\eta_db\delta^2\).
The functional
\[
 D\longmapsto\int_{R_a}
       |z^*Dz-c_d\operatorname{tr}D|\,d\sigma
\]
is a norm on Hermitian matrices: if its polynomial vanishes on the
open cell, then \(D=c_d\operatorname{tr}(D)I\), and
\(1-dc_d=\eta_d>0\) forces \(D=0\).
It is therefore at least \(\kappa_d\|D\|\).
In a group attaining diameter \(\delta\), choose \(i,j\) attaining it.
The average minus the minimum is at least \(|f_i-f_j|/|I_a|\).
Moreover, a unit vector in \(\operatorname{ran}P_j\) gives
\(\|C_i-C_j\|\ge a\delta\). This yields the negative linear term.
If \(\delta=0\), the balanced frame is a basis refinement, and
\eqref{geo:haar-cell} gives \(F_d\le0\) directly.
\end{proof}

\paragraph{Normal cones of the Haar body.}
Suppose \(W\) supports \(\mathcal H_{d,m}\) at the noisy image of a
commuting POVM \(Q_i=\sum_a p_i(a)P_a\).
Its known response equals \(p_i(a)\) on \(R_a\).
Every positive branch must therefore attain the pointwise maximum
throughout that cell. Phase variation on interfaces \(p_a=p_b\)
and variation of coordinate moduli show that selected branches have
the canonical form
\begin{equation}
 W_{i(a)}=M+\lambda P_a,\qquad \lambda\ge0,
 \label{geo:normal}
\end{equation}
and every other branch lies below the envelope
\[
 z^*W_jz\le z^*Mz+\lambda\max_a|z_a|^2.
 \label{geo:envelope}
\]
Indeed, the difference of two selected branches must be
\(\lambda_{ab}(P_a-P_b)\); the cocycle relation makes all
\(\lambda_{ab}\) equal, with sign fixed by the maximizing cells.

An unused branch is controlled quantitatively at the limiting
normal. If \(z^*Az\le\max_a|z_a|^2\) for every unit vector, then
\(A\preceq I\), and averaging over flat-amplitude phase states gives
\(\operatorname{tr}A\le1\). Thus
\[
 (\eta_d+c_d)I-\Phi_{\eta_d}(A)
 =\eta_d(I-A)+c_d(1-\operatorname{tr}A)I\succeq0.
\]
This slack is singular only if \(A=P_a\).
In fact singularity forces \(\operatorname{tr}A=1\) and \(Av=v\);
the envelope makes \(v\) a basis vector \(e_a\).
The nonnegative function \(z^*(P_a-A)z\) on \(R_a\) has zero
integral by \eqref{geo:haar-cell}. It vanishes on an open set,
so \(A=P_a\).

\paragraph{Support reduction for qutrits.}
Every supporting witness at the noisy image of a point of
\(\mathsf K_m\) also supports at a commuting POVM.
First take a reducible point and pinch \(W\) to its block,
obtaining \(\overline W\). The unchanged score, the flagged theorem,
and convexity of Haar support give
\[
 h(W)=\operatorname{score}_{\overline W}(E)
 \le q(\overline W)\le h(\overline W)\le h(W).
\]
All inequalities are equalities. Choose an extreme qubit optimizer
for \(\overline W\). If it has one or two rays it is commuting.
If an active contact coefficient \(a_i=0\), that contact annihilates
the entire qubit block and a single block effect is also optimal.
In the remaining three- or four-ray case, all \(a_i>0\) and the rays
are distinct. The strict reserve \eqref{geo:strict-reserve} makes
the pair-cell bound strictly negative, since \(v\ne0\) and
\(\sqrt{|w|^2+t^2|v|^2+2t|w\times v|}\) is strictly increasing in \(t>0\).
Equation \eqref{geo:contact-identity} would give \(h>q\), a contradiction.
A commuting optimizer therefore exists; as it is fixed by the
pinching, it has the same score for the original \(W\).
For a convex combination of reducible POVMs, a vanishing support
deficit forces every positive component's deficit to vanish.
This proves the reduction for \(\mathsf K_m\).
For \(\mathsf P_{d,m}\), the same reduction follows immediately
by taking a commuting component.

\begin{proof}[Proof of Theorem~\ref{geo:operational}]
Both center sets are compact and closed under merging labels and
deleting zero effects. We argue by induction on \(m\), beginning
with the single effect \(I\).
If nonsimulable \(E^n\) approached a center point \(E^0\), choose
separating witnesses \(W^n\), subtract their common matrix average,
and normalize their tuple norm to one. A subsequence converges to
a nongauge support \(W^0\) at \(E^0\).
The support reduction supplies a commuting support point, so
\eqref{geo:normal} and the envelope apply.

If \(\lambda>0\), the unique optimal quantum dual at the limit is
\[
 Y^0=\Phi_{\eta_d}(M)+\lambda(\eta_d+c_d)I.
\]
Feasibility follows from the envelope lemma. The selected basis
labels force every basis diagonal of an optimal dual to attain its
minimal value; positivity then kills all off-diagonal entries,
proving uniqueness.
Nearby optimal duals \(Y^n\) are bounded and converge to \(Y^0\).
Every limiting slack is either positive definite or
\(\eta_d\lambda(I-P_a)\).
Complementary slackness for optimal quantum POVMs makes the active
rays fall into \(d\) shrinking clusters about the basis, with
uniform positive pinching on their kernel complements. Every basis
group is present, since the active effects sum to \(I\).
Discarding inactive labels decreases Haar support without changing
the quantum value. Strict separation and
\eqref{geo:contact-identity} then force \(F_d(C^n)>0\).
This contradicts Lemma~\ref{geo:cusp} as the actual cluster diameter
tends to zero.

If \(\lambda=0\), all selected branches equal \(M\), while
\(W_j^0\preceq M\). Nongauge normalization provides a strict branch
\(D=M-W_j^0\succeq0\), \(D\ne0\), for which
\(\Phi_{\eta_d}(D)\succ0\). Equality of the limiting score forces
\(E_j^0=0\). For large \(n\), merge that label into a selected branch:
the quantum score does not decrease, while deleting the witness
branch does not increase Haar support. Strict separation survives
with \(m-1\) labels. The limit is a center point with its zero label
deleted, contradicting induction.

Thus no nonsimulable sequence can approach any point of the compact
center set. Compactness gives a uniform positive distance for each
fixed \(d,m\) (or fixed \(m\) in the qutrit case).
\end{proof}

\begin{corollary}[Exact boundary within the center classes]
\label{geo:boundary}
After deleting zero effects, a POVM in \(\mathsf P_{d,m}\), or in
\(\mathsf K_m\) for \(d=3\), maps to the boundary of the Haar
postprocessing body at the stated visibility if and only if it is
a projective refinement of one orthonormal basis:
\[
 E_i=s_iP_{a(i)},\qquad s_i>0,\qquad
 \sum_{i:a(i)=a}s_i=1.
\]
Every other such POVM admits a pointwise increase in visibility.
\end{corollary}

\begin{proof}
The Haar body is full dimensional in the affine space \(\sum A_i=I\):
small bounded quadratic perturbations of constant responses realize
all small zero-sum Hermitian tuples by the Haar second moment.
A nontrivial supporting tuple has the normal form above.
The case \(\lambda=0\) forces a zero effect. When \(\lambda>0\),
the explicit optimal dual and complementary slackness force each
nonzero effect into one of its single basis-ray kernels.
Conversely \(W_i=P_{a(i)}\) supports the image of any projective
refinement, since both support and score equal \(H_d\).
For every remaining point, interiority and continuity permit a
small increase of \(\eta\).
\end{proof}

\subsection{The Euclidean selection inequality}
\label{geo:angles}

This subsection proves \eqref{geo:four-selection} and records all
constants used in the contact theorem. Put
\[
 f_t(\theta)=t\sin\theta-1-\cos\theta,\qquad
 g_t(\theta)=4t\sin\theta-2(1-t^2)(1-\cos\theta),\qquad
 H_t=g_t+4f_t.
\]

\paragraph{Planar balanced triples.}
Let three unit vectors in a plane enclose the origin, with cyclic gaps
\(\alpha,\beta,\gamma\in[0,\pi]\) summing to \(2\pi\).
For pair \(i,j\) and singleton \(k\), define \(w,v,A\) as in
\eqref{geo:pair-cell}, and set
\[
 D_k=|w|^2+t^2|v|^2+2t|w\times v|-A^2.
\]
All cyclic oriented sine terms are nonnegative. Expanding gives
\begin{equation}
 D_k=t^2(a_i-a_j)^2+a_ia_jg_t(\theta_{ij})
 +4a_k\{a_if_t(\theta_{ki})+a_jf_t(\theta_{jk})\}.
 \label{geo:planar-deficit}
\end{equation}
Take \(t=3/4\), \(\theta_0=2\arctan(1/t)\), and
\(\theta_1=2\arctan t=\pi-\theta_0\).
One has \(f_t\ge0\) on \([\theta_0,\pi]\) and
\(g_t\ge0\) on \([0,4\arctan t]\).
If two gaps are at least \(\theta_0\), choose their common vertex as
singleton; the remaining gap is at most \(4\arctan t\), so every
term in \eqref{geo:planar-deficit} is nonnegative.

Otherwise exactly one gap \(\alpha\) is at least \(\theta_0\), as
\(\theta_0<2\pi/3\). Its other gaps lie in
\([\theta_1,\theta_0]\). Add the deficits at the two endpoints \(i,j\)
of \(\alpha\):
\[
 D_i+D_j=t^2\{(a_j-a_k)^2+(a_i-a_k)^2\}
 +8a_ia_jf_t(\alpha)+a_ia_kH_t(\gamma)+a_ja_kH_t(\beta).
\]
On that middle interval,
\(\sin\theta\ge2t/(1+t^2)\) and
\(\cos\theta\le(1-t^2)/(1+t^2)\), hence
\[
 H_t(\theta)\ge
 \frac{4(t^4+3t^2-2)}{1+t^2}=\frac1{100}.
\]
Thus some deficit is nonnegative. Monotonicity in \(t\) gives
\eqref{geo:four-selection} for every balanced subset of at most
three rays, including limiting degeneracies.

\paragraph{A nearest-facet inequality.}
Now suppose four rays form a tetrahedron with the origin in its
interior; all other cases have a balanced subset of at most three.
Choose a facet of least plane distance \(h>0\) from the origin,
with outward normal \(e\), and write its vertices
\[
 n_i=he+Ru_i,\qquad R=\sqrt{1-h^2},\qquad u_i\in S^1\subset e^\perp.
\]
The radius-\(h\) ball lies in the tetrahedron and touches this facet
in its relative interior: equality with another facet would force
its unit normal to equal \(e\).
The projected triangle therefore encloses zero. Its cyclic gaps
\(\alpha,\beta,\gamma\) are in \((0,\pi)\), and
\[
 \Pi=\cos(\alpha/2)\cos(\beta/2)\cos(\gamma/2)>0,\qquad
 p=\frac{h^2}{R^2}.
\]
We claim
\begin{equation}
 \frac{h}{1-h}\le\sqrt{2\Pi},\qquad
 p\le\frac{2\Pi}{1+2\sqrt{2\Pi}}.
 \label{geo:facet}
\end{equation}
Write the fourth vertex as \(r-se\), \(0<s\le1\).
For an edge let \(\nu\) be its outward planar unit normal,
\(b=R\cos(\theta/2)>0\) its chord distance, and
\(\ell=2R\sin(\theta/2)\) its length. The adjacent facet equation is
\[
 (s+h)\nu\cdot x+(\nu\cdot r-b)z=sb+h\nu\cdot r.
\]
Its right side is positive, since the origin is in the interior.
Its distance from the origin is at least \(h\); squaring yields
\[
 \nu\cdot r\ge
 \frac{h(s+h)}{2b}-\frac{(s-h)b}{2h}.
\]
Multiply by \(\ell\), sum, and use the polygon identity
\(\sum\ell\nu=0\). This forces \(s>h\) and gives
\[
 h^2\frac{s+h}{s-h}\le\frac{\sum\ell b}{\sum\ell/b}
 =2R^2\Pi.
\]
The equality follows from the half-gap sum \(\pi\) and the identities
\(\sum\sin2A=4\sin A\sin B\sin C\),
\(\sum\tan A=\tan A\tan B\tan C\).
Since \(s\le1\), this gives \(h^2\le2(1-h)^2\Pi\), proving
\eqref{geo:facet}. In particular, concavity of \(\log\cos\) gives
\(\Pi\le1/8\) and \(h\le1/3\).

\paragraph{Lifting the planar deficit.}
Set \(t=9/10\). If \(w_0,v_0\) are the vectors for the projected
triangle, the normal component of \(w\times v\) has magnitude
\(R^2|w_0\times v_0|\), while
\[
 w=Rw_0+h(a_i+a_j-2a_k)e,\qquad
 v=Rv_0+h(a_i-a_j)e.
\]
Consequently
\[
 D_k\ge R^2D_k^{\rm planar}
 +h^2\{t^2(a_i-a_j)^2-8a_k(a_i+a_j)\}.
\]
Dropping the last positive square and using
\eqref{geo:planar-deficit},
\begin{equation}
 \frac{D_k}{R^2}\ge
 t^2(a_i-a_j)^2+a_ia_jg_t(\theta_{ij})
 +4a_k\{a_i(f_t(\theta_{ki})-2p)
             +a_j(f_t(\theta_{jk})-2p)\}.
 \label{geo:lifted}
\end{equation}
Call a gap large when it is at least $2\pi/3$. We prove the following
three bounds from \eqref{geo:facet}:
\begin{align}
 f_t(\theta)-2p&\ge0
 &&\text{for either large gap if two are large},
 \label{geo:coeff1}\\
 f_t(\alpha)-2p&\ge-1/200
 &&\text{for every }\alpha\ge2\pi/3,
 \label{geo:coeff2}\\
 H_t(\beta)-8p,\ H_t(\gamma)-8p&\ge1/25
 &&\text{if }\alpha\text{ is the unique large gap}.
 \label{geo:coeff3}
\end{align}

For \eqref{geo:coeff1}, put \(B=\theta/2\in[\pi/3,\pi/2]\),
\(z=\cos B\), and
\[
 r=\frac{\sqrt3}{2}\sin B-\frac z2,\qquad
 A=\frac{\sqrt3}{2}\sin B-z.
\]
The second large gap implies \(2\Pi\le zr\).
Both \(r,A\) increase with \(B\), while \(z\) decreases.
At \(B=\pi/3\), \(16rA^2=z\), so \(16rA^2\ge z\) throughout.
By \eqref{geo:facet} and monotonicity of \(u/(1+2\sqrt u)\),
\[
 p\le\frac{zr}{1+2\sqrt{zr}}\le zA
 =\frac12 f_{\sqrt3/2}(2B).
\]
The middle inequality is equivalent, by positive squaring, to the
inequality just proved. Since \(9/10>\sqrt3/2\),
\eqref{geo:coeff1} follows.

For \eqref{geo:coeff2}, put \(z=\cos(\alpha/2)\in[0,1/2]\).
The product of the other two half-gap cosines is at most
\((1-z)/2\). Thus
\[
 f_t(\alpha)-2p
 \ge2z\left[\frac9{10}\sqrt{1-z^2}-q(z)\right],\qquad
 q(z)=z+\frac{1-z}{1+2\sqrt{z(1-z)}}.
\]
On \(0\le z\le1/100\), use \(q(z)\le1\) and
\(\sqrt{1-z^2}\ge1-z^2\), obtaining
\[
 f_t(\alpha)-2p
 \ge-\frac{10009}{5000000}>-\frac1{200}.
\]
On \(1/100\le z\le1/4\), let \(k=\sqrt{z(1-z)}\).
The sign of \(q'\) is the sign of
\[
 N=(1-z)(4z-1)+4k^3,\qquad N'=5-8z+6k(1-2z)>0.
\]
Hence \(q\) has no interior maximum. Its endpoint values are at
most \(1/100+99/118<17/20\) and \(2/3<17/20\), whereas
\((9/10)\sqrt{1-z^2}\ge(9/10)\sqrt{15/16}>17/20\).
On \(1/4\le z\le1/2\), \(N\ge0\), so \(q(z)\le q(1/2)=3/4\);
meanwhile \((9/10)\sqrt{1-z^2}\ge9\sqrt3/20>3/4\).
These bounds prove \eqref{geo:coeff2}.

For \eqref{geo:coeff3}, the two middle gaps belong to
\([\pi/3,2\pi/3]\). For either such \(\theta\), maximizing the other
two half-gap cosine factors subject to one being large gives
\[
 p\le2\Pi\le\frac{\sqrt3\sin\theta-1-\cos\theta}{4}.
\]
Since
\(H_{9/10}=(36/5)\sin\theta-219/50-(181/50)\cos\theta\),
\[
 H_{9/10}(\theta)-8p\ge
 \left(\frac{36}{5}-2\sqrt3\right)\sin\theta
             -\frac{119}{50}-\frac{81}{50}\cos\theta.
\]
This expression is concave on the stated interval. Its smaller
endpoint value is \(18\sqrt3/5-619/100>1/25\), the last inequality
following after positive squaring from \(388800>388129\).

\paragraph{Completing selection.}
If two gaps are at least \(2\pi/3\), choose their common vertex
as singleton. The remaining gap is at most \(2\pi/3\), where
\(g_{9/10}\ge0\); \eqref{geo:coeff1} makes \eqref{geo:lifted}
nonnegative.
Otherwise let \(a,b\) be the amplitudes at the endpoints of the
unique large gap, and \(c\) the third amplitude. Add the endpoint
deficits, and use \eqref{geo:coeff2}--\eqref{geo:coeff3}:
\[
 \frac{D_i+D_j}{R^2}\ge
 \frac{81}{100}\{(a-c)^2+(b-c)^2\}
      -\frac{ab}{25}+\frac{c(a+b)}{25}
 =\frac{79}{200}(a+b-2c)^2
  +\frac{83}{200}(a-b)^2+\frac{c^2}{25}\ge0.
\]
Some deficit is nonnegative, proving \eqref{geo:four-selection}.
All amplitudes may vanish; degeneracies are covered either by
continuity or reduction to a smaller balanced subset.

\paragraph{The optional explicit cusp constant.}
For \(d=3\), the shifted diagonal Gram matrix on \(R_1\) is
\[
 M=\frac1{3888}
 \begin{pmatrix}251&-13&-13\\-13&20&-7\\-13&-7&20\end{pmatrix}.
\]
The exact Haar integrals give this matrix; off-diagonal phase
averages have larger unmultiplied coefficients, at least \(7/648\).
For its positive eigenvalues,
\(\lambda_{\min}(M)\ge4\det M/(\operatorname{tr}M)^2=325/338724\).
Also
\(\|z^*Dz-(7/36)\operatorname{tr}D\|_\infty\le(19/12)\|D\|\).
Dividing the \(L^2\) lower bound by this \(L^\infty\) upper bound
gives \(\kappa_3=325/536313\) in Lemma~\ref{geo:cusp}.



\section{Computational content and verification}
\label{sec:verification}

The computational materials are intended to be usable mathematical
objects: a reader can supply measurement counts, a finite channel, a graph,
a Hermitian generator, or rational points in a simplex and obtain an output
with a stated meaning.
The same materials support verification of finite identities occurring in
the proofs. We distinguish those two purposes from exploratory numerical
sampling.

\subsection{Reusable interfaces}

\paragraph{Ensemble certification.}
The Hesse tool takes the number of triple-copy $Q$ clicks, the nine
single-copy SIC counts, and a confidence level. It returns the bound in
\eqref{geo:confidence}, the binomial upper confidence limit, and the
one-copy moment error. A planning mode returns sufficient sample budgets
for a requested accuracy. Seeded synthetic experiments illustrate the
interface, but the confidence statement is proved for the matched-replica
access model. It does not assume that the unknown ensemble is a finite
mixture of nine states. The support witness and the weight-conditioning
matrix are exposed separately so that the approach can be adapted to a
different finite coherent ensemble.

\paragraph{Channel and graph certificates.}
The network tool accepts rational real or complex Choi matrices, a
faithful rank-one reference, and proposed sandwich constants. It checks
Hermiticity and positive semidefiniteness by exact congruence elimination,
including zero-pivot cases. It then builds the edge influence matrix and
checks a rational positive vector $w$ with $Cw\le rw$. The output records
the coarse gap $1-r$ and the separate microscopic coverage multiplier.
For homogeneous AKLT inputs, the exact certificates include
\[
 (\Delta,n)=(3,4):\quad\operatorname{gap}H_{\rm block}\ge\frac3{16},
 \qquad
 (\Delta,n)=(31,6):\quad\operatorname{gap}H_{\rm block}\ge\frac7{432}.
\]
These are blocked-parent bounds; a specified microscopic Hamiltonian
still uses its finite-patch gap in Lemma~\ref{q:lem-seed}. Rational
certificates can be checked without diagonalizing the full spin system.

\paragraph{Positive modular and information-matrix approximations.}
The modular tool constructs positive Gauss and Radau clocks from exact
Bessel moments or supplied positive moment data. The output includes
positive nodes, Gamma shapes, a heat drift when needed, the derivative
order in the error bound, and an exact rounding certificate obtained from
Sturm isolation and rational interval arithmetic. A sparse Pauli routine
computes nested commutators by combining integer coefficients before
taking an absolute sum. This yields a certified operator-norm upper bound
for use in Theorem~\ref{op:quadrature-theorem}.

For example, for
\[
 G_0=\frac{37\sum_jZ_jZ_{j+1}+29\sum_jX_j+13\sum_jZ_j}{100},
 \qquad A=Z_{\rm center},
\]
the second, fourth, and sixth commutator bounds are respectively
$1.3456$, $6.80685216$, and $49.483070071552$. These displayed decimals
are terminating exact rational values. Their support sizes are three,
five, and seven sites, so they agree on sufficiently long open chains.
The calculation concerns the supplied $A$ and $G_0$; it does not infer
locality of a relative tangent from locality of an unnormalized perturbation.

One sampled use of a compiled rule requires its Gamma draws, one Gaussian
draw, and one evolution by $e^{iTG}$. Their times add before evolution.
For $G=K/2$ the evolution is $e^{iTK/2}$. At the hyperbolic endpoint and
semigroup time one, a direct logistic sampler is already available; the
advantage of the common clock construction is its fractional-time
semigroup, certified deterministic approximations, and fixed-state
information brackets. Access to $K$, state preparation, feature inner
products, and Hamiltonian simulation precision remain explicit inputs.

\paragraph{Coherent spectral combinatorics.}
The spectral tool enumerates restricted-growth words for set partitions,
constructs the directed cycle quotients, and evaluates rooted cofactors by
fraction-free integer elimination. It produces rational moments at
arbitrary integer complex dimension and rational sampling density. A
separate extension check verifies the noncrossing layer, Narayana
multiplicities, the aspect-ratio factors, and pairing bounds used for
unbounded support. Optional Monte Carlo generates coherent Gram matrices
at finite degree, with exact finite-degree expectations as the comparison.
The random experiment is separated from the default exact calculation.

\paragraph{Certified Slater-energy assignments.}
The transport tool accepts positive rational squared moduli for $N_k$
projective points. A numerical assignment proposes a matching of these
points to the degree-$Dk$ exponent grid. A separate rational calculation
then certifies optimality: multiplicative dual variables majorize every
monomial weight and attain equality on the proposed matching. This proves
the optimum without treating floating-point comparisons as exact.
Arb ball arithmetic encloses the logarithms and the two endpoints in
Proposition~\ref{tor:finite-certificate}. The output bounds the angular
mean of the actual determinant energy for the supplied moduli. It is
useful for comparing radial configurations and for investigating toric
energy landscapes. The certificate has precisely this scope; it does not
bound every phase configuration or an exponential moment of the energy.

\paragraph{Sampling geometric potentials.}
The nonlinear Green theorem supplies a further algorithmic interface,
specified mathematically in Corollary~\ref{cy:interpolation-sampling}.
Given evaluations of a section basis, its Hermitian norm, and independent
samples from a smooth positive volume, solve a homogeneous interpolation
system, evaluate the resulting section, and subtract the mean logarithm
of its squared Hermitian norm on fresh calibration samples. Averaging $M$ such functions estimates the
finite-level potential. With $J$ calibration points per section, its
expected squared $L^2$ error is at most
$2C_a(\mu)(1+1/J)/(a^2M)$, uniformly in the level. The polynomial-level
bias converges, but its rate and finite-precision interpolation conditioning
remain questions for an implementation. This result is an algorithmic
consequence of the proof; the executable suites do not simulate a
Calabi--Yau metric or certify a finite-level bias.

\subsection{What the certificates establish}

Table~\ref{tab:claim-support} summarizes the support for the principal
claims. The external inputs are cited where they enter the proofs. The
machine-readable registry expands these rows into individual statements,
hypotheses, dependencies, and arithmetic assumptions.

\begingroup
\small
\begin{longtable}{p{.22\linewidth}p{.42\linewidth}p{.26\linewidth}}
\caption{Mathematical support and the role of finite verification.}
\label{tab:claim-support}\\
\toprule
Main result & Decisive argument & Computational support\\
\midrule
\endfirsthead
\toprule Main result & Decisive argument & Computational support\\
\midrule
\endhead
Balanced threshold, Theorem~\ref{geo:sharp}
& Jet semi-invariants, balanced determinantal couplings, and fixed-degree
quantization prove the sharp strict-subcritical coefficient.
& Exact jet maps and monomial examples check constants and conventions.\\[5pt]
Cubic classification and Hesse certification,
Theorem~\ref{geo:cubic} and Section~\ref{geo:hesse-cert}
& A quartic rank witness supplies full tensor balance; Cartan theory
classifies support. A pointwise polynomial inequality gives robustness.
& Exact tensor, witness, moment, and Gram-matrix identities.\\[5pt]
Singer surface, Section~\ref{geo:k3}
& Blow-up and cover calculations, automorphisms, and an integral
discriminant-form computation determine the K3 lattice.
& Cyclotomic cancellation and intersection-lattice arithmetic.\\[5pt]
Coherent spectra and toric energy,
Sections~\ref{sec:coherent-spectrum} and \ref{tor:section}
& Local Gaussian graph integrals yield tree moments; exposed Newton
vertices and Jensen's formula bound the torus-averaged logarithm.
& Partition enumeration and exact assignment duals, with certified
logarithmic endpoints.\\[5pt]
Equilibria and determinant LDPs,
Sections~\ref{sec:equilibria}, \ref{sec:repulsive-constraints},
and \ref{att:entropy-section}
& Positive moment repair, pluripotential duality, conditional compactness,
entropy coarse-graining, and two-sided mean-energy transfer.
& Finite repair identities check normalization; the limit theorems have
analytic proofs.\\[5pt]
Nonlinear Green formula, Section~\ref{cy:section}
& The determinant Hessian is a positive hierarchy of coherent moments.
Convex secants identify the limiting derivative with the inverse
Monge--Amp\`ere potential.
& A proved sampling-variance bound separates statistical error from
the convergent, nonquantified level bias.\\[5pt]
Curvature and susceptibility, Section~\ref{ac:section}
& Singular anticanonical metrics are regularized before applying
H\"ormander's estimate; projective changes of variables give the Ward identity.
& The identities are proved at finite particle number.\\[5pt]
Modular and gap certificates,
Sections~\ref{sec:positive-calculus} and \ref{q:sec-networks}
& Positive random-time representations, local commutators, two-port
Schur cancellation, and projection-angle estimates.
& Rational clocks, Arb brackets, Choi positivity, and graph inequalities;
bounded matrix diagnostics check implementations.\\[5pt]
Measurement neighborhoods, Theorem~\ref{geo:operational}
& Convex dual contact identities, anisotropic flag estimates, and
compactness at reducible qutrit strata.
& Exact finite-angle inequalities; the neighborhood conclusion uses
the full analytic contact proof.\\
\bottomrule
\end{longtable}
\endgroup

The exact calculations cover Moore determinants, quartic contractions,
finite Hesse moments and Gram matrices, Singer character cancellation and
intersection forms, finite jet moment maps, scalar response identities,
Choi positivity, and rational graph inequalities. They are particularly
useful where a transpose convention, normalization, or cancellation could
change the statement of a theorem. Finite jet examples verify the
implementation of the all-degree argument, rather than replacing it.

Floating-point checks have a different role. Direct physical tensor
contractions test the CP convention; numerical principal angles test the
quotient-metric formula; Fourier quadrature tests modular kernels against
matrix divided differences. Dense eigenvalue experiments and the sampled
Hesse scatter plot illustrate formulas and can expose implementation
errors. Their output is labeled numerical evidence and is not used as a
certificate of a universal inequality.

The packaged run passed all 25 suites: 18 exact-algebra suites, three
certified-arithmetic suites, and four bounded floating-point diagnostic
suites. Monte Carlo experiments were disabled in this run. Exact coverage
includes 26,442 set partitions through order nine, an independent
5,295-partition spectral audit, 50 finite posets, and 18 Hesse identities.
The clock tools checked 32 rational quadrature isolations, 72 Lee--Yang
rules, and 16 Arb-certified information brackets; the eighth-order BKM
example has certified bracket width below $1.893\times10^{-23}$. These
counts describe finite computational coverage, not the scope of a theorem.
The assignment suite additionally checks 720 Newton-exposure assignments,
319 exact multinomial mode inequalities, and 720 exhaustive small
assignment comparisons; its default rational example verifies all 100
dual inequalities.

For the analytic probability results, the evidence is the proof itself.
The attractive LDP follows from the two-sided estimate
\eqref{att:quadratic-bound}, compact conditional densities, the exact
coarse-graining entropy inequality, and finite-reference energy transfer.
No finite simulation certifies a thermodynamic limit. Likewise, the
curvature and Ward identities are derived before any limit is taken.
The theorem registry records the classical external results and each
intermediate lemma, so their verification can proceed independently of
the numerical illustrations.

The aggregate runner records the command, return status, elapsed time,
arithmetic category, and result path for every check. The report includes
the Python and package versions and fixes random seeds. Exact rational
and algebraic calculations rely on the correctness of the supplied code
and the arithmetic libraries. No proof-assistant kernel has checked the
full analytic or geometric arguments in this version. The accompanying
formalization guide identifies the definitions, theorem dependencies,
external theorems, and finite certificates needed for such a development;
its machine-readable registry records this verification status explicitly.

\subsection{Figures and reproducibility}

All thirteen figures are generated from the equations in the text by one
script, with PDF and PNG outputs and a formula-and-data manifest. The
plotting style and the surface inset use the supplied Matplotlib identity
package; font notices are retained with that code. No example data from
the styling package are used. The simplex projections, curve operations,
and core-wire architecture are labeled as geometric illustrations; the
remaining curves and finite tables are computed from explicit formulas.
The figures and compiled pages were visually inspected for mathematical
labeling, caption consistency, clipping, and legibility.

The main source, bibliography, proof modules, figure code, executable
supplements, result files, and verification guide are distributed together.
The mathematical proofs are contained in the paper; the dependency guide
makes their interfaces easier to formalize and reuse.

\section{Further directions}
\label{sec:future}

\subsection{Beyond the section integrability threshold}

The attractive LDP is established for $0\le\gamma<1/3$ on the projective
plane, and below the corresponding normalized-potential threshold on any
positive polarized compact complex manifold. On projective space it also
holds in shrinking design tubes of every fixed order. The next objective
is a neighborhood of the canonical coupling $\gamma=1$, where the balanced
coefficient and the microscopic spectral law suggest a continuation toward
reduced Gibbs measures in complex geometry
\citep{AndreassonBermanSvensson2026}. Let $\mathcal D_1$ be the balanced
fibre on $\mathbb P^2$, and condition the reference product measure on
$\|M_1(L_N)-I_3/3\|_{\rm HS}<N^{-a}$, with fixed $0<a<1/2$.
Write this conditioned measure as $Q_N^{(1)}$, with
$N=N_k=\binom{3k+2}{2}$ as in Section~\ref{sec:equilibria}.

\begin{conjecture}[Attractive reduced Laplace principle]
\label{conj:attractive}
There is an $\epsilon>0$ such that, for
$\gamma\in(1-\epsilon,1+\epsilon)$ and every continuous
$\Phi:\mathcal P(\mathbb P^2)\to\R$, the partitions are finite for
all sufficiently large $k$ and
\begin{equation}
 \lim_{k\to\infty}\frac1N\log
 \int e^{N\gamma E_k-N\Phi(L_N)}\,dQ_N^{(1)}
 =-\inf_{\substack{\mu\in\mathcal D_1\\H(\mu)<\infty}}
       [H(\mu)-\gamma E(\mu)+\Phi(\mu)].
 \label{eq:attractive-target}
\end{equation}
\end{conjecture}

This would extend the constrained attractive LDP and its pressure to the
canonical range, including continuous empirical tilts. The first-moment coefficient
$12/7$ supplies macroscopic coercivity beyond the canonical coupling
$\gamma=1$. A microscopic route is to prove, for the attractive tilted
Gibbs laws $P_{k,\gamma,\Phi}$ and the continuous energy envelopes,
\[
 \lim_{r\to\infty}\limsup_{k\to\infty}
 \mathbb E_{P_{k,\gamma,\Phi}}
       [E_k-E_r(L_N)]_+=0.
\]
Lower comparison controls the opposite side of this residual. Through
\eqref{eq:coherent-logarithmic-energy}, the new estimate concerns
eigenvalues on exponentially small scales, rather than fixed positive
spectral moments. The directed-tree law supplies a reference-measure
benchmark; controlling its change under attractive reweighting is an
additional requirement. Extending the coupling
range toward the full macroscopic threshold is a further question.
The proof below $1/3$ identifies a specific place to seek an improvement:
compactness of the conditional densities after exploiting the global
moment constraint. The finite assignment theorem supplies an independent
computable description of torus-averaged energy, but a stronger
exponential estimate would be needed to carry it to the canonical range.

\subsection{Hermitian algebraic support and higher designs}

The cubic classification and the Hesse certificate invite a general theory
of positive design measures on determinantal varieties. A useful target
is to identify formats in which a rank witness converts some scalar
marginals into full moment-map balance, then classify the relevant closed
orbits. The equality cases $(3,3)$ and $(4,2)$ in the rank-witness
calculation show that this mechanism has a specific representation-theoretic
boundary. Beyond it, additional equations or higher moments are needed.
Quantitative ideal-to-distance bounds would turn those equations into
finite-copy ensemble certificates. Exact thresholds for higher-design
entropy fibres and the geometry of their concentrating supports are
natural problems in their own right.

The Singer construction suggests studying cyclic covers attached to
weighted designs with special vanishing divisors. Here the divisor
relation, automorphisms, and integral discriminant form determine the K3
lattice. One can ask which positive moment constraints select similarly
distinguished surfaces and whether lattice data classify the associated
ensemble faces. Such a classification would join Hermitian interpolation
to the geometry of surfaces through explicit constructions.

\subsection{Undecorated networks and spatial modular response}

The wire criterion covers arbitrary cores and uniformly primitive
attachments, including decorated AKLT models on every bounded-degree
simple graph. For undecorated two-dimensional parents, the corresponding
objective is a finite collar or regional-response estimate that implies
the required subspace angles with a uniform microscopic energy scale.
Boundary-state factorization provides a complementary setting for this
question \citep{Kastoryano2026}. The local-ratio and interaction-norm
examples show why a spatial argument must control the actual compatible
response or the support growth of modular evolution. Complete positivity
alone gives ordinary norm contraction, not an interaction-norm estimate.
This distinction is also relevant to effective Hamiltonians and strong
mixing criteria for quantum Gibbs states
\citep{BluhmCapelPerezHernandez2025}.

\subsection{Positive information algorithms and sharper contact geometry}

The clock construction admits any certified positive zero measure of the
stated form. Further input families from Lee--Yang theory could enlarge
the class of quantum score filters with positive finite approximations.
The practical questions are then quantitative: tight observable-specific
commutator bounds, optimal allocation of deterministic and sampling error,
and efficient access to the required weighted inner products. The
information-matrix bracket is a concrete target for these improvements.

The nonlinear Green formula gives a complementary computational problem
in complex geometry: obtain an effective polynomial-level bias rate and
stable interpolation rules for the random-section estimator. The
coherent-moment Hessian in \eqref{cy:hessian-hierarchy} offers positive
finite truncations for studying this question. The sampling variance is
already controlled uniformly in the level; geometric approximation error
and numerical conditioning are the additional quantities needed for a
practical metric reconstruction.

For measurements, the qualitative neighborhood at $5/12$ can be refined
by deriving explicit radii as functions of the number of outcomes and the
distance to the classified contact strata. The genuinely irreducible
qutrit region is a separate geometric problem. Combining its dual
contact inequalities with finite approximation methods for compatibility
\citep{PortoDesignollePokuttaQuintino2026} would make the theory more
directly usable in quantum-information experiments and optimization.

\section*{AI use and research contributions}
\addcontentsline{toc}{section}{AI use and research contributions}

This project was directed by a human researcher. The human contribution
included identifying the initial spectral-gap question, setting the goal
of finding theory with significance beyond that example, choosing the
research priorities, guiding the exploratory computations, and deciding
which conjectures and connections merited further investigation. The
direction of the project was repeatedly revised in response to the
human researcher's questions, criticism, and insistence on substantive
connections between different areas of mathematics.

ChatGPT and collaborating AI agents contributed extensively to the
mathematical work. Their roles included early-stage conjecture generation,
literature searches, construction and examination of examples, searches
for counterexamples, development of proposed proofs, independent checks
of arguments, implementation of symbolic and numerical computations,
and preparation of the exposition and figures. Early exploratory
calculations were used to generate and test possibilities. The results
retained in the paper are supported by the arguments and evidence
identified in Section~\ref{sec:verification} and the accompanying theorem
registry; exploratory agreement by itself is not used as a proof.

SAIR's Open Math OS was used principally for autoformalization
\citep{SAIROpenMathOS2026}. This complemented the proof development and
computational work described here. The verification record distinguishes
prose proofs, exact finite calculations, certified arithmetic, numerical
diagnostics, and proof-assistant verification. The present release does
not claim that a proof-assistant kernel has checked the complete paper.
Claims of proof-assistant verification are limited to theorem entries
with an accompanying checked artifact and explicit dependencies in the
registry. Use of an autoformalization platform by itself is not such a
claim.

\section*{Acknowledgements}
\addcontentsline{toc}{section}{Acknowledgements}

We thank the SAIR Foundation for Open Math OS and the infrastructure
supporting the autoformalization work. The mathematical figures use the
Matplotlib identity kit made available for this project; its code and
font notices are retained in the accompanying materials.

\clearpage
\bibliographystyle{plainnat}
\bibliography{references}
\end{document}
